\documentclass[12pt,b4paper]{article}

% basic packages
\usepackage{bigints}
% \usepackage[T1]{fontenc}
%%%%%%%%%%%%%%%%%%%%%%%%%%%%%%%%%%%%%%%%%
\newcommand\myfontsizeTitle{\fontsize{100pt}{100pt}\selectfont \bf \sffamily}
% \renewcommand\familydefault{\sfdefault}
\allowdisplaybreaks
\parindent 0in
\parskip 1.5ex
\usepackage{float}
% define a list of equations
\usepackage{tocloft}    % for list of equations
\usepackage{ragged2e}   % to undo \centering
\usepackage{hyperref}   % to make references hyperlinks
\usepackage{glossaries} 
\newcommand{\listequationsname}{\Large{List of Equations}}
\newlistof{myequations}{equ}{\listequationsname}
\newcommand{\myequations}[1]{
   \addcontentsline{equ}{myequations}{\protect\numberline{\theequation}#1}
}
\usepackage{array}
\newcolumntype{P}[1]{>{\centering\arraybackslash}p{#1}}
\usepackage{tikz}
\usetikzlibrary{calc, positioning, shapes, backgrounds, fit, arrows}
\usepackage{siunitx}
\usepackage{contour}
\usepackage{pgf-spectra}
\usepackage{longtable}
\usepackage{latexsym}
\usepackage{siunitx}
\usepackage{setspace}
\usepackage{xcolor}
\usepackage{listings}
% Define color for mind map
\definecolor{codegreen}{rgb}{0,0.6,0}
\definecolor{codegray}{rgb}{0.5,0.5,0.5}
\definecolor{codepurple}{rgb}{0.58,0,0.82}
\definecolor{backcolour}{rgb}{0.98,0.98,0.98}
\definecolor{americanrose}{rgb}{1.0, 0.31, 0.0}
\definecolor{americanrose}{rgb}{1.0, 0.01, 0.24}
\definecolor{bleudefrance}{rgb}{0.19, 0.55, 0.91}
\definecolor{cornflowerblue}{rgb}{0.39, 0.58, 0.93}
\definecolor{darkmidnightblue}{rgb}{0.0, 0.2, 0.4}
\usepackage{hyperref}
\hypersetup{
    backref=true,       
    pagebackref=true,               
    hyperindex=true,                
    colorlinks=true,                
    breaklinks=true,
    % linkcolor=[rgb]{0.62,0.22,0.24},
    linkcolor=americanrose,
    filecolor=magenta,      
    urlcolor=americanrose,
    citecolor=americanrose,
    citebordercolor=Violet,
    filebordercolor=Red,
    linkbordercolor=Blue
    }
\usepackage{geometry}
\geometry{
 right=1in,
 left=1in,
 top=1in,
 bottom=1in,
 includeheadfoot,
 headsep=6mm,
 headskip=20mm}
\setlength{\parindent}{0em}
\setlength{\parskip}{0.5em}
\usepackage[pdftex]{graphicx}
\usepackage{amsmath}
\usepackage{amssymb}
\usepackage{pifont}% http://ctan.org/pkg/pifont
\newcommand{\cmark}{\ding{51}}%
\newcommand{\xmark}{\ding{55}}%
\usepackage{amsthm}
\usepackage{epsfig}
\usepackage{mathpazo}
\usepackage{graphicx}
\usepackage{fancyhdr}
\usepackage[super]{nth}
\usepackage[symbol]{footmisc}
\usepackage{caption}
\usepackage[rawfloats=true]{floatrow} 
\restylefloat{figure}
\captionsetup{format=hang,{font={stretch=1.3}}}
\renewcommand{\figurename}{\scshape\nouppercase Figure}
\renewcommand{\tablename}{\scshape\nouppercase Table}
\setlength{\footnotesep}{1.5\baselineskip}
\let\origfootnote\footnote
\renewcommand{\footnote}[1]{%
   \begingroup
   \renewcommand{\footnotesize}{\small}%
   \origfootnote{#1}%
   \endgroup}
\usepackage[sc]{mathpazo}

% page formatting
\usepackage{fancyhdr}
\pagestyle{fancy}
\renewcommand{\sectionmark}[1]{\markright{\textsf{\arabic{section}. #1}}}
\renewcommand{\subsectionmark}[1]{}
\lhead{\scshape Page~~\textemdash~~\textbf{\thepage}\hfill\scshape\nouppercase{\rightmark}}
\chead{}
\rhead{}
\lfoot{}
\cfoot{}
\rfoot{}
\setlength{\headheight}{14pt}

\linespread{1.03} % give a little extra room
\setlength{\parindent}{0.2in} % reduce paragraph indent a bit
\setcounter{secnumdepth}{2} % no numbered subsubsections
\setcounter{tocdepth}{2} % no subsubsections in ToC

\begin{document}

% make a title page
\thispagestyle{empty}
\bigskip \
\vspace{0.1cm}
\noindent
	\begin{figure}
	    \centering
	    \includegraphics[width=.4\linewidth]{bracu_logo.pdf}\end{figure}
     \vspace{-.3
     in}
	\begin{center}
		 % {\fontsize{30}{30}\selectfont \sffamily\bf BRAC University}\par
   %       \vspace{0.05in}
         {\fontsize{22}{22} \selectfont \bf \sffamily Department of Mathematics and Natural Sciences} \par
         \vspace{0.06in}
         \large\sc Pragati Sarani, Merul Badda, Dhaka, 1212\par
	\end{center}
\vskip 15pt
\begin{center}
    {\fontsize{20}{20} \selectfont Lecture Notes on \fontsize{20}{20} \selectfont \bf \sffamily \textcolor{darkmidnightblue}{PHY-112}}
    \vskip 16pt
    {\fontsize{25}{25}\selectfont \sffamily\bf \textcolor{darkmidnightblue}{Principles of Physics-II}}
    \vskip 15pt
    \begin{tabular*}{.96\textwidth}{@{\extracolsep{\fill}}lr}
    \hline\hline\\
%%%%%%%%%%%%%%%%%%%%%%%%%%%%%%%%%%%%%%%%%%%%%%%%%%%%%%%%%%%%%%

% Modify information %%%%%%%%%%%%%%%%%%%%%%%%%%%%%%%%%%%%%%%%%
\large
\textbf{Instructor}: \href{mailto:ext.akiful.islam@bracu.ac.bd}{\textsc{Akiful Islam}} & \large \textbf{FIN}: \textcolor{americanrose}{AZW}\\[10pt]
\large \textbf{E-mail}: \href{mailto:ext.akiful.islam@bracu.ac.bd}{\texttt{ext.akiful.islam@bracu.ac.bd}} & \large \textbf{Consultation Room}: \textcolor{americanrose}{Floor-4, Block-G}\\
&  \\
\hline\hline
\end{tabular*}
% {\fontsize{18}{18} \selectfont \rmfamily Akiful Islam} 
% \vskip 6pt
% {\fontsize{15}{15} \selectfont \ttfamily ext.akiful.islam@bracu.ac.bd} 
\vskip 20pt
\end{center}
This course introduces Electromagnetism, a field of study that explores the interactions between Electric and Magnetic fields. Additionally, it delves into the Electromagnetic Force, recognized as one of the four fundamental forces of nature, and examines the reactions of charged particles when subjected to this force.

The primary topics to be covered in this course encompass Electric Charge, Force, Field, Potential, Capacitance, Resistance, Electric Current, Kirchhoff's laws, Ohm's Law, Electric power, Magnetic Force, Field, Flux, Induction, Faraday's Law, Lenz's Law, Ampère's Law, Gauss's Law, Maxwell's equations, Electromagnetic Oscillations, Impedance, Resonance, Transformers, and Electromagnetic Waves.

The course strongly emphasizes fostering an intuitive understanding of the core principles of Physics. This is accomplished through extensive problem-solving exercises of varying complexity. The course will utilize a practical, calculus-based methodology, making it imperative that students have a solid grasp of calculus before enrollment.

\textsc{\bf Disclaimer}: These notes are based on the PHY-112 course at BracU, taught by the author/instructor during the Summer, Fall 23, and Spring 24 (ongoing) semesters. The author does not assert any originality or copyright ownership over this work. The primary sources for these notes were:
\begin{enumerate}
\setlength\itemsep{0em}
    \item Halliday, D., Resnick, R., \& Walker, J. (2023). \textit{\textcolor{americanrose}{Principles of Physics}}, (Extended, \nth{10} Edition).
    \item Young, Hugh D., Roger A. Freedman. \textit{\textcolor{americanrose}{University Physics with Modern Physics}}, (\nth{15} Edition), 2020.
\end{enumerate}
These notes are not abstract/rigorous for a Physics/Mathematics major. To err is human. Please report any errors found to \href{mailto:ext.akiful.islam@bracu.ac.bd}{\texttt{ext.akiful.islam@bracu.ac.bd}}.
% make a table of contents
\newpage
\tableofcontents
\thispagestyle{empty}
\newpage
\listoffigures \pagebreak
% \listoftables
\newpage
\setstretch{1.1}
% main content
\section*{Prerequisites I: Vector Algebra}
Vector algebra deals with the manipulation and analysis of vectors. It includes the operations of vector \textbf{addition}, \textbf{subtraction}, scalar \textbf{multiplication}, dot (aka scalar) product, and cross (aka vector) product. \textcolor{americanrose}{It is assumed that students have a strong foothold on the prerequisites from having done PHY-111 (Principles of Physics-1) before.}
\thispagestyle{empty}
\subsection*{Vector Addition and Subtraction}
Vector addition is the process of adding two or more vectors together to \textbf{form a new vector}. The sum of two or more vectors is the \textbf{resultant vector}.

Vector \textbf{subtraction} is subtracting one vector from another. It is the reverse operation of vector addition.
\subsection*{Analytical Approach of Vector Addition}
The idea is to resolve all the vectors \textit{to be summed} into their $x$ and $y$-coordinate axes components. Then, add \textit{like} components ($\hat{i}$ with $\hat{i}$, $\hat{j}$ with $\hat{j}$, $\hat{k}$ with $\hat{k}$) to construct the resultant vector.

We use the concept of \textit{vector resolution}. It means to \textbf{resolve} or break apart the vector into all possible basis directions it can contribute to. These contributions are the \textit{components} of the main vector. 3D parts are not shown in the diagrams below.
\begin{figure}[H]
    \centering
    \includegraphics[width=.42\linewidth]{figs/vec_add.pdf}
    \caption{A system with four vectors is shown here.}
    \label{fig:anal_vec_def}
\end{figure}
Assume the following setup in Figure \ref{anal_vec_def}. We have four vectors $\vec{A}$, $\vec{B}$, $\vec{C}$, $\vec{D}$ making $\theta_A$, $\theta_B$, $\theta_C$, $\theta_D$ with the $+x$-axis, according to our counter-clockwise convention of angle measurement. 
\begin{table}[H]
    \centering
    \begin{tabular}{|c|c|c|}
        $\text{Vectors}$ & $x$-components & $y$-components \\\hline
        $\vec{A}$ & $A_x\cos{\theta_A}$ & $A_y\sin{\theta_A}$ \\
        $\vec{B}$ & $B_x\cos{\theta_B}$ & $B_y\sin{\theta_B}$ \\
        $\vec{C}$ & $C_x\cos{\theta_C}$ & $C_y\sin{\theta_C}$ \\
        $\vec{D}$ & $D_x\cos{\theta_D}$ & $D_y\sin{\theta_D}$
    \end{tabular}
    \caption{Resolution of all four vectors into their coordinate components.}
    \label{anal_vec_def}
\end{table}
The resultant will have its $x$-component as the sum of all vector contributions in the $x$-direction and the $y$-component in the designated $y$-direction.
\begin{align*}
    R_x &= A_x\cos{\theta_A}+B_x\cos{\theta_B}+C_x\cos{\theta_C}+D_x\cos{\theta_D}\\
    R_y &= A_y\sin{\theta_A}+B_y\sin{\theta_B}+C_y\sin{\theta_C}+D_y\sin{\theta_D}
\end{align*}
The direction of the resultant will be given by 
\begin{align*}
    \theta=\left\{
                \begin{array}{ll}
                  \displaystyle\tan^{-1}\left|\frac{R_y}{R_x}\right|&;\,\,R_x>0,\,R_y>0\,\,\text{(\nth{1} quadrant)}\\[15pt]
                  \displaystyle\pi-\tan^{-1}\left|\frac{R_x}{R_y}\right|&;\,\,R_x<0,\,R_y>0\,\,\text{(\nth{2} quadrant)}\\[15pt]
                  \displaystyle\pi+\tan^{-1}\left|\frac{R_y}{R_x}\right|&;\,\,R_x<0,\,R_y<0\,\,\text{(\nth{3} quadrant)}\\[15pt]
                  \displaystyle2\pi-\tan^{-1}\left|\frac{R_x}{R_y}\right|&;\,\,R_x>0,\,R_y<0\,\,\text{(\nth{4} quadrant)}
                  \end{array}
              \right.
\end{align*}
Nothing new here. It's just the same argument we used for the Polar to Cartesian coordinate transformation.
\thispagestyle{empty}
\subsubsection*{Geometric Approach of Vector Addition}
\begin{figure}[H]
    \centering
    \includegraphics[width=.42\linewidth]{figs/vec_add_geom.pdf}
    \caption{A system with four vectors is shown here.}
    \label{fig:vec_add_gem}
\end{figure}
The idea is to draw all the vectors according to the given problem and geometrically add them using the \textbf{tip-to-tail} method. The tip of the first vector will be on top of the \textbf{tail} of the second vector, followed by the \textbf{tip} of the second vector on top of the tail of the third vector, and so on and on until all vectors are accounted for.

Then, using the \textit{geometric} vector sum to find the resultant vector.
\thispagestyle{empty}
Assume the same setup above. Keep placing the vectors on top of each other using the tip-to-tail method. Finally, add a vector pointing from the tail of the first vector to the tip of the last vector in the series. That is the resultant vector. Although the analytical method of vector addition may not always place the resultant in its actual quadrant, the geometrical method does.
\thispagestyle{empty}
\section*{Vector Multiplication}
The last and most crucial part of vector algebra. Unlike scalar multiplication, vectors follow a different concept when it comes to multiplication. We need to understand its physical intuition before doing any math.
\subsection*{Scalar (Dot) Multiplication}
The dot product is a scalar quantity obtained by multiplying the magnitudes of the two vectors and the cosine of the angle between them. This multiplication takes the \textbf{projection} of one vector in a direction \textbf{parallel} to another vector.

For example, take the following vectors $\Vec{A}$ and $\Vec{B}$. Assume they make an angle $\theta$ with each other. There can be two ways we calculate scalar multiplication.
\begin{figure}[H]
    \centering
    \includegraphics[width=.7\linewidth]{figs/dot_prod_def.pdf}
    \caption{Geometrical definition of scalar multiplication.}
    \label{fig:dot_prod_def}
\end{figure}
\subsubsection*{Geometrical}
Measure the magnitude of the vectors and simply plugin their value in the following formula:
\begin{align}
    \Vec{A}\cdot\Vec{B}=AB\cos{\theta}=BA\cos{\theta}=\Vec{B}\cdot\Vec{A}
\end{align}
$B\cos{\theta}$ is the projection of $\Vec{B}$ in a direction parallel to $\Vec{A}$.
\subsubsection*{Analytical} Similar to the vector addition, we first resolve the vectors in the system coordinate axes and then take the products of \textit{like} components and sum all the products.
\begin{align}
    \Vec{A}\cdot\Vec{B}=A_xB_x + A_yB_y + A_zB_z
\end{align}
\thispagestyle{empty}
\subsubsection*{Limiting Cases of Scalar Multiplication}
\begin{enumerate}
    \item $\theta=0^\circ$: The dot product is maximum.
    \item $90^\circ>\theta\geq 0^\circ$: The dot product is \textbf{positive} in this range.
    \item $\theta=90^\circ$: The dot product is zero.
    \item $180^\circ>\theta> 90^\circ$: The dot product is \textbf{negative} in this range.
    \item $\theta=180^\circ$: The dot product is minimum.
    \item $270^\circ>\theta> 180^\circ$: The dot product is \textbf{negative} in this range.
    \item $\theta=270^\circ$: The dot product is zero.
    \item $360^\circ\geq\theta>270^\circ$: The dot product is \textbf{positive} in this range.
    \item $\theta=360^\circ$: The dot product is maximum. Superimposes with the \nth{1} case.
\end{enumerate}
\subsubsection*{How to use the Dot Product?}
\begin{enumerate}
    \item If vectors are given in \textit{vector} notation, You may use the \textbf{analytical} method to measure their scalar product.
    \item If vectors are given in \textit{verbal} notation, You may use the \textbf{geometrical} method to measure their scalar product.
    \item In either notation, the angle between the involved vectors may be found using:
    \begin{align}
        \theta=\cos^{-1}\left(\frac{\Vec{A}\cdot\Vec{B}}{AB}\right)=\cos^{-1}\left\{\frac{A_xB_x + A_yB_y + A_zB_z}{\left(\sqrt{A_x^2+A_y^2+A_z^2}\right)\left(\sqrt{B_x^2+B_y^2+B_z^2}\right)}\right\}
    \end{align}
    Remember to place both vectors on top of each other in such a way that their tails overlap.
\end{enumerate}
\textsc{Takeaway}: Dot multiplication is commutative. That is, $\Vec{A}\cdot\Vec{B}=\Vec{B}\cdot\Vec{A}$.
\subsection*{Vector (Cross) Multiplication}
The cross product is a vector quantity perpendicular to the plane of the two vectors being multiplied, and its magnitude is equal to the product of the magnitudes of the two vectors and the sine of the angle between them. The direction of the cross product is determined by the \textbf{right-hand rule} (Fig. \ref{fig:cross_prod_cycle}). This multiplication takes the \textbf{projection} of one vector in a direction \textbf{perpendicular} to another vector.

Take the same two vectors $\Vec{A}$ and $\Vec{B}$ again, making an angle $\theta$ with each other. There are two ways we can calculate vector multiplication.

\subsubsection*{Analytical} Similar to the vector addition, we first resolve the vectors in the system coordinate axes and then take the products of all possible combinations of \textit{unlike} components and follow the following format:
\begin{align}
    \Vec{A}\times\Vec{B}&=\left(A_yB_z-A_zB_y\right)\hat{i}+\left(A_zB_x-A_xB_z\right)\hat{j}+\left(A_xB_y-A_yB_x\right)\hat{k}\\[5pt]
    &=\begin{pmatrix}
        \hat{i} & \hat{j} & \hat{k} \\
        A_x & A_y & A_z \\
        B_x & B_y & B_z
    \end{pmatrix}
\end{align}
\begin{figure}[H]
    \centering
    \includegraphics[width=.7\linewidth]{figs/cross_prod_def.pdf}
    \caption{Geometrical definition of cross multiplication.}
    \label{fig:cross_prod_def}
\end{figure}
\subsubsection*{Geometrical}
Measure the magnitude of the vectors and simply plugin their value in the following formula:
\begin{align}
    \Vec{A}\times\Vec{B}&=AB\sin{\theta}\hat{\eta}_{|\Vec{A}\times\Vec{B}|}\\
    \Vec{B}\times\Vec{A}&=BA\sin{\theta}\hat{\eta}_{|\Vec{B}\times\Vec{A}|}\\
    \hat{\eta}_{|\Vec{A}\times\Vec{B}|}&=-\hat{\eta}_{|\Vec{B}\times\Vec{A}|}
\end{align}
$B\sin{\theta}$ is the projection of $\Vec{B}$ in a direction perpendicular to $\Vec{A}$. The cross product lies in the plane perpendicular to the plane where the parent vector lies in.
\thispagestyle{empty}
\subsubsection*{Limiting Cases of Vector Multiplication}
\begin{enumerate}
    \item $\theta=0$: The cross product is zero.
    \item $90^\circ\geq\theta>0^\circ$: The cross product is \textbf{positive} in this range.
    \item $\theta=90^\circ$: The cross product is maximum.
    \item $180^\circ>\theta>90^\circ$: The cross product is \textbf{positive} in this range.
    \item $\theta=180^\circ$: The cross product is zero.
    \item $270^\circ>\theta>180^\circ$: The cross product is \textbf{negative} in this range.
    \item $\theta=270^\circ$: The cross product is maximum in the opposite direction to the \nth{3} case.
    \item $360^\circ>\theta>270^\circ$: The cross product is \textbf{positive} in this range.
    \item $\theta=360^\circ$: The cross product is zero. Superimposes with the \nth{1} case.
\end{enumerate}
\subsubsection*{How to use the Cross Product?}
\begin{enumerate}
    \item If vectors are given in \textit{vector} notation, You may use the \textbf{analytical} method to measure their cross product.
    \item If vectors are given in \textit{verbal} notation, You may use the \textbf{geometrical} method to measure their cross-product.
    \item In either notation, the angle between the involved vectors may be found using:
    \begin{align}
        \theta=\sin^{-1}\left\{\frac{\left|\Vec{A}\times\Vec{B}\right|}{AB}\right\}
    \end{align}
\end{enumerate}
\textsc{Takeaway}: Cross multiplication is anti-commutative. That is, $\Vec{A}\times\Vec{B}=-\Vec{B}\times\Vec{A}$.
\subsubsection*{Rule of Thumb to Cross Multiplication}
Unless explicitly mentioned otherwise, always \textbf{follow the right-handed system} Fig. \ref{fig:cross_prod_def} when doing maths with cross products.
\begin{figure}[H]
    \centering
    \includegraphics[width=.56\linewidth]{figs/cross_prod_cycle.pdf}
    \caption{The coordinate axes unit vectors follow the cyclic relation in 3D space. The Left-handed system is just the mirror opposite of the Right-handed system. The reflection symmetry lies in the $x$-$y$ plane.}
    \label{fig:cross_prod_cycle}
\end{figure}
\thispagestyle{empty}
\section*{Vector Division Makes no Sense}
Vector division is impractical. Because unlike the mathematical operations of addition and subtraction, which can be applied to vectors with clear geometric interpretations, the operation of division does not have a precise geometric interpretation for vectors.

However, if we divide a vector with a positive scalar, the vector gets down-scaled, more like shrinks, while keeping the same direction. Do the same with a negative scalar; the vector shrinks in the opposite direction.
\thispagestyle{empty}
\section*{Prerequisites II: Fields}
A field is a physical quantity that fills space and can vary in magnitude and direction at each point. Fields describe how different entities or properties, such as forces or energy, are distributed throughout space.

Fields are typically defined by their value at each point in space and can be represented mathematically as functions or vector quantities. The value of a field at a particular point determines the magnitude and direction of its effect on other objects or particles in that region.
\subsection*{Vector Fields}
They are characterized by a vector quantity assigned to each point in space and time. A vector field can be visualized by imagining arrows at each point in space, where the direction and length of the arrow represent the direction and magnitude of the vector at that point. The vector field provides information about the direction and strength of the vector quantity at any given point in the region. Examples of vector fields include electric, magnetic, and velocity fields.
\begin{figure}[H]
    \centering
    \includegraphics[width=.9\linewidth]{figs/types-of-fields.pdf}
    \caption{Vector fields require the magnitude and the direction at every point in the coordinate grid. The scalar field, however, needs only the magnitude represented by the pseudo color map.}
    \label{fig:fields-def}
\end{figure}
\subsection*{Scalar Fields}
They are characterized by a scalar quantity assigned to each point in space and time. Scalar fields can be visualized by imagining a numerical value assigned to each point in space. The value can vary from point to point, representing the magnitude of the scalar quantity at that specific location. Examples of scalar fields include temperature distribution, pressure distribution, and gravitational potential.
\thispagestyle{empty}
\section*{Vector Calculus}
Vector calculus is a branch of mathematics that deals with vector fields and their derivatives. It provides a robust framework for describing and analyzing various physical phenomena that involve quantities with both magnitude and direction, such as forces, electric and magnetic fields, fluid flow, and more.

There are three main components of vector calculus: gradient, divergence, and curl. Each of these operations captures different aspects of vector fields and offers valuable insights into the behavior and properties of these fields.
\begin{figure}[H]
    \centering
    \includegraphics[width=.9\linewidth]{figs/grad-div-curl.pdf}
    \caption{(Left) Gradient measures a scalar field's spatial change/rate. (Middle) Divergence measures the divergence and convergence of the vector field from a point source. (Right) Curl measures the vector field's rotation around a point source.}
    \label{fig:grad-div-curl}
\end{figure}
\subsection*{Gradient} The gradient represents the spatial variation or the rate of change of a scalar field. A scalar field assigns a scalar value (e.g., temperature, pressure, electric potential) to every point in space. The gradient of a scalar field yields a vector field that points in the direction of the steepest increase of the scalar field. In simple terms, the gradient tells us how and in which direction a scalar quantity changes rapidly.
\begin{align}
    \Vec{\nabla}\left(\text{Scalar Field}\right)\to\text{Vector}.
\end{align}
The gradient helps us understand how a scalar quantity changes in space, providing insights into heat, fluid, or electric potential flow.
\thispagestyle{empty}
\subsection*{Divergence} The divergence measures the outflow of a vector field, i.e., the tendency of a vector field to flow outward or inward from a given point. It quantifies the spreading or convergence of vector field lines. If the divergence at a point is positive, the vector field is \textbf{diverging} or spreading out from that point. If the divergence is negative, it suggests that the field is \textbf{converging} or coming together at that point. The divergence is associated with the flow of a vector field through a closed surface and is related to the presence of sources or sinks within that region.
\begin{align}
    \Vec{\nabla}\cdot\left(\text{Vector Field}\right)&\to\text{Scalar},\\
    \text{where }\Vec{\nabla}\cdot\left(\text{Vector Field}\right)>0&\Rightarrow\text{Diverging field}\notag\\
    \Vec{\nabla}\cdot\left(\text{Vector Field}\right)<0&\Rightarrow\text{Converging field}\notag\\
    \Vec{\nabla}\cdot\left(\text{Vector Field}\right)=0&\Rightarrow\text{Solenoidal field}\notag
\end{align}
The divergence informs us about the behavior of a vector field in terms of its sources (positive divergence) or sinks (negative divergence), which is crucial for understanding fluid flow, electric charge distribution, or the behavior of gravitational fields. 
\subsection*{Curl} The curl measures the circulation of a vector field, i.e., the tendency of a vector field to circulate or form vortices around a given point. It characterizes the rotational behavior of a vector field. If the curl at a point is non-zero, it signifies the presence of swirling motion or rotation in the vector field. The curl is associated with the circulation of a vector field around a closed loop and is often encountered in the study of fluid dynamics and electromagnetism.
\begin{align}
    \Vec{\nabla}\times\left(\text{Vector Field}\right)&\to\text{Vector},\\
    \text{where }\Vec{\nabla}\times\left(\text{Vector Field}\right)>0&\Rightarrow\text{Counterclockwise rotation}\notag\\
    \Vec{\nabla}\times\left(\text{Vector Field}\right)<0&\Rightarrow\text{Clockwise rotation}\notag\\
    \Vec{\nabla}\times\left(\text{Vector Field}\right)=0&\Rightarrow\text{No rotation}\notag
\end{align}
The curl captures the rotational aspects of a vector field, shedding light on phenomena such as fluid vortices or the behavior of magnetic fields around current-carrying wires.
\section*{Types of Integrals and Vector Calculus}
Line, surface, and volume integrals are mathematical tools for calculating physical quantities associated with one-dimensional lines, two-dimensional surfaces, and three-dimensional volumes. 
\thispagestyle{empty}
\subsection*{Line Integrals} 
They are also known as \textbf{path integrals}. They are used to calculate quantities along curves or paths in space. They are typically used to determine quantities associated with a line, such as work, circulation, or flux.

A line integral can be visualized as \textbf{the addition of} the contributions of a quantity along a curve. Think of walking along a curved path and measuring some physical quantity as you move. The line integral accumulates the contributions of the quantity along the entire path, taking into account the direction and magnitude at each point.
\subsection*{Surface Integrals}
They are used to calculate quantities that are distributed over two-dimensional surfaces. They are commonly employed to determine flux, representing the flow of a vector quantity through a surface.

Imagine a fluid flowing through a surface to grasp the physical intuition behind surface integrals. The surface integral calculates the total amount of fluid passing through the surface, considering the direction and magnitude of the flow at each point on the surface. It measures how much vector quantity passes through the surface per unit area.
\subsection*{Volume Integrals} 
They are used to calculate quantities associated with three-dimensional volumes. They are frequently used to determine quantities like total mass, charge, energy, or electric current.

Physically, volume integrals can be understood as \textbf{the summation of} the contributions of a quantity throughout a three-dimensional region. Picture a quantity distributed within a solid object. The volume integral integrates the contributions of the quantity over the entire volume, considering the variation of the quantity within the region.
\subsection*{Stoke's Theorem}
This theorem relates \textbf{the circulation of a vector field around a closed curve to the flux of the curl of the vector field through the surface bounded by the curve}. It establishes \textbf{a connection between a line integral around a closed curve and a surface integral over the region bounded by that curve}.
\begin{figure}[H]
    \centering
    \includegraphics[width=.82\linewidth]{figs/stokes-theorem.pdf}
    \caption{Stoke's theorem states the contribution of a vector field through a surface measured across the boundary line of that surface is the same if it were measured using the surface integral of the curl of the vector field.}
    \label{fig:stokes-theorem}
\end{figure}
Mathematically, this theorem is stated as follows:
\begin{align}
    \oint_C \vec{V}\cdot d\vec{l}=\iint_S (\Vec{\nabla}\times\vec{V})\cdot d\vec{S},
\end{align}
where $\displaystyle\oint_C$ represents the line integral around the closed curve $C$, $V$ is the vector field, $dr$ is the differential vector along the curve, $\displaystyle\iint_S$ represents the surface integral over the surface $S$ bounded by the curve, curl $\displaystyle\Vec{\nabla}\times\vec{V}$ is the curl of the vector field $V$, and $dS$ is the differential surface area vector.
\subsection*{Divergence Theorem}
\thispagestyle{empty}
Also known as \textbf{Gauss' theorem}, it relates \textbf{the flux of a vector field through a closed surface to the divergence of the vector field within the volume enclosed by that surface}. It establishes a \textbf{connection between a surface integral over a closed surface and a volume integral over the region bounded by that surface}.
\begin{figure}[H]
    \centering
    \includegraphics[width=.84\linewidth]{figs/divergence-theorem.pdf}
    \caption{Divergence theorem states the contribution of a vector field through a surface measured using the surface integral is the same if it were measured using the volume integral of the divergence of the vector field.}
    \label{fig:divergence-theorem}
\end{figure}
Mathematically, this theorem is stated as follows:
\begin{align}
    \iint_S \vec{V}\cdot d\vec{S}=\iiint_V (\Vec{\nabla}\cdot\vec{V})\cdot dV,
\end{align}
Where $\displaystyle\iint_S$ represents the surface integral over the closed surface $S$, $V$ is the vector field, $dS$ is the differential surface area vector, $\displaystyle\iiint_V$ represents the volume integral over the region $V$ enclosed by the surface, $\displaystyle\Vec{\nabla}\cdot\vec{V}$ is the divergence of the vector field $V$, and $dV$ is the differential volume element.
\thispagestyle{empty}
\clearpage
%%%%%%%%%%%%%%%%%%%%%%%%%%%%%%%%%%%%%%%%%%%%%%%%%%%%%%%%%%%%%%%%%%%%%
\bigskip \
\vspace{10cm}
\vskip 20pt
\begin{center}
    \thispagestyle{empty}
    \huge{\selectfont \bf \sffamily Part I: $\boldsymbol{\vec{E}}$-fields and all that}
\end{center}
\addcontentsline{toc}{section}{\Large Part I: $\boldsymbol{\vec{E}}$-fields and all that}
\vfill
\clearpage
%%%%%%%%%%%%%%%%%%%%%%%%%%%%%%%%%%%%%%%%%%%%%%%%%%%%%%%%%%%%%%%%%%%%%
\section{Electrostatics}
Electrostatics is a branch of physics that studies electric charges at rest and their interactions. This also includes the study of charge distribution that stays static in time. It focuses on the behavior of stationary electric charges and the electric fields they create.
\thispagestyle{empty}
\subsection{Electric Charge}
Electric charge is a \textbf{fundamental property of matter} that determines how objects interact through electromagnetic forces. It is a property carried by particles such as protons and electrons, which are the building blocks of atoms.

At its core, electric charge measures the imbalance of positive and negative particles within an object. It is denoted by the symbol $q$ and is typically measured in coulombs (C) units. There are two types of electric charges: positive and negative.

\textsc{Note}: \textbf{The magnitude of the charge of the electron or proton is a natural unit of charge}.
\subsection*{What Makes Positive Charges \textit{positive} and Negative Charges \textit{negative}?}
Electric charges determine the ability of elementary particles to experience electromagnetic forces. Objects with like charges (either positive or negative) repel each other, while objects with opposite charges attract each other. This fundamental force of attraction and repulsion is responsible for various phenomena we encounter in everyday life, from the spark produced by static electricity to the operation of electronic devices. \textbf{That's it!!! The signs of electric charges are just a convention}. No divine intervention here.
\subsection{Properties of Electric Charge}
\begin{itemize}
    \item Conservation of Charge: Charge is conserved, meaning it cannot be created or destroyed. The total amount (algebraic sum) of charge in an isolated system remains constant. This principle allows us to analyze charge interactions and predict the outcomes of various electrical processes.
    \item Quantization of Charge: Electric charge comes in discrete units. An electron carries the smallest unit of charge and is equal to approximately $-1.602 \times 10^{-19}$ coulombs. This quantization explains why we observe whole number multiples of the electron's charge in most charge interactions.
    \item Additivity of Charge: The total charge of an object is the algebraic sum of the individual charges within it. Charges can be added or subtracted, leading to a net charge for the object. For example, if we have two objects with charges +2C and -1C, the net charge of the system is +1C.
    \item Charge Interactions: Charged objects exert forces on each other through electric fields. Like charges repel each other, while opposite charges attract. These interactions are fundamental to understanding the behavior of charged particles and the formation of electric fields.
    \item Coulomb's Law: Coulomb's Law describes the force between two charged objects. It states that the force between two charges is directly proportional to the product of their magnitudes and inversely proportional to the square of the distance between them. This mathematical relationship helps quantify the strength of the electrostatic force.
\end{itemize}
\section{Electric Charge Distribution}
A \textit{point charge}, denoted $q$, is an idealized charged object with negligible physical size. Any system larger than a single point charge is called a \textit{charge distribution}. These are usually classified as either \textit{discrete} and/or \textit{continuous} distributions.
\subsection{Discrete Charge Distribution}
A discrete charge distribution is modeled as a collection of individual point charges. The total charge of such a distribution is the algebraic sum of individual point charges.
\begin{align}
    Q=\sum_i^Nq_i.
\end{align}
\begin{figure}[H]
    \centering
    \includegraphics[width=.78\linewidth]{figs/charge-distribution.pdf}
    \caption{(a) Discrete charge distribution. The total charge is the sum of individual point charges. (b) Continuous charge distributions. The total charge is found by integrating over the charge elements, $dq$.}
    \label{fig:charge-distribution}
\end{figure}
\subsection{Continuous Charge Distribution}
A continuous charge distribution is modeled as charge spread \textit{continuously} over a spatial region. We may split the system into infinitesimal elements of charge, $dq$, and treat these elements as point charges. The total charge is then found by integrating over the elements:
\begin{align}
    Q=\bigintsss dq.
\end{align}
\subsection*{Charge Densities}
While dealing with a continuous charge distribution, we can have the distribution in all three possible spatial dimensions. It is convenient to express the element $sq$ in terms of \textit{charge density}, then.

A line charge $\lambda$ [\unit{\coulomb\per\meter}], is a charge per unit length:
\begin{align}
    \lambda=\frac{dq}{dl}.
\end{align}
A surface charge $\sigma$ [\unit{\coulomb\per\meter\squared}], is a charge per unit area:
\begin{align}
    \sigma=\frac{dq}{da}.
\end{align}
A volume charge $\rho$ [$\unit{\coulomb\meter}^{-3}$], is a charge per unit volume:
\begin{align}
    \rho=\frac{dq}{dV}.
\end{align}
Elements of charge and the total charge of a continuous distribution are thus written as:
\begin{align}
    \text{distribution}\equiv
    \begin{cases}
      \displaystyle \frac{dq}{dl} & \text{in one dimension (1D)}\\[10pt]
      \displaystyle \frac{dq}{da} & \text{in two dimension (2D)}\\[10pt]
      \displaystyle \frac{dq}{dV} & \text{in three dimension (3D)}
    \end{cases} 
\end{align}
\begin{align}
    Q=\int dq=
    \begin{cases}
      \lambda dl & \text{in one dimension (1D)}\\[5pt]
      \sigma da & \text{in two dimension (2D)}\\[5pt]
      \rho dV & \text{in three dimension (3D)}
    \end{cases} 
\end{align}
\section{Coulomb's Law of Electrostatics}
Coulomb's Law is a fundamental principle in physics that describes the force between two electrically charged particles. It states that \textbf{the magnitude of the electrostatic force between two charged particles is directly proportional to the product of their charges and inversely proportional to the square of the distance between them}.
\begin{figure}[H]
    \centering
    \includegraphics[width=.7\linewidth]{figs/coulomb-force.pdf}
    \caption{Electrostatic force $\vec{F}_E$ between two point charges $q_1$ and $q_2$. This force can be attractive or repulsive depending on the signs of the point charges.}
    \label{fig:coulomb-force}
\end{figure}
Intuitively, Coulomb's Law can be understood as follows: If you have two charged objects, whether they are positively or negatively charged, they will either attract or repel each other. The strength of this attraction or repulsion depends on the magnitude of the charges and the distance between them. Like charges (positive-positive or negative-negative) repel each other, while opposite charges (positive-negative) attract each other.

Mathematically, Coulomb's Law (Only magnitude) can be expressed as:
\begin{align}
    F_E = \frac{1}{4\pi\epsilon_0}\frac{|q_1q_2|}{r^2},
\end{align}
where $F$ is the magnitude of the electrostatic force between the charges $q_1$ and $q_2$, $r$ is the distance between them, and $\displaystyle k=\frac{1}{4\pi\epsilon_0}=8.998\times10^9\,\unit{\newton\meter\squared}\text{C}^{-2}$ is the electrostatic constant.

In vector form, the equation reduces to
\begin{align}
    \Vec{F}_E&=\frac{1}{4\pi\epsilon_0}\frac{q_1q_2}{r^3}\Vec{r}.\\
    &=\frac{1}{4\pi\epsilon_0}\frac{q_1q_2}{r^2}\hat{r}
\end{align}
\subsection*{Comparing Electrostatic Force with Gravitational Force}
Coulomb's Law has a similar form to Newton's Law of gravitation, which gives the gravitational force, $\vec{F}_G$, exerted on a mass $m_2$ by a mass $m_1$ a distance $r$ apart:
\begin{align}
    \vec{F}_G=\frac{Gm_1m_2}{r^2}\hat{r}.
\end{align}
Since mass is always positive, the force of gravity is always attractive. Electric charges can be positive or negative, resulting in repulsive or attractive electrostatic forces.

If we calculate the ratio of the magnitudes of the gravitational force to the electrostatic force, we get the following:
\begin{align}
    \left|\frac{\vec{F}_G}{\vec{F}_E}\right|=4\pi\epsilon_0G\left(\frac{m_q}{q}\right)^2.
\end{align}
We can get an approximation for this by setting the mass and charge of elementary particles like the proton or electron we get:
\begin{align}
    \left|\frac{\vec{F}_G}{\vec{F}_E}\right|\sim10^{-36}.
\end{align}
This means it is safe to ignore gravitational effects in many electrostatic setups.

\textsc{Takeaway}: Each and every one of us possesses such a strong force on an elementary level, keeping our existence together that trumps the force that binds the galaxy and stars together.
\section*{The Principle of Superposition for Coulomb Forces}
The total electric force on a charged particle with charge $Q$ is the vector sum of the individual forces exerted on that particle by other $n$ amounts of charged particles in the distribution.
\begin{figure}[H]
    \centering
    \includegraphics[width=.78\linewidth]{figs/force-superposition-charge-distribution.pdf}
    \caption{(a) Discrete charge distribution. The total force is the vector sum of individual forces. (b) Continuous charge distributions. The total force is found by integrating over the force elements, $d\vec{F}$.}
    \label{fig:charge-distribution}
\end{figure}
The total force on $Q$ is:
\begin{align}
    \vec{F}_Q&=\vec{F}_{1Q}+\vec{F}_{2Q}+\vec{F}_{3Q}+\ldots\,\,\ldots\vec{F}_{nQ}
\end{align}
This type of linear addition of individual components to determine the total applies to many so-called \textit{linear systems} and is referred to as the \textit{principle of superposition}. The total force on $Q$ is given by the sum of the $n$ individual Coulomb interactions over charges $q_i$:
\begin{align}
    \vec{F}_Q&=\frac{Qq_1}{4\pi\epsilon_0r_{1Q}^2}\hat{r}_{1Q}+\frac{Qq_2}{4\pi\epsilon_0r_{2Q}^2}\hat{r}_{2Q}+\frac{Qq_3}{4\pi\epsilon_0r_{3Q}^2}\hat{r}_{3Q}+\ldots\,\,\ldots\frac{Qq_n}{4\pi\epsilon_0r_{nQ}^2}\hat{r}_{nQ}\notag\\
    &=\sum_i^N\frac{Qq}{4\pi\epsilon_0r_{iQ}^2}\hat{r}_{iQ},
\end{align}
where $\vec{r}_{iQ}=\vec{r}_Q-\vec{r}_i$. $\vec{r}_Q$ is the position vector of the observer charge. $\vec{r}_i$ is the position vector of the $i^{\text{th}}$ charge in the distribution, measured in the Cartesian coordinate system.

For a continuous distribution, the sum becomes an integral over elements of charge $dq$:
\begin{align}
    \vec{F}_Q=\bigintss\frac{Qq}{4\pi\epsilon_0r_{dqQ}^2}\hat{r}_{dqQ},
\end{align}
where $\vec{r}_{dqQ}=\vec{r}_Q-\vec{r}_{dq}$. $\vec{r}_Q$ is the position vector of the observer charge. $\vec{r}_{dq}$ is the position vector of the charge element $dq$ in the distribution, measured in the Cartesian coordinate system.

\subsection{Electric Field Intensity/Strength}
Electric field intensity, also known as \textbf{electric field strength} or simply \textbf{electric field}, is a fundamental concept in physics that describes the force experienced by a positive test charge placed in an electric field. It is defined as the force per unit positive charge exerted on the test charge.

Mathematically, electric field intensity $E$ at a point in space is given by the equation:
\begin{align}
    \vec{E}=\lim_{q_0\to 0}\frac{\vec{F}}{q_0}=\frac{1}{4\pi\epsilon_0}\frac{|q|}{r^2},
\end{align}
where $E$ is the electric field intensity, $F$ is the electrostatic force experienced by the test charge, and $q$ is the magnitude of the test charge. The limit is there to ensure the existence of the field even in the absence of $Q$,
\begin{figure}[H]
    \centering
    \includegraphics[width=.85\linewidth]{figs/electric-field.pdf}
    \caption{The electric field $\vec{E}$ felt by a point charge $q_0$ by an isolated point charge $q$.}
    \label{fig:coulomb-force}
\end{figure}
In vector form, the equation reduces to
\begin{align}
    \Vec{E}&=\frac{1}{4\pi\epsilon_0}\frac{q}{r^3}\Vec{r}.\\
    &=\frac{1}{4\pi\epsilon_0}\frac{q}{r^2}\hat{r}
\end{align}
\textsc{Note}:\textbf{ The electric field at a point does not depend on the incoming charge at \textit{that} point}. The electric field at a certain point is present whether or not there is a second point charge to experience the field.

Electric field intensity is a vector quantity with both magnitude and direction. The direction of the electric field at a given point is \textbf{the direction in which a positive test charge would experience a force if placed at that point}.
Another crucial relation is as follows:
\begin{align}
    \vec{F}_E=q\vec{E},
\end{align}
where $\Vec{F}_E$ is the Coulomb force felt by the charge $q$.
\subsection*{The Superposition Principle for Electric Fields}
Like the Coulomb force, the electric field also obeys the principle of superposition. The electric field experienced by a point charge $Q$ due to a charge distribution $q_i$ or $dq$ is given by:
\begin{align}
    \vec{E}=\frac{\vec{F}}{Q}&=\frac{q_1}{4\pi\epsilon_0r_{1r}^2}\hat{r}_{1r}+\frac{q_2}{4\pi\epsilon_0r_{2r}^2}\hat{r}_{2r}+\frac{q_3}{4\pi\epsilon_0r_{3r}^2}\hat{r}_{3r}+\ldots\,\,\ldots\frac{q_n}{4\pi\epsilon_0r_{nr}^2}\hat{r}_{nr}\notag\\
    &=\sum_i^N\frac{q_n}{4\pi\epsilon_0r_{ir}^2}\hat{r}_{ir},
\end{align}
where $\vec{r}_{ir}=\vec{r}-\vec{r}_i$ are the separation coordinates. $\vec{r}$ is the position vector of the point where the observation is being made. $\vec{r}_i$ is the position vector of the $i^{\text{th}}$ source charge, measured in the Cartesian coordinate system.

We can translate this to a continuous source charge distribution:
\begin{align}
    \vec{E}=\bigintss\frac{dq}{4\pi\epsilon_0r_{dqr}^2}\hat{r}_{dqr},
\end{align}
where $\vec{r}_{dqr}=\vec{r}-\vec{r}_{dq}$ are the separation coordinates. $\vec{r}$ is the position vector of the point where the observation is being made. $\vec{r}_{dq}$ is the position vector of the $dq$ source charge element, measured in the Cartesian coordinate system.
\subsection{Electric Field Lines}
They are imaginary lines used to visualize and represent the direction and strength of an electric field. They are drawn \textbf{tangent to the electric field vector} at each point and point away from positive charges and towards negative charges. 
\begin{figure}[H]
    \centering
    \includegraphics[width=.78\linewidth]{figs/electric-field-lines.pdf}
    \caption{The direction of the electric field at any point is tangent to the field line through that point.}
    \label{fig:electric-field-lines}
\end{figure}
\textsc{Note}: \textbf{Electric field lines are not trajectories. No two electric field lines will intersect.}

\textbf{Electric field lines provide a qualitative representation of the electric field's behavior and direction}. The density of field lines indicates the strength of the electric field, with closer lines representing a stronger field. The direction of the field lines shows the direction in which a positive test charge would move if placed in the field.

\subsection*{Why can't two electric field lines intersect each other?}
If two different field lines were to intersect, it would mean that a test charge would feel two different fields pointing in two different directions, which is forbidden.
\begin{figure}[H]
    \centering
    \includegraphics[width=\linewidth]{figs/filed-lines-charges.pdf}
    \caption{Electric field lines for two different charge distributions. The magnitude of $\vec{E}$ generally differs at different points along a given field line.}
    \label{fig:electric-field-lines}
\end{figure}
\subsection{Electric Field due to a Charged Particle}
\begin{figure}[H]
    \centering
    \includegraphics[width=.86\linewidth]{figs/electric-field-charges.pdf}
    \caption{A point charge $q$ produces an electric field $\vec{E}$ at all points in space. The field strength decreases with increasing distance.}
    \label{fig:electric-field-charges}
\end{figure}
\begin{align}
    \Vec{E}=\frac{1}{4\pi\epsilon_0}\frac{|q|}{r^2}\hat{r}.
\end{align}
We write $|q|$ to avoid the danger of getting a negative $E$ when $q$ is negative and then think the negative sign has something to do with direction.
\section{Electric Dipole}
An electric dipole is a system consisting of two equal and opposite point charges separated by a small distance. The charges in an electric dipole have opposite signs, typically denoted as $q=+ne^-$ and $q=-ne^-$, and are located close to each other. The separation between the charges is denoted as $d=2a$, where $a$ represents half the distance between the charges, which is quite essential in the calculation.
\begin{figure}[H]
    \centering
    \includegraphics[width=\linewidth]{figs/dipole-moment-def.pdf}
    \caption{(Left) The dipole moment $\vec{p}$ of the dipole points from the negative charge to the positive charge. (Right) The pattern of electric field lines around an electric dipole.}
    \label{fig:dipole-moment-def}
\end{figure}
\textsc{Note}: \textbf{The reason to take the distance between the point charges to be $2a$ is to avoid a non-fractional mess in the calculation}.

Electric dipoles also exhibit an electric dipole moment vector $\Vec{p}$, which points from the negative to the positive charge. The direction of the dipole moment indicates the orientation of the dipole. The magnitude of the dipole moment measures the strength of the dipole. Mathematically, it can be expressed as:
\begin{align}
    \vec{p}=q\cdot2\vec{a}.
\end{align}
\textsc{Note}: \textbf{While point charges have a net charge and produce an electric field around them, electric dipoles have a zero net charge since the charges cancel each other out}.

There are no theoretical limits on the separation distances of the two charges in an electric dipole. The smallest meaningful separation between charges in a physical system will depend on the scale of the particles involved and the conditions under which the dipole operates. We can set a limit $r\gg a$ for practical purposes, where $r$ is the distance from the common center of the dipole charges.
\subsection{Force on an Electric Dipole}
When placed in an electric field $\vec{E}$, each charge experiences a force. The torque on the dipole is the sum of the torques on the two charges.
\begin{figure}[H]
    \centering
    \includegraphics[width=.8\linewidth]{figs/torque-electric-dipole.pdf}
    \caption{(a) The positive charge experiences the force $F_+$, along $\vec{E}$. The torque points inward with $r_+$ being the moment arm. (b) The negative charge experiences the force $F_-$, opposite to $\vec{E}$. The torque still points inward, with $r_-$ being the moment arm. (c) The net force on this electric dipole is zero, but a torque directed into the page tends to rotate the dipole clockwise.}
    \label{fig:torque-electric-dipole}
\end{figure}
The positive charge experiences a Coulomb force $\vec{F}_+=+Q\vec{E}$ that points in the field direction. The negative charge of the dipole also feels an equal but opposite force $\vec{F}_-=-Q\vec{E}$ to that of $Q_+$. The net force on the dipole thus $\vec{F}_\text{net}=\vec{F}_++\vec{F}_-=0$. This does not mean the dipole is motionless. 
\subsection{Torque on an Electric Dipole}
The moment arm vector $\vec{r}_+=a\hat{d}$ for the positive charge points from the axis of rotation (the midpoint of the dipole) to the positive charge, perpendicular to the separation vector $d=2a$. The torque for the positive charge would then be:
\begin{align}
    \tau_+=\vec{r}_+\times\vec{F}_+=a\hat{d}\times Q\vec{E}=(Qa)\hat{d}\times\vec{E}=\frac{1}{2}\left(\vec{p}\times\vec{E}\right).
\end{align}
Similarly, the negative charge will feel a torque of the form, $\displaystyle\tau_-=\vec{r}_-\times\vec{F}_-=-L\hat{d}\times -Q\vec{E}=(Qa)\hat{d}\times\vec{E}=\frac{1}{2}\left(\vec{p}\times\vec{E}\right)$. The total torque of the dipole system is the vector sum of the two torque.
\begin{align}
    \vec{\tau}=\vec{\tau}_++\vec{\tau}_-=\frac{1}{2}\left(\vec{p}\times\vec{E}\right)+\frac{1}{2}\left(\vec{p}\times\vec{E}\right)=\vec{p}\times\vec{E}.
\end{align}
When placed in an external electric field, an electric dipole experiences a torque that tends to align it with the electric field. The torque exerted on the dipole is proportional to the product of the electric field $\vec{E}$ and the dipole moment $\Vec{p}$. This torque tends to rotate the dipole such that its dipole moment aligns with the electric field.
\subsection{Potential Energy in an Electric Dipole}
Since the dipole system experiences a net torque $\vec{\tau}\,[\unit{\newton\meter}]$ that rotates the dipole in an electric field, it gains some rotational potential energy. 
\begin{align}
    U_\text{dipole}&=-\bigintss\tau d\theta\notag\\
    &=-\bigints_{\displaystyle\,\,\frac{\pi}{2}}^{\displaystyle\,\,\theta}pE\sin\theta\,\,d\theta\notag\\
    &=\bigg[pE\cos\theta\bigg]_{\displaystyle\,\,\frac{\pi}{2}}^{\displaystyle\,\,\theta}\notag\\
    &=\vec{p}\cdot\vec{E}
\end{align}
\section{Electric Field Intensity Measurements for Charge Distribuions}
\subsection{Discrete I: Electric Field due to an Electric Dipole (Parallel on the Dipole Axis)}
\begin{figure}[H]
    \centering
    \includegraphics[width=.85\linewidth]{figs/electric-field-dipole.pdf}
    \caption{An electric dipole. The dipole's two charges result from the electric field vectors at point $P$ on the dipole axis. Point $P$ is at distances $r_+=x-a$ and $r_-=x+a$ from the individual charges that make up the dipole. Here, $x\gg a$.}
    \label{fig:electric-field-dipole}
\end{figure}
Measuring the field $x$ distance away from the bisecting point of the dipole:
\begin{align}
    \vec{E}_\text{dipole}&=\vec{E}_++\vec{E}_-\notag\\
    &=\frac{1}{4\pi\epsilon_0}\left(-\frac{q}{r_+^2}+\frac{q}{r_-^2}\right)\hat{i}\notag\\
    &=\frac{q}{4\pi\epsilon_0}\left\{\frac{1}{(x-a)^2}-\frac{1}{(x+a)^2}\right\}\hat{i}\notag\\
    &=\frac{q}{4\pi\epsilon_0}\left[\frac{(x-a)^2-(x+a)^2}{\left\{(x+a)(x-a)\right\}^2}\right]\hat{i}\notag\\
    &=\frac{q}{4\pi\epsilon_0}\left[\frac{4xa}{(x^2-a^2)^2}\right]\hat{i}\notag\\
    &=\frac{q}{4\pi\epsilon_0}\left[\frac{2\times2xa}{x^4}\right]\hat{i}\notag\\
    &=\frac{1}{4\pi\epsilon_0}\times2\times\frac{q\times 2a}{x^3}\hat{i}\notag\\
    &=\frac{1}{4\pi\epsilon_0}\frac{2{p}}{x^3}\hat{i}\notag\\
    &=\frac{1}{2\pi\epsilon_0}\frac{{p}}{x^3}\hat{i}.
\end{align}
The field points from the direction of the dipole moment. i.e., toward the positive charge.

To account for the arbitrary direction of the observation point $P$ along the dipole axis. We replace $x$ with $r$, giving us:
\begin{align}
    \vec{E}_\text{dipole}=\frac{1}{2\pi\epsilon_0}\frac{\vec{p}}{r^3}\hat{r}_{(-)\to(+)}=\frac{1}{2\pi\epsilon_0}\frac{\vec{p}}{r^3}\hat{p}\label{eq:field-on-the-dipole-axis}
\end{align}
The electric field of a dipole thus points in the direction of the electric dipole moment, regardless of whatever orientation we try. It further establishes the \textit{why} behind the convention of $\vec{p}$'s direction.
\subsection{Discrete II: Electric Field due to an Electric Dipole (Perpendicular to the Dipole Axis)}
\begin{figure}[H]
    \centering
    \includegraphics[width=.85\linewidth]{figs/electric-field-dipole-off-the-axis.pdf}
    \caption{An electric dipole produces an electric field at point $P$ off the dipole axis. Point $P$ is at a distance $r$ perpendicularly off the dipole axis. Here, $r\gg a$.}
    \label{fig:electric-field-dipole}
\end{figure}
Measuring the field $r$ distance, perpendicularly away from the bisecting point of the dipole:
\begin{align}
    \vec{E}_\text{dipole}&=\vec{E}_++\vec{E}_-\notag\\
    &=\left\{E_+(x)+E_-(x)\right\}\hat{i}+\left\{E_+(y)+E_-(y)\right\}\hat{j}\notag\\
    &=\left\{E_+(x)+E_-(x)\right\}\cos\alpha\hat{i}+0\hat{j}\notag\\
    &=\frac{2}{4\pi\epsilon_0}\times\frac{q}{y^2+a^2}\times\frac{a}{\sqrt{y^2+a^2}}\hat{i}\notag\\
    &=\frac{q\times 2a}{4\pi\epsilon_0}\frac{1}{\left(y^2+a^2\right)^{\frac{3}{2}}}\hat{i}\notag\\
    &=\frac{1}{4\pi\epsilon_0}\frac{{p}}{\left(y^2+a^2\right)^{\frac{3}{2}}}\left\{-\hat{p}\right\}.
\end{align}
The field points from the direction of the dipole moment. i.e., toward the positive charge.

To account for the arbitrary direction of the observation point $P$ along the dipole axis. We replace $y$ with $r$, giving us:
\begin{align}
    \vec{E}_\text{dipole}=\frac{1}{4\pi\epsilon_0}\frac{{p}}{\left(r^2+a^2\right)^{\frac{3}{2}}}\left\{-\hat{p}\right\}
\end{align}
In the limit $r\gg a$, we get the following:
\begin{align}
    \vec{E}_\text{dipole}=\frac{1}{4\pi\epsilon_0}\frac{{p}}{r^3}\left\{-\hat{p}\right\}.\label{eq:field-perp-to-the-dipole-axis}
\end{align}
We can see that Eqs. (\ref{eq:field-perp-to-the-dipole-axis}) and (\ref{eq:field-on-the-dipole-axis}) are identical except for a factor of two, which shows up only when the field is measured along the dipole axis.

In both cases, it is evident that the electric field due to an electric dipole falls off as $\displaystyle\frac{1}{r^3}$. This is called the \textit{Inverse Cube Law}, unlike the well-known \textit{Inverse Square Law} we encountered for a single-point charge.
\subsection{Continuous I: Electric Field due to a Charged Line (Finite) Segment}
A line segment, of length $2L$, is uniformly charged with a charge density $\displaystyle\lambda=\frac{Q}{2L}\,\unit{\coulomb\per\meter}$. We divide the charged line segment into infinitesimal point charge elements $\lambda dy$. 
\begin{figure}[H]
    \centering
    \includegraphics[width=.65\linewidth]{figs/line-segment-charge-elec-field.pdf}
    \caption{Measuring the field $x$ distance, \textit{perpendicularly} away from the center point of a line segment of length $2a$, uniformly charged with $+Q$ charge.}
    \label{fig:line-segment-charge-elec-field}
\end{figure}
At $y=+L$, we can use one of these elements and measure the infinitesimal field $d\vec{E}$ produced at point $P$. We then resolve this $d\vec{E}$ into its spatial components $dE_x$ and $dE_y$. This line element has an identical twin on the opposite end of the segment at $y=-L$. This element also produces a $d\vec{E}$ of its own, having $dE_x$ and $dE_y$. This twin point charge element cancels their $y$-components while adding $x$-components. We get a similar outcome if we exhaust all possible charge elements on this line. Each symmetric charge element gives a summed $x$-components and canceled $y$-components. We now integrate these infinitesimal field components between the limits of $y=-L$ to $y=+L$ and calculate the net $E$ at $P$.  
\begin{align}
    \vec{E} &= \bigintss d\vec{E} = \bigintss dE_x \hat{i} \\\notag
    &=\frac{1}{4\pi\epsilon_0}\bigintss\frac{\lambda dy}{\sqrt{x^2+y^2}}\cos\alpha \, \hat{i}\notag\\
    &=\frac{1}{4\pi\epsilon_0}\bigintss\frac{\lambda dy}{x^2+y^2}\times\frac{x}{\sqrt{x^2+y^2}}\,\hat{i}\notag\\
    &=\frac{\lambda x}{4\pi\epsilon_0}\bigintss\frac{ dy}{\left(x^2+y^2\right)^{3/2}}\,\hat{i}.\label{eq:line-charge-pause}
\end{align}
\textsc{Note}: We are taking the charge $Q$ to be positive in the following calculation. For the same case where $Q$ is taken to be negative, we flip the direction of the electric field at each point, even the final net direction, for that matter.

In Eq. (\ref{eq:line-charge-pause}), we substitute $y=x\tan{u}$ with $\displaystyle u=\tan^{-1}\left(\frac{y}{x}\right)$.
\begin{align*}
    y&=x\tan{u}\\
    dy&=x\sec^2{u}du.\\
    x^2+y^2&=x^2+\left(x\tan{u}\right)^2\\
    &=x^2\left(1+\tan^2{u}\right)\notag\\
    &=x^2\sec^2{u}.
\end{align*}
With these substitutions, Eq. (\ref{eq:line-charge-pause}) now becomes:
\begin{align}
    \vec{E}&=\frac{\lambda x}{4\pi\epsilon_0}\bigintss\frac{x\sec^2{u}}{\left(x^2\sec^2{u}\right)^{3/2}}du\notag\\
    &=\frac{\lambda x}{4\pi\epsilon_0}\bigintss\frac{x\sec^2{u}}{x^3\sec^3{u}}du\notag\\
    &=\frac{\lambda x}{4\pi\epsilon_0}\bigintss\frac{du}{x^2\sec{u}}\notag\\
    &=\frac{\lambda x}{4\pi\epsilon_0x^2}\bigintss\cos{u}{du}\notag
\end{align}
Here, we must apply the limits of integration for $u$ and not $y$ because we swapped the integrating variable $y$.
\begin{table}[H]
    \centering
    \begin{tabular}{P{1.0cm}|P{3.0cm}}
    $y$ & $\displaystyle u=\tan^{-1}\left(\frac{y}{x}\right)$ \\[8pt]\hline\hline\\
    $+L$ & $\displaystyle\tan^{-1}\left(\frac{+L}{x}\right)$\\[12pt]
    $-L$ & $\displaystyle\tan^{-1}\left(\frac{-L}{x}\right)$\\
\end{tabular}
\end{table}
We now evaluate the integral with the properly set limits:
\begin{align}
    \vec{E}&=\frac{\lambda x}{4\pi\epsilon_0x^2}\bigints_{\displaystyle\tan^{-1}\left(\frac{-L}{x}\right)}^{\displaystyle\tan^{-1}\left(\frac{+L}{x}\right)}\cos{u}{du}\,\hat{i}\notag\\[6pt]
    &=\frac{\lambda x}{4\pi\epsilon_0x^2}\Biggl[\sin{u}\Biggr]_{\displaystyle\tan^{-1}\left(\frac{-L}{x}\right)}^{\displaystyle\tan^{-1}\left(\frac{+L}{x}\right)}\,\hat{i}\notag\\[6pt]
    &=\frac{\lambda x}{4\pi\epsilon_0x^2}\left[\sin\left\{\tan^{-1}\left(\frac{L}{x}\right)\right\}-\sin\left\{\tan^{-1}\left(\frac{-L}{x}\right)\right\}\right]\,\hat{i}\notag\\
    &=\frac{\lambda x}{4\pi\epsilon_0x^2}\left[\sin\left\{\sin^{-1}\left(\frac{L}{\displaystyle x\sqrt{\frac{L^2}{x^2}+1}}\right)\right\}-\sin\left\{\sin^{-1}\left(\frac{-L}{\displaystyle x\sqrt{\frac{L^2}{x^2}+1}}\right)\right\}\right]\,\hat{i}\notag\\
    &=\frac{\lambda x}{4\pi\epsilon_0x^2}\left(\frac{L}{x\sqrt{\displaystyle\frac{L^2}{x^2}+1}}-\frac{-L}{\displaystyle x\sqrt{\frac{L^2}{x^2}+1}}\right)\,\hat{i}\notag\\
    &=\frac{\lambda x}{4\pi\epsilon_0x^2}\left(\frac{2L}{\displaystyle x\sqrt{\displaystyle\frac{L^2}{x^2}+1}}\right)\,\hat{i}\notag\\
    &=\frac{\lambda }{4\pi\epsilon_0x^2}\left(\frac{2L}{\displaystyle\sqrt{\frac{L^2+x^2}{x^2}}}\right)\notag\\
    &=\frac{\lambda\times2L}{4\pi\epsilon_0x\sqrt{L^2+x^2}}\,\hat{i}\\
    &=\dfrac{1}{4 \pi \epsilon_0} \dfrac{Q}{x\sqrt{x^2 + L^2}}\,\hat{i}.
\end{align}
For an observation made in any arbitrary orientation of the line charge segment:
\begin{align}
    \vec{E}&=\frac{\lambda\times2L}{4\pi\epsilon_0r\sqrt{r^2+L^2}}\,\hat{r}\\
    &=\dfrac{1}{4 \pi \epsilon_0} \dfrac{Q}{r\sqrt{r^2 + L^2}}\,\hat{r}.
\end{align}
$\hat{r}$ denotes the direction (radially outward) of the field (due to $+Q$) \textit{away} from the source \textit{toward} the observation point.

As promised, for the case of $-Q$, the field will point radially inward toward the source. The field strength will remain the same; only the direction will earn a negative sign before it to that of the positive case.

\textsc{Limiting Cases}: \begin{enumerate}
\itemsep0em 
    \item $r \gg L$: It reduces to the field due to a single-point charged particle. 
    \begin{align}
    \vec{E} = \dfrac{1}{4 \pi \epsilon_0} \dfrac{Q}{r^2}\,\hat{r}.
    \end{align}
    \item $r\to\infty$: 
    \begin{align}
    \vec{E} = 0.
    \end{align}
    \item $r=L$: 
    \begin{align}
    \vec{E} = \dfrac{1}{2 \pi \epsilon_0} \dfrac{Q}{r^2}\,\hat{r}.
    \end{align}
    \item $r\ll L$: This case gives us the Electric Field due to an Infinite Line of Charge.
    \begin{align}
        \vec{E} = \frac{1}{2\pi\epsilon_0}\frac{\lambda}{r}\hat{r}.\label{eq:infinite-line-charge-field}
    \end{align}
\end{enumerate}  


\subsection{Continuous II: Electric Field due to an Infinite Line of Charge}
\textsc{Note}: We are taking the charge $Q$ to be positive in the following calculation. For the same case where $Q$ is taken to be negative, we simply flip the direction of the electric field at each point, even the final net direction, for that matter.
\begin{figure}[H]
    \centering
    \includegraphics[width=.75\linewidth]{figs/infinite-line-charge-elec-field.pdf}
    \caption{Measuring the field $x$ distance away from an infinite line uniformly charged with $+Q$ charge.}
    \label{fig:infinite-line-charge-elec-field}
\end{figure}
The calculations for this are the same for the charged line segment. The integration limit should be $y=-\infty$ to $y=+\infty$ instead of $y=-L$ to $y=+L$. Here, we must apply the limits of integration for $u$ and not $y$ because we swapped the integrating variable $y$. This type of integration usually requires a splitting up of the limits.
\begin{table}[H]
    \centering
    \begin{tabular}{P{1.0cm}|P{3.0cm}}
    $y$ & $\displaystyle u=\tan^{-1}\left(\frac{y}{x}\right)$ \\[8pt]\hline\hline\\
    $+\infty$ & $\displaystyle+\frac{\pi}{2}$\\[10pt]
    $-\infty$ & $\displaystyle-\frac{\pi}{2}$\\[15pt]
    $0$ & $\displaystyle0$\\
\end{tabular}
\end{table}
We directly import the steps from the line segment field derivation (\ref{eq:line-charge-pause}) and evaluate the integral with the properly set limits:
\begin{align*} 
    \vec{E}&= \dfrac{1}{4 \pi \epsilon_0} \bigintss_{\displaystyle\,\,y=-\infty}^{\displaystyle\,\,y=0}\dfrac{\lambda xdy}{(x^2 + R^2)^{3/2}} \hat{i}+\dfrac{1}{4 \pi \epsilon_0} \bigintss_{\displaystyle\,\,y=0}^{\displaystyle\,\,y=+\infty}\dfrac{\lambda xdy}{(x^2 + R^2)^{3/2}} \hat{i} \\[4pt]
    &=\frac{\lambda x}{4\pi\epsilon_0x^2}\bigints_{\displaystyle\,\,-\frac{\pi}{2}}^{\displaystyle\,\,0}\cos{u}{du}\,\hat{i}+\frac{\lambda x}{4\pi\epsilon_0x^2}\bigints_{\displaystyle \,\,0}^{\displaystyle\,\,+\frac{\pi}{2}}\cos{u}{du}\,\hat{i} \\[4pt]
    &=\frac{\lambda}{4\pi\epsilon_0x}\left(\bigints_{\displaystyle-\frac{\pi}{2}}^{\displaystyle 0}\cos{u}{du}+\bigints_{\displaystyle \,\,0}^{\displaystyle+\frac{\pi}{2}}\cos{u}{du}\right)\hat{i}\notag\\[4pt]
    &=\dfrac{1}{4 \pi \epsilon_0}\frac{2\lambda}{x}\hat{i}.
\end{align*}
For an observation made in any arbitrary orientation of the line charge segment:
\begin{align}
    \vec{E}&=\dfrac{\lambda}{2\pi\epsilon_0r}\hat{r}.
\end{align}
$\hat{r}$ denotes the direction (perpendicularly outward) of the field (due to $+Q$) \textit{away} from the source \textit{toward} the observation point.

As promised, for the case of $-Q$, the field will point radially inward toward the source. The field strength will remain the same; only the direction will earn a negative sign in front of it to that of the positive case.

\textsc{Note}: We will have a \textit{given} line charge density $\lambda$ in this case. Unlike the previous cases where we could define $\displaystyle\lambda=\frac{\text{Total Charge}}{\text{Total Length}}$, we can't do the same here due to the infinite length. 
\subsection{Continuous III: Electric Field due to a Ring of Charge}
\begin{figure}[H]
    \centering
    \includegraphics[width=.75\linewidth]{figs/ring-charge-elec-field.pdf}
    \caption{Measuring the field $x$ distance away from the center of a ring of radius $a$, uniformly charged with $+Q$ charge.}
    \label{fig:ring-charge-elec-field}
\end{figure}
The ring, having radius $R$, is uniformly charged with a density $\displaystyle\lambda=\frac{Q}{2\pi R}\,[\unit{\coulomb\meter}^{-2}]$. We divide the charged ring into infinitesimal point charge elements $\lambda ds$ across the ring's circumference. 

\textsc{Note}: We are taking the charge $Q$ to be positive in the following calculation. For the same case where $Q$ is taken to be negative, we simply flip the direction of the electric field at each point, even the final net direction, for that matter.

At $s=0$ (suppose it denotes $y=+R$ in the diagram), we can use one of these elements and measure the infinitesimal field $d\vec{E}$ produced at point $P$. We then resolve this $d\vec{E}$ into its spatial components $dE_x$ and $dE_y$. This line element has an identical twin on the opposite end of the segment at $y=-R$ across the circumference. This element also produces a $d\vec{E}$ of its own, having $dE_x$ and $dE_y$. This twin point charge element cancels their $y$-components while adding $x$-components. We get a similar outcome if we exhaust all possible charge elements on this line. Each symmetric charge element gives a summed $x$-components and canceled $y$-components. We now integrate these infinitesimal field components between the limits of $s=0$ to $s=2\pi R$ and calculate the net $E$ at $P$. 
\begin{align}
    \vec{E} &= \bigintss d\vec{E} = \bigintss dE_x \hat{i} \\\notag
    &=\frac{1}{4\pi\epsilon_0}\bigints_{0}^{2\pi R}\frac{\lambda ds}{\sqrt{x^2+R^2}}\cos\alpha\,\hat{i}\notag\\
    &=\frac{1}{4\pi\epsilon_0}\bigints_{0}^{2\pi R}\frac{\lambda ds}{x^2+R^2}\times\frac{x}{\sqrt{x^2+R^2}}\,\hat{i}\notag\\
    &=\frac{\lambda x}{4\pi\epsilon_0\left(x^2 + R^2\right)^{3/2}}\bigints_{0}^{2\pi R}dS\,\hat{i}\notag\\
    &=\frac{2\pi\lambda Rx}{4\pi\epsilon_0\left(x^2 + R^2\right)^{3/2}}\,\hat{i}\notag\\
    &=\frac{\lambda Rx}{2\epsilon_0\left(x^2 + R^2\right)^{3/2}}\,\hat{i}.
\end{align}
For an observation made in any arbitrary orientation of the line charge segment:
\begin{align}
    \vec{E}&=\frac{\lambda Rr}{2\epsilon_0\left(r^2 + R^2\right)^{3/2}}\,\hat{r}.\label{eq:ring-charge-field-def}\notag\\
    &=\frac{Qr}{4\pi\epsilon_0\left(r^2 + R^2\right)^{3/2}}\,\hat{r}.
\end{align}
\textsc{Limiting Cases}: \begin{enumerate}
\itemsep0em 
    \item $r\gg R$: 
    \begin{align}
    \vec{E} = \dfrac{1}{4 \pi \epsilon_0} \dfrac{Q}{r^2}\,\hat{r}.
    \end{align}
    This reduces to a field due to a single-point charge.
    \item $r=0$: 
    \begin{align}
    \vec{E} = 0.
    \end{align}
    The field is zero at the center of the ring.
    \item $r=R$: 
    \begin{align}
    \vec{E} = \frac{1}{\left(2\right)^{3/2}}\dfrac{1}{4 \pi \epsilon_0} \dfrac{Q}{r^2}\,\hat{r}.
    \end{align}
    \item $r\ll R$: 
    \begin{align}
    \vec{E} = \dfrac{1}{4 \pi \epsilon_0} \dfrac{Qr}{R^2}\,\hat{r}.
    \end{align}
    This reduces to a field due to a single-point charge.
\end{enumerate}
$\hat{r}$ denotes the direction (axially outward) of the field (due to $+Q$) \textit{away} from the source \textit{toward} the observation point.

As promised, for the case of $-Q$, the field will point radially inward toward the source. The field strength will remain the same; only the direction will earn a negative sign in front of it compared to that of the positive case.
\subsection{Continuous IV: Electric Field due to a Uniformly Charged Disk}
\textsc{Note}: We are taking the charge $Q$ to be positive in the following calculation. For the same case where $Q$ is taken to be negative, we simply flip the direction of the electric field at each point, even the final net direction for that matter.
\begin{figure}[H]
    \centering
    \includegraphics[width=.75\linewidth]{figs/solid-disk-charge-elec-field.pdf}
    \caption{Measuring the field $x$ distance away from the center point of a disk of radius $R$, uniformly charged with $+Q$ charge.}
    \label{fig:solid-disk-charge-elec-field}
\end{figure}
We divide the charged disk into a collection of infinitesimal ring charge elements. We already know the electric field produced by a ring charge distribution from Eq. (\ref{eq:ring-charge-field-def}). We start by taking one ring charge element of radius $r^{'}$ and integrate it from $r^{'}=0$ to $r^{'}=R$ to measure the net $\vec{E}$ produced at $P$. A ring charge produces a field along the central axis as per Eq. (\ref{eq:ring-charge-field-def}). In order to integrate the ring charge element, we need to vary it infinitesimally. Take $dQ=\sigma dA=2\pi\sigma r^{'}dr^{'}$, where $\sigma$ is the surface charge density of the disk, $2\pi r^{'}$ is the circumference of the initial ring charge, and $dr^{'}$ is the variation.
\begin{align}
    \vec{E} &= \frac{1}{4\pi\epsilon_0}\bigintss \frac{dQ}{x^2+r^{'2}}\times\frac{x}{\sqrt{x^2+r^{'2}}}\hat{i}\notag\\
    &= \frac{1}{4\pi\epsilon_0}\bigintss \frac{2\pi\sigma r^{'}dr^{'}}{\left(x^2+r^{'2}\right)^{\frac{3}{2}}}\hat{i}\notag\\
    &=\frac{x\sigma\pi}{4\pi\epsilon_0}\bigints_{r^{'}=0}^{r^{'}=R}\frac{2r^{'}dr^{'}}{\left(x^2+r^{'2}\right)^{\frac{3}{2}}}\hat{i}\notag\\
    &=\frac{\sigma x}{4\epsilon_0}\left[\frac{\left(x^2+r^{'2}\right)^{-\frac{1}{2}}}{-\frac{1}{2}}\right]_{r^{'}=0}^{r^{'}=R}\hat{i}\notag\\[12pt]
    \because\bigintss X^MdX&=\frac{X^{M+1}}{M+1}\,\text{with }X=x^2+r^{'2},\,m=-\frac{3}{2},\,dX=2r^{'}dr^{'}\notag\\[12pt]
    &=\frac{2\sigma x}{4\epsilon_0}\bigg[\frac{1}{x}-\frac{1}{\sqrt{x^2+R^2}}\bigg]\hat{i}\notag\\
    &=\frac{\sigma}{2\epsilon_0}\bigg(1-\frac{x}{\sqrt{x^2+R^2}}\bigg)\hat{r}.
\end{align}
For an observation made in any arbitrary orientation of the line charge segment:
\begin{align}
    \vec{E}&=\frac{\sigma}{2\epsilon_0}\bigg(1-\frac{r}{\sqrt{r^2+R^2}}\bigg)\hat{r}.\\
    &=\frac{Q}{2\epsilon_0\pi R^2}\bigg(1-\frac{r}{\sqrt{r^2+R^2}}\bigg)\hat{r}.
\end{align}
$\hat{r}$ denotes the direction (axially outward) of the field (due to $+Q$) \textit{away} from the source \textit{toward} the observation point.

As promised, for the case of $-Q$, the field will point radially inward toward the source. The field strength will remain the same; only the direction will earn a negative sign in front of it compared to that of the positive case.

\textsc{Limiting Cases}: \begin{enumerate}
\itemsep0em 
    \item $R\gg r$: This case reduces to a field due to an infinite sheet of charge. We will return to this important result later in the course. 
    \begin{align}
    \vec{E} = \frac{\sigma}{2\epsilon_0}.
    \end{align}
    \item $R\ll r$: Observing too far from the source essentially will result in zero field intensity.
    \begin{align}
    \vec{E} = 0.
    \end{align}
\end{enumerate}
\subsection{Continuous V: Electric Field due to an Infinite Sheet of Charge}
\begin{figure}[H]
    \centering
    \includegraphics[width=\linewidth]{figs/single-sheet-charge-elec-field.pdf}
    \caption{Measuring the field from a random distance away from an infinitely long sheet of charge.}
    \label{fig:single-sheet-charge-elec-field}
\end{figure}
\begin{align}
    \vec{E} = \frac{\sigma}{2\epsilon_0}.\label{eq:infty-sheet}
\end{align}
This field does not depend on the distance from the sheet. It is uniform everywhere perpendicular to the sheet, away from it.

\textsc{Note}: There is no such thing as an infinite sheet of charge, but if the dimensions of the sheet are much larger than the observation distance $r$ of the field from the sheet, the field is very nearly given by Eq. (\ref{eq:infty-sheet}).
\subsection{Continuous VI: Electric Field due to Two Oppositely Charged Infinite Sheets}
\begin{figure}[H]
    \centering
    \includegraphics[width=.75\linewidth]{figs/sheet-charge-elec-field.pdf}    \caption{Measuring the field in between and away from two infinite plane sheets with uniform surface charge densities $+\sigma$ and $-\sigma$ are placed parallel to each other with separation $d$.}
    \label{fig:sheet-charge-elec-field}
\end{figure}
\begin{align}
    \vec{E}=\vec{E}_1+\vec{E}_2=\left\{
                \begin{array}{ll}
                  0&\qquad\text{above the upper (positive) sheet}\\[5pt]
                  \displaystyle \frac{\sigma}{\epsilon_0}\hat{n}_{(+)\to(-)}&\qquad\text{between the sheet}\\[8pt]
                  0&\qquad\text{below the lower (negative) sheet}
                \end{array}
              \right.
\end{align}
Where both the sheets have the same magnitude at all points, independent of the distance from either sheet: 
\begin{align*}
    E_1=E_2=\frac{\sigma}{2\epsilon_0}.
\end{align*}
% \subsection{Electric Field Density}
% Electric flux density, also known as \textbf{electric displacement}, is a concept that extends the idea of electric field lines to include the effects of electric charges and electric polarization in materials. It measures how electric field lines originate or terminate on charges, including free charges and charges induced in the material.

% When we have electric charges, whether free or induced in a material, they create electric fields around them. These electric fields give rise to electric field lines that emanate from positive charges and terminate on negative charges. However, when we introduce materials into the picture, such as dielectrics or insulators, the presence of these materials can cause electric charges to become distributed or polarized within the material.

% Electric flux density considers both the free charges and the charges induced in the material. It represents the total electric displacement created by both types of charges. Considering the material's influence, it measures how the electric field lines originate or terminate on charges.
\section{Gauss's Law of Electrostatics}
Gauss's Law of Electrostatics is an alternative form of Coulomb's Law. The latter provides information about the electric field once the source electric charge distribution is given. Gauss's Law does the opposite. It provides information on the amount of charge, given the field. 

It states that \textbf{the total electric flux through any closed surface is equal to the net electric charge enclosed by that surface divided by the permittivity of free space}.
\subsection{Electric Flux}
It measures \textbf{the number of electric field lines passing through a given surface}. It is defined as the dot product of the electric field vector $E$ and the area vector $A$ of the surface.

Mathematically, electric flux $\Phi_E$ is given by the equation:
\begin{align}
    \Phi_E = \sum\text{Electric field lines crossing through a chosen area},
\end{align}
\begin{figure}[H]
    \centering
    \includegraphics[width=.7\linewidth]{figs/electric-flux-def.pdf}
    \caption{Electric field vectors and field lines pierce an imaginary, spherical Gaussian surface that encloses a particle with charge $+Q$ in (a) and $-Q$ in (b).}
    \label{fig:electric-flux-charges}
\end{figure}
Electric flux provides information about the \textit{strength} and \textit{distribution} of the electric field through a surface. It is \textbf{proportional to the number of electric field lines} passing through the surface, indicating the intensity of the electric field.

The flux is zero if the number of electric field lines piercing the surface equals the number of field lines piercing out of the surface. We can quickly get a head start on the flux problem just by drawing a relevant diagram. Depending on the angle between $\vec{E}$ and $\vec{A}$ or $d\vec{A}$, electric flux can be positive, zero, or negative. 
\subsection{Flux of a Uniform Electric Field}
\begin{figure}[H]
    \centering
    \includegraphics[width=.75\linewidth]{figs/electric-flux-def-with-area.pdf}
    \caption{(a) An electric field vector pierces a small square patch on a flat surface. (b) Only the $x$ component pierces the patch; the $y$ component skims across it. (c) The area vector of the patch is perpendicular to the patch, with a magnitude equal to the patch’s area.}
    \label{fig:electric-flux-area}
\end{figure}
The flux is mathematically calculated as the total contribution of the electric field across the Gaussian surface area, precisely given by the following dot product.
\begin{align}
    \Phi_E&=\vec{E}\cdot\Vec{A}=E_xA_x+E_yA_y+E_zA_z\\
    &=E(A\cos\theta)\notag\\
    &=A(E\cos\theta)\notag.
\end{align}
Just like the notion that we use in dot product, the parallel projection of $\vec{E}$ onto $\vec{A}$ gives the same flux as the projection of $\vec{A}$ onto $\vec{E}$. Here, $\vec{E}$ is uniform across the Gaussian surface area.
\subsection{Flux of a Non-Uniform Electric Field}
In this case, we cannot take the same field $\vec{E}$ across a common cross-sectional area. We must measure the field flux across all possible differential area elements and then integrate all contributions.
\begin{align}
    \Phi_E&=\int E\cos\theta dA\\
    &=\int E_\parallel\, dA\notag\\
    &=\int\vec{E}\cdot d\vec{A}
\end{align}
The differential element is defined arbitrarily across the Gaussian surface. We take all the possible combinations of $\vec{E}\cdot d\vec{A}$ and then integrate them through the whole area.
\begin{figure}[H]
    \centering
    \includegraphics[width=.8\linewidth]{figs/gauss's-law.pdf}
    \caption{A Gaussian surface of arbitrary shape immersed in an electric field. The surface is divided into small squares of area $A$. The electric field vectors and the area vectors for three representative squares, marked (a), (b), and (c), are shown.}
    \label{fig:gauss-law}
\end{figure}
Mathematically, Gauss's Law can be expressed as:
\begin{align}
    \oint \Vec{E}\cdot d\Vec{a} &= \frac{Q_\text{enclosed}}{\varepsilon_0}.\qquad\text{(Integral Form)}\label{gauss-integral}
\end{align}
This Law demands that \textbf{the total electric flux passing through a closed surface is proportional to the total electric charge enclosed within that surface}. In other words, it relates the behavior of the electric field to the distribution of electric charge.

\textsc{Intrigue}: The prime reason for choosing a specific Gaussian Surface comes once we define a term called \textit{Equipotential Surface}. It is to be introduced soon with proper context. Stand by!

Gauss's Law's ability to provide a simplified method for calculating electric fields in highly symmetrical situations is critical. The electric flux and the enclosed charge can be determined by choosing an \textit{appropriate} Gaussian surface that \textit{matches the symmetry of the problem}. This allows for calculating the electric field at any point of interest without considering the details of the charge distribution.
\subsection{Applications of Gauss's Law}
\subsubsection*{Spherical Symmetry}
We want to measure $\vec{E}$ in three regions divided by the limits: $r<R, r=R, r>R$, where $r$ is the radius of the chosen Gaussian surface, $R$ is the radius of the spherical charge distribution or the spherical shell charge distribution, both uniformly charged with an amount $Q$.
\begin{figure}[H]
    \centering
    \includegraphics[width=.67\linewidth]{figs/gauss-spherical.pdf}
    \caption{(a) Charge $Q$ is uniformly distributed across a shell with a radius of $R$. (b) The dots represent a spherically symmetric charge distribution with a radius of $R$, whose volume charge density $\rho$ is only a function of the distance from the center. The charged object is not a conductor, so the charge is assumed to be fixed in position. A concentric spherical Gaussian surface with $r > R$ is shown. A similar Gaussian surface with $r > R$ is shown in (c).}
    \label{fig:gauss-law-spherical}
\end{figure}
\textsc{Case #1}: \textbf{Spherical Shell Charge Distribution}: 
\begin{tabular}{cl}
    \textbf{Region} & \textbf{Electric Field} \\[3pt]
    $r<R$ & $E=0$ \\[5pt]
    $r=R$ & $\displaystyle E=\frac{Q}{4\pi\epsilon_0 R^2}$ \\[8pt]
    $r>R$ & $\displaystyle E=\frac{Q}{4\pi\epsilon_0 r^2}\sim E\propto\frac{1}{r^2}$
\end{tabular}

\textsc{Case #2}: \textbf{Spherical Charge Distribution}: 
\begin{tabular}{cl}
    \textbf{Region} & \textbf{Electric Field} \\[3pt]
    $r<R$ & $\displaystyle E=\frac{Qr}{4\pi\epsilon_0 R^3}$ \\[6pt]
    $r=R$ & $\displaystyle E=\frac{Q}{4\pi\epsilon_0 R^2}$ \\[8pt]
    $r>R$ & $\displaystyle E=\frac{Qr}{4\pi\epsilon_0 R^3}\sim E\propto\frac{1}{r^2}$
\end{tabular}

We used the uniform behavior of the charge distribution, meaning the volume charge density is constant within the distribution. Our chosen Gaussian surface encloses a fraction of the total charge $\displaystyle Q_\text{enc}=\frac{Qr^3}{R^3}$.

\textsc{Case #3}: \textbf{Excess Charge Distribution on a Solid Spherical Conductor}: 

Inside a conductor, the charges are free to move due to the presence of mobile electrons or ions. When a conductive sphere is charged, the excess charge distributes itself on the sphere's outer surface. This occurs because like charges repel each other, and the charges on the surface of the sphere naturally move as far away from each other as possible.

If there were any excess charges inside the sphere, it would create an electric field within the conductor. The charges inside the sphere would experience a force due to this electric field, causing them to redistribute until the electric field inside the conductor becomes zero. This redistribution of charges continues until the system reaches electrostatic equilibrium.
\subsubsection*{Cylindrical Symmetry}
\begin{figure}[H]
    \centering
    \includegraphics[width=.8\linewidth]{figs/gauss-cylindrical.pdf}
    \caption{A Gaussian surface in the form of a closed cylinder surrounds a section of a very long, uniformly charged, cylindrical plastic rod. Only the boundary wall across the cylinder's length gives a non-zero flux. Both ends of the cylinder contribute nothing to the flux. That is why we ignored the $2\pi r^2$ term from the area.}
    \label{fig:gauss-law-cylindrical}
\end{figure}
An infinitely long charged wire with a charge density $\lambda$ exhibits cylindrical symmetry in its vicinity. We choose a cylinder as our Gaussian surface to preserve symmetry in our calculations. The cylinder has a radius $r$ at both ends and a length of $l$ units, with an area of $2\pi rl + 2\pi r^2$.

Gauss's law states that
\begin{align}
    EA&=\frac{Q_\text{enc}}{\epsilon_0}\notag\\
    \Rightarrow E\times2\pi rl &=\frac{\lambda l}{\epsilon_0}\notag\\
    \Rightarrow E&=\frac{\lambda}{2\pi\epsilon_0 r}\hat{r}.
\end{align}
Here $r$ represents the distance from which the field is being measured. We may arbitrarily choose this in any direction. It predominantly doesn't have to be the $x$ or $y$ axis.
\subsubsection*{Planar Symmetry}
An infinitely long charged sheet with a charge density $\sigma$ shows a head-on cylindrical symmetry around its vicinity. To maintain symmetry in our calculation, we consider a cylinder as our Gaussian surface. The cylinder has a radius $r$ on both ends and is $l$ units long with an area of $2\pi rl + 2\pi r^2$.
\begin{figure}[H]
    \centering
    \includegraphics[width=.72\linewidth]{figs/gauss-planar.pdf}
    \caption{(a) Perspective view and (b) side view of a portion of a very large, thin plastic sheet, uniformly charged on one side to surface charge density $\sigma$. A closed cylindrical Gaussian surface passes through the sheet perpendicular to it. This time, only both ends of the cylinder contribute to the flux. That is why we shall ignore the $2\pi rl$ term from the area.}
    \label{fig:gauss-law-planar}
\end{figure}
Gauss's Law demands
\begin{align}
    EA&=\frac{Q_\text{enc}}{\epsilon_0}\notag\\
    \Rightarrow 2EA &=\frac{\sigma A}{\epsilon_0}\notag\\
    \Rightarrow E&=\frac{\sigma}{2\epsilon_0}.
\end{align}
\subsection{Vector Calculus with $\boldsymbol{\vec{E}}$-Field}
\subsection*{Outflow of $\boldsymbol{\vec{E}}$-Field}
The divergence of $\vec{E}$-field calculates the out(in)flow out of a closed surface.
\begin{align}
    \vec{\nabla}\cdot\vec{E} &= \frac{\rho_\text{enclosed}}{\epsilon_0}.\qquad\text{(Differential Form)}\label{gauss-differential}
\end{align}
\textsc{Note}: We have used the Divergence/Gauss's theorem to the integral form of Gauss's Law (\ref{gauss-integral}) to get the differential form (\ref{gauss-differential}).
\subsection*{Circulation of $\boldsymbol{\vec{E}}$-Field}
In electrostatics, the electric fields have zero curl.
\begin{align}
    \vec{\nabla}\times\vec{E}=0.
\end{align}
This directly implies, from Stokes's theorem, that in electrostatics,
\begin{align}
    \oint \vec{E}\cdot d\vec{l}=0.
\end{align}
The rest is another story. We shall deal with that in the latter half of the course.
\section{Work Done due to Electrostatic Force}
\begin{figure}[H]
    \centering
    \includegraphics[width=.85\linewidth]{figs/elec-potential-energy.pdf}
    \caption{A test charge $q_0$ moves from one point to another along the path shown in a non-uniform electric field. During a displacement, an electric force acts on the test charge. This force points in the direction of the field line at the location of the test charge.}
    \label{fig:electric-potential-energy}
\end{figure}
When a charge is moved against the electric field, work is done by an external force to overcome the repulsive forces between the moving charge and the other charges in the field. This work done is stored as potential energy in the system.
\begin{align}
    W=\bigintss dW&=\bigintss\vec{F}\cdot d\vec{r}\notag\\
    &=q_0\bigintss_{\text{start}}^{\text{final}}\vec{E}\cdot d\vec{r}.
\end{align}
The above formula gives you the work done to move one charge from one point to another within an electric field. There needs not be a target charge at the final point. You only need one charge to displace in the field. 
\begin{figure}
    \centering
    \includegraphics[width=.85\linewidth]{figs/electric-field-work-down.pdf}
    \caption{A test charge $q_0$ moves from one point to another along the path shown in a uniform electric field produced by a point charge $q$.}
    \label{fig:electric-potential-energy}
\end{figure}
\begin{align}
    W&=\bigintss dW\notag\\
    &=\bigintss\vec{F}\cdot d\vec{r}\notag\\
    &=\frac{1}{4\pi\epsilon_0}\bigintss_{\text{start}}^{\text{final}}\frac{q_0q}{r^2}\cdot dr\notag\\
    &=\frac{q_0q}{4\pi\epsilon_0}\bigintss_{\text{start}}^{\text{final}}\frac{dr}{r^2}\notag\\
    &=\frac{q_0q}{4\pi\epsilon_0}\left(\frac{1}{r_\text{initial}}-\frac{1}{r_\text{final}}\right)
\end{align}
The above formula gives you the work done to move one charge from one point to another within the vicinity of a charge, producing an electric field. There needs to be a target charge at the final point. You need two or more charges to displace in the field. 

If the general displacement in which the start and end points do not lie in the same radial line, the work done is given by
\begin{align}
    W=\bigintss_{\text{start}}^{\text{final}}F\cos\theta dl=\bigintss_{\text{start}}^{\text{final}}\frac{q_0q}{r^2}\cos\theta dl,
\end{align}
where $dr=\cos\theta dl$ is the radial component of the line element.
\subsection*{What do signs of work done mean here?}
Imagine you are moving a test charge (positive charge) within a field produced by a positive charge. The test charge will feel a repulsive force during this displacement. If the displacement is toward the target positive charge, the work done comes out negative. In a different case where the same test charge was to be moved toward a negative charge, the force is attractive, making the work done positive.
\subsection{To (Dis)Assemble Many Point Charges in the presence of One Point Charge}
\begin{figure}[H]
    \centering
    \includegraphics[width=.55\linewidth]{figs/work-done-for-multiple-charges.pdf}
    \caption{The potential energy associated with a charge $q_0$ at point a depends on the other charges $q_1$, $q_2$, and $q_3$ and on their distances $r_1$, $r_2$, and $r_3$ from point $q_0$.}
    \label{fig:work-done-several-point}
\end{figure}
Suppose a charge $q_0$ moves in an electric field $\vec{E}$ made or caused by several point charges $q_1,q_2,q_3,\ldots\,\ldots$ at distances $r_1,r_2,r_3,\ldots\,\ldots$ from $q_0$. The superposition principle will find the total work done.
\begin{align}
    W=\frac{q_0}{4\pi\epsilon_0}\left(\frac{q_1}{r_1}+\frac{q_2}{r_2}+\frac{q_3}{r_3}+\ldots\,\ldots\right)=\frac{q_0}{4\pi\epsilon_0}\sum_i\frac{q_i}{r_i}.
\end{align}
\textsc{Note}: The $r_{i}$ in the above equation are distances between the $i^{\text{ th}}$ charge and $q_0$. The source charge $q_0$ must be located from the origin of a chosen coordinate system first.
\subsection{To (Dis)Assemble a Discrete Point Charge Distribution from Scratch}
Imagine instead of the motion of the charge $q_0$; we now try to calculate work done to form the collection of point charges $q_1,q_2,q_3,\ldots\,\ldots$ all separated from each other by infinite distances and then bringing them together. The distance between $q_i$ and $q_j$ is $r_{ij}$. The work done now is given by
\begin{align}
    W=\frac{1}{4\pi\epsilon_0}\sum_{i<j,\,i\ne j}\frac{q_iq_j}{r_{ij}}.
\end{align}
\textsc{Note}: The $r_{ij}$ in the above equation are distances between the $i^{\text{ th}}$ and $j^{\text{ th}}$ charges. Those charges must be located from the origin of a chosen coordinate system. Thus, $\vec{r}_{ij}=\vec{r}_i-\vec{r}_j$.

From the conservation of mechanical energy and the work-energy theorem, we get $\Delta U=-W$.
\section{Electric Potential Energy}
Physically, \textbf{Electric potential energy is the work done to bring a charge from infinity (where the electric field is considered to be zero) to a particular location in an electric field}. {Electrical potential energy} $U$ is required to move within an electric field from one point to another. It is the negative of the work done to execute the same displacement.
\begin{align}
    \Delta U&=U_\text{final}-U_\text{initial}\notag\\
    &=-\frac{q_0q}{4\pi\epsilon_0}\left(\frac{1}{r_\text{initial}}-\frac{1}{r_\text{final}}\right)\notag\\
    &=\frac{q_0q}{4\pi\epsilon_0}\left(\frac{1}{r_\text{final}}-\frac{1}{r_\text{initial}}\right)\\
    &=-q_0\bigintss_{\text{start}}^{\text{final}}\vec{E}\cdot d\vec{r}\notag\\
    \Rightarrow\Delta U&=q_0\bigintss_{\text{final}}^{\text{start}}\vec{E}\cdot d\vec{r}.
\end{align}
\section{Electric Potential}
Electric potential describes \textbf{the amount of electric potential energy per unit charge at a specific point in an electric field}. It represents the \textit{electrical potential energy} that a positive test charge would possess if placed at that point in the field.

\textbf{The electric potential} at a specific point in an electric field \textbf{is defined as the electric potential energy per unit charge}. It is denoted by the symbol $V$ and is measured in volts $V$. 
\begin{figure}[H]
    \centering
    \includegraphics[width=.9\linewidth]{figs/electric-potential-increase-decrease.pdf}
    \caption{If you move in the direction of $\vec{E}$, electric potential V decreases; if you move in the direction opposite \vec{E}, $V$ increases.}
    \label{fig:potential-change-direction}
\end{figure}
The electric potential at a point is a \textit{scalar quantity} and can be positive or negative, depending on the relative positions and magnitudes of the charges involved.
\begin{align}
    \Delta V=\frac{\Delta U}{q_0}=\frac{-W}{q_0}.
\end{align}
\subsection{Electric Potential Difference}
Also known as voltage difference, it measures the difference in electric potential between two points in an electric field. It represents the work done per unit charge moving a charge between those two points.

It is defined as the amount of work required to move a charge from one point to another in an electric field. It is directly related to the electric field, which determines the force experienced by a charge. The electric field pushes or pulls the charge, and the potential difference measures the energy transfer as the charge moves through the field.

Consider a region where there is a uniform electric field. If a positive test charge is moved against the direction of the electric field, from point $A$ to point $B$, work is done against the electric field, and the charge gains potential energy. The electric potential at point $A$ is higher than at point $B$, and the potential difference $V$ between the two points is given by:
\begin{align}
    \Delta V&=V_B - V_A \notag\\
    &= \frac{U_B-U_A}{q_0}\notag\\
    &=\frac{-W_{AB}}{q_0}\\
    &=\frac{\displaystyle-\bigintss_A^B \vec{F}\cdot d\vec{r}}{q_0}\notag\\
    &=\frac{\displaystyle\bigintss_B^A q_0\vec{E}\cdot d\vec{r}}{q_0}\notag\\
    &=\bigintss_B^A\vec{E}\cdot d\vec{r}.
\end{align}
where $\Delta V$ represents the change in electric potential, $V_B$ is the potential at point B (higher), and $V_A$ is the potential at point A (lower).

The relationship between electric potential difference and the electric field (uniform) is given by:
\begin{align}
    \Delta V=Ed,
\end{align}
where $d$ is the separation distance between the two points. This equation shows that the potential difference is directly proportional to the electric field strength and the separation distance.
\subsection{Equipotential Surfaces and Gauss's Law (Again!)}
These are imaginary surfaces in an electric field where all points on the surface have the same electric potential. In other words, they are surfaces that connect points with the same electric potential value.
\begin{figure}[H]
    \centering
    \includegraphics[width=.7\linewidth]{figs/equipotential_of_dipole.pdf}
    \caption{Contour map of the range of equipotential surfaces of an electric dipole. These surfaces get dense because the electric field is strong between the two component charges. This is so because a strong field means the places where you would get the same field reading change drastically, giving closely spaced equipotential surfaces.}
    \label{fig:equipotential_of_dipole}
\end{figure}
A few important properties of equipotential surfaces are:
\begin{itemize}
    \item \textbf{Gaussian surfaces}: All the Gaussian surfaces that we considered during field calculation were chosen to make up equipotential surfaces. That way, we can integrate within the surface without changing the field component anywhere.
    \item \textbf{Perpendicular to Electric Field Lines}: Equipotential surfaces are always perpendicular to the electric field lines at every point. This means the electric field lines intersect the equipotential surfaces at right angles.
    \item \textbf{No Potential Difference}: The electric potential remains constant within an equipotential surface. This implies that no work is required to move a charge along the surface because the potential difference is zero.
    \item \textbf{Uniform Potential}: All points on an equipotential surface have the same electric potential. The potential values may vary between different equipotential surfaces, but the potential remains constant within each surface.
    \item \textbf{Closer Spacing Indicates Stronger Field}: The spacing between equipotential surfaces indicates the strength of the electric field. Closer spacing between the surfaces indicates a stronger electric field, while broader spacing indicates a weaker field.
\end{itemize}
They represent regions in an electric field where the work done in moving a charge between any two points on the surface is zero. Since the electric potential is constant along an equipotential surface, no work is required to move a charge within that surface because the potential difference is zero. This implies that the electric field is perpendicular to the equipotential surface at every point.
 \section{Electric Potential Measurements for Charge Distributions}
\subsection{Discrete I: Electric Potential due to a Charged Particle}
\begin{figure}[H]
    \centering
    \includegraphics[width=.8\linewidth]{figs/electric-potential-charged-particle.pdf}
    \caption{The particle with positive charge $q$ produces an electric field and an electric potential $V$ at point $P$. Moving a test charge $q_0$ from $P$ to infinity, we find the potential. The test charge is shown at a distance $r$ from the particle during differential displacement.}
    \label{fig:potential-point-charge}
\end{figure}
To find the potential of a point charge $Q$, we need to bring the test charge $q_0$ from $P$ toward infinity. We need to remember that $V=0$ at $r\to\infty$.
\begin{align}
    V_f-V_i&=-\bigintss_{{r=R}}^{{r=\infty}} Edr\notag\\
    \Rightarrow 0-V&=-\frac{Q}{4\pi\epsilon_0}\bigintss_{{r=R}}^{{r=\infty}} \frac{dr}{r^2}\notag\\
    &=\frac{Q}{4\pi\epsilon_0}\left[\frac{1}{r}\right]_{r=R}^{r=\infty} \notag\\
    &=-\frac{Q}{4\pi\epsilon_0R}\notag\\
    \therefore V&=\frac{Q}{4\pi\epsilon_0R}.
\end{align}
\subsection{Discrete II: Electric Potential due to a Collection of Charged Particle}
Using the superposition principle, we need to find the net electric potential at a point due to a collection of point charges.
\begin{align}
    V=\sum_i^nV_i=\frac{1}{4\pi\epsilon_0}\sum_i^n\frac{q_i}{r_i}.
\end{align}
\subsection{Continuous I: Electric Potential due to a Continuous Charge Distribution}
When the charge distribution is continuous, we cannot sum the potential contribution but integrate it.

Let us take the $V=0$ at $r\to\infty$. If we treat the element of charge $dq$ as a particle, then we can express the potential $dV$ due to $dq$:
\begin{align}
    dV=\frac{1}{4\pi\epsilon_0}\frac{dq}{r},
\end{align}
$dq$ can be either positive or negative.

The total electric potential would then be
\begin{align}
    V=\bigintss dV=\frac{1}{4\pi\epsilon_0}\bigintss\frac{dq}{r}.
\end{align}
The integral must be taken over the entire charge distribution. Note that because
the electric potential is a scalar, there are no vector components to consider.
\subsection{Continuous II: Electric Potential due to a Charged Conducting Sphere}
We can easily infer this using our ideas on the electric field of a spherical shell charge distribution or an excess charge distributed on a conducting sphere case. We found $\vec{E}=0$ inside the shell(or sphere) in both cases. We found a maximum at the shell(or sphere) radius and points away from the source, $\vec{E}$ dies off as $\displaystyle\frac{1}{r^2}$. We use these values to calculate the work done per unit test charges to move it within this $\vec{E}$ in the three regions.
\begin{center}
\begin{tabular}{cll}
    \textbf{Region} & \textbf{Electric Field} & \textbf{Electric Potential} \\[3pt]
    $r<R$ & $\displaystyle E=0$ & $\displaystyle V=\frac{Q}{4\pi\epsilon_0R}$ \\[6pt]
    $r=R$ & $\displaystyle E=\frac{Q}{4\pi\epsilon_0 R^2}$ & $\displaystyle V=\frac{Q}{4\pi\epsilon_0R}$ \\[8pt]
    $r>R$ & $\displaystyle E=\frac{Q}{4\pi\epsilon_0 r^2}$ & $\displaystyle V=\frac{Q}{4\pi\epsilon_0r}$
\end{tabular}
\end{center}
With no field inside the shell(or sphere), no work is done on a test charge that moves from any point to any other point inside the sphere, giving us a constant $V$. $V$ becomes the maximum where $E$ is the maximum and later dies off as you move farther outward.
\subsection{Continuous III: Electric Potential due to a Finitely Long Line of Charge}

Consider a thin nonconducting rod of length $2L$ with a positive charge of uniform linear density $\lambda$.
\begin{align}
    V=\bigintss dV&=\frac{1}{4\pi\epsilon_0}\bigintss\frac{dq}{r}\notag\\
    &=\frac{1}{4\pi\epsilon_0}\bigintss_{y=-L}^{y=L}\frac{\lambda dy}{\sqrt{x^2+y^2}}\label{eq:line-charge-potential-pause}
\end{align}
In Eq. (\ref{eq:line-charge-potential-pause}), we substitute $y=x\tan{u}$ with $\displaystyle u=\tan^{-1}\left(\frac{y}{x}\right)$.
\begin{align*}
    y&=x\tan{u}\\
    dy&=x\sec^2{u}du.\\
    x^2+y^2&=x^2+\left(x\tan{u}\right)^2\\
    &=x^2\left(1+\tan^2{u}\right)\notag\\
    &=x^2\sec^2{u}.
\end{align*}
\begin{figure}[H]
    \centering
    \includegraphics[width=.7\linewidth]{figs/line-segment-charge-elec-potential.pdf}
    \caption{A thin, uniformly charged rod of length $2L$, produces an electric potential $V$ at point $P$. An element is treated as a particle with a small charge $dQ=\lambda dy$.}
    \label{fig:line-segment-charge-elec-potential}
\end{figure}
With these substitutions, Eq. (\ref{eq:line-charge-potential-pause}) now becomes:
\begin{align}
    V&=\frac{\lambda}{4\pi\epsilon_0}\bigintss\sec{u}du\notag
\end{align}
Here, we must apply the limits of integration for $u$ and not $y$ because we swapped the integrating variable $y$.
\begin{table}[H]
    \centering
    \begin{tabular}{P{1.0cm}|P{3.0cm}}
    $y$ & $\displaystyle u=\tan^{-1}\left(\frac{y}{x}\right)$ \\[8pt]\hline\hline\\
    $+L$ & $\displaystyle\tan^{-1}\left(\frac{+L}{x}\right)$\\[12pt]
    $-L$ & $\displaystyle\tan^{-1}\left(\frac{-L}{x}\right)$\\
\end{tabular}
\end{table}
With these substitutions, Eq. (\ref{eq:line-charge-potential-pause}) now becomes:
\begin{align}
    V&=\frac{\lambda}{4\pi\epsilon_0}\bigint_{\displaystyle\tan^{-1}\left(\frac{-L}{x}\right)}^{\displaystyle\tan^{-1}\left(\frac{+L}{x}\right)}\sec{u}du\notag\\
    &=\frac{\lambda}{4\pi\epsilon_0}\left(\Bigg[\ln{\left|\sec{u}+\tan{u}\right|}\Bigg]_{\displaystyle\tan^{-1}\left(\frac{-L}{x}\right)}^{\displaystyle\tan^{-1}\left(\frac{+L}{x}\right)}\right)\notag\\
    &=\frac{\lambda}{4\pi\epsilon_0}\left(\ln\left|\sec{\left[\tan^{-1}\left(\frac{L}{x}\right)\right]}+\tan{\left[\tan^{-1}\left(\frac{L}{x}\right)\right]}\right|-\ln\left|\sec{\left[\tan^{-1}\left(\frac{-L}{x}\right)\right]}+\tan{\left[\tan^{-1}\left(\frac{-L}{x}\right)\right]}\right|\right)\notag\\
    &=\frac{\lambda}{4\pi\epsilon_0}\left(\ln\left|\sec{\left[\sec^{-1}\left(\frac{\sqrt{x^2+L^2}}{x}\right)\right]}+\tan{\left[\tan^{-1}\left(\frac{L}{x}\right)\right]}\right|\right.\notag\\
    &\left.\hspace{120pt}-\ln\left|\sec{\left[\sec^{-1}\left(\frac{\sqrt{x^2+L^2}}{x}\right)\right]}+\tan{\left[\tan^{-1}\left(\frac{-L}{x}\right)\right]}\right|\right)\notag\\
    &=\frac{\lambda}{4\pi\epsilon_0}\left(\ln\left|\frac{\sqrt{x^2+L^2}}{x}+\frac{L}{x}\right|-\ln\left|\frac{\sqrt{x^2+L^2}}{x}+\left(\frac{-L}{x}\right)\right|\right)\notag\\
    &=\frac{\lambda}{4\pi\epsilon_0}\left(\ln\left|\frac{\sqrt{x^2+L^2}+L}{x}\right|-\ln\left|\frac{\sqrt{x^2+L^2}-L}{x}\right|\right)\notag\\
    &=\frac{\lambda}{4\pi\epsilon_0}\left(\ln{\left|\sqrt{x^2+L^2}+L\right|}-\ln{x}-\ln{\left|\sqrt{x^2+L^2}-L\right|}+\ln{x}\right)\notag\\
    &=\frac{\lambda}{4\pi\epsilon_0}\left(\ln\left|{\frac{\sqrt{x^2+L^2}+L}{\sqrt{x^2+L^2}-L}}\right|\right).
\end{align}
\subsection{Continuous IV: Electric Potential due to an Infinitely Long Line of Charge}
We may take this from an electric field point of view. But unlike the finite charged line segment, defining a potential bringing a charge from infinity and placing it near the source is impractical. Instead, we bring in the charge from a finite distance where $V=0$ and move to another finite distance.
\begin{figure}[H]
    \centering
    \includegraphics[width=.8\linewidth]{figs/infinite-line-charge-elec-potential.pdf}
    \caption{A thin, uniformly charged rod of infinite length produces an electric potential $V$ at point $P$. The charged line has a uniform charge density $\lambda$.}
    \label{fig:infinite-line-charge-elec-potential}
\end{figure}
Consider a positively charged infinitely long line of charge with positive line charge density $\lambda$. The line has field lines emanating radially outward from it. The potential of any point $a$ with respect to $b$, radial distances $r_a$, and $r_b$ from the line charge is
\begin{align}
    V_a-V_b&=\bigintss_a^b E dr\notag\\
    &=\frac{\lambda}{2\pi\epsilon_0}\bigintss_a^b\frac{dr}{r}\notag\\
    &=\frac{\lambda}{2\pi\epsilon_0}\ln\left|{\frac{r_b}{r_a}}\right|.
\end{align}
Consider a finite starting distance where $V=0$. Then, we would get the potential for the point $a$ with a radial distance $r_a=r$ from the source. If we take the starting point $r_b=r_0$.
\begin{align}
    V_b-0=V&=\frac{\lambda}{2\pi\epsilon_0}\ln{\frac{r_0}{r}}.
\end{align}
\subsection{Continuous V: Electric Potential due to an Infinitely Long Charged Cylinder}
This is the extended case of potential due to a charged conducting sphere, except the electric field resembles that of an infinitely long charged line. Similar to the finite charged line segment, defining a potential bringing a charge from infinity and placing it near the source is impractical. Instead, we bring in the charge from a finite distance where $V=0$ and move to another finite distance.
\begin{figure}[H]
    \centering
    \includegraphics[width=.82\linewidth]{figs/infinite-cylindrical-elec-potential.pdf}
    \caption{A cylinder of radius $R$, positively charged on its outer surfaces to a uniform surface charge density $\sigma$. Find the potential $V$ at point $P$ on the central axis of the cylinder.}
    \label{fig:potential-cylinder}
\end{figure}
\begin{center}
\begin{tabular}{cll}
    \textbf{Region} & \textbf{Electric Field} & \textbf{Electric Potential} \\[3pt]
    $r<R$ & $\displaystyle E=0$ & $\displaystyle V=\frac{\lambda}{2\pi\epsilon_0}\ln{\frac{R}{r}}$ \\[6pt]
    $r=R$ & $\displaystyle E=\frac{\lambda}{2\pi\epsilon_0 R}$ & $\displaystyle V=0$ \\[8pt]
    $r>R$ & $\displaystyle E=\frac{\lambda}{2\pi\epsilon_0 r}$ & $\displaystyle V=\ln{\left(\frac{r_a}{r_b}\right)\quad r_a>r>r_b$
\end{tabular}
\end{center}
\subsection{Continuous VI: Electric Potential due to a Charged Ring}
\begin{figure}
    \centering
    \includegraphics[width=.7\linewidth]{figs/ring-charge-elec-potential.pdf}
    \caption{A ring of radius $R$, charged on its circumference to a uniform line charge density $\lambda$. Find the potential $V$ at point $P$ on the central axis of the disk.}
    \label{fig:ring-charge-elec-potential}
\end{figure}
Consider a ring of radius $R$ with a positive charge $Q$ of uniform line density $\lambda$ (Fig. \ref{fig:ring-charge-elec-potential}).
\begin{align}
    V=\bigintss dV=&=\frac{1}{4\pi\epsilon_0}\bigintss\frac{dq}{r}\notag\\
    &=\frac{1}{4\pi\epsilon_0}\frac{1}{\sqrt{x^2+R^2}}\bigintss dq\notag\\
    &=\frac{1}{4\pi\epsilon_0}\frac{Q}{\sqrt{x^2+R^2}}.
\end{align}
\subsection{Continuous VII: Electric Potential due to a Charged Flat Disk}
\begin{figure}[H]
    \centering
    \includegraphics[width=.7\linewidth]{figs/solid-disk-charge-elec-potential.pdf}
    \caption{A flat disk of radius $R$, charged on its top surface to a uniform surface charge density $\sigma$. Find the potential $V$ at point $P$ on the central axis of the disk.}
    \label{fig:potential-disk-charge}
\end{figure}
Consider a flat disk of radius $R$ with a positive charge of uniform surface density $\sigma$. In this case $dq=\sigma(2\pi r)(dr)$.
\begin{align}
    V=\bigintss dV&=\frac{1}{4\pi\epsilon_0}\bigintss\frac{dq}{r}\notag\\
    &=\frac{1}{4\pi\epsilon_0}\bigintss_{r=0}^{r=R}\frac{2\pi\sigma rdr}{\sqrt{x^2+r^2}}\notag\\
    &=\frac{\sigma}{2\epsilon_0}\bigintss_{r=0}^{r=R}\frac{rdr}{\sqrt{x^2+r^2}}\label{eq:sheet-charge-potential-pause}
\end{align}
In Eq. (\ref{eq:sheet-charge-potential-pause}), we substitute $\displaystyle u=x^2+r^2$ such that $du=2rdr$.

With these substitutions, Eq. (\ref{eq:sheet-charge-potential-pause}) now becomes:
\begin{align}
    V&=\frac{\sigma}{2\epsilon_0}\bigintss\frac{du}{2\sqrt{u}}\notag
\end{align}
Here, we must apply the limits of integration for $u$ and not $r$ because we swapped the integrating variable $r$.
\begin{table}[H]
    \centering
    \begin{tabular}{P{1.0cm}|P{3.0cm}}
    $y$ & $\displaystyle u=x^2+r^2$ \\[8pt]\hline\hline\\
    $0$ & $\displaystyle x^2$\\[12pt]
    $R$ & $x^2+R^2$\\
\end{tabular}
\end{table}
With these substitutions, Eq. (\ref{eq:sheet-charge-potential-pause}) now becomes:
\begin{align}
    V&=\frac{\sigma}{2\epsilon_0}\bigg[\sqrt{u}\bigg]_{u=x^2}^{u=x^2+R^2}\notag\\[6pt]
    &=\frac{\sigma}{2\epsilon_0}\left[\sqrt{x^2+R^2}-x\right].
\end{align}
\section{Capacitance and Capacitors}
\begin{figure}[H]
    \centering
    \includegraphics[width=\linewidth]{figs/capacitance-def.pdf}
    \caption{(a) A parallel-plate capacitor comprises two plates of area $A$ separated by a distance d.The charges on the facing plate surfaces have the same magnitude, $q$, but opposite signs. (b) As the field lines show, the electric field due to the charged plates is uniform in the central region between the plates. The field is not uniform at the edges of the plates, as indicated by the \textbf{fringing} of the field lines.}
    \label{fig:capacitance}
\end{figure}
Capacitors are electronic components designed to store and release electrical energy. They consist of two conductive plates separated by a dielectric material. When a voltage difference is applied across the plates, an electric field is established, and the capacitor stores electric charge, which can be later released.

When a capacitor charges, its plates have charges of equal magnitudes but
opposite signs: $+Q$ and $-Q$. 

\textsc{Note}: We call the charge of a capacitor $Q$. Remember, $Q$ is not the net charge on the capacitor, which is zero.

Because the plates are conductors, they are equipotential surfaces; both share the same electric potential. There is a potential difference between the two plates applied using a battery source.

\textsc{Note}: Although we will calculate the potential difference across the two plates, we shall use $V$ and not $\Delta V$. Why? Historical reasons and conventions.

The capacitance of the two capacitors is then defined by
\begin{align}
    C=\frac{Q}{V}.
\end{align}
The capacitance measures how much charge must be put on the plates to produce a specific potential difference between them: \textit{The greater the capacitance, the more charge is stored}.
The SI unit for capacitance is the \textit{farad} ($F$) $\displaystyle1\,\text{farad}=1F=1\,\text{coulomb per volt}=\frac{\text{C}}{\text{V}}$.
\subsection{Calculating Capacitance}
We follow the following two steps:
\begin{enumerate}
    \item \textbf{Calculate the Charge of the capacitor using Gauss's Law}
    \begin{align}
        \epsilon_0\oint\vec{E}\cdot d\vec{A}&=Q_\text{enc}=Q.\notag\\
        Q&=\epsilon_0E\oint dA\notag\\
        \therefore Q&=\epsilon_0EA.
    \end{align}
    \item \textbf{Calculate the Voltage of the capacitor}
    \begin{align}
        V&=V_f-V_i=-\bigintss_i^f\vec{E}\cdot d\vec{r}\notag\\
        \therefore V&=\bigintss_{r=0}^{r=d}Edr=Ed.
    \end{align}
    Thus, the primary formula for capacitance (of a parallel-plate capacitor) is \begin{align}
        C=\frac{\epsilon_0A}{d}.
    \end{align}
    $\epsilon_0$ is the permittivity of free space in this equation. It measures the tendency of resistance to an electric field when a dielectric material is placed in the field. For any dielectric material, $\epsilon$ is the permittivity. For vacuum/air/free space, $\epsilon=\epsilon_0$.

    \textsc{Note}: The Capacitance of a capacitor only depends on the geometry of the conductor plates.
\end{enumerate}
The dielectric material, an insulator, prevents the direct flow of current between the plates. Instead, it allows the build-up of an electric field, causing the capacitor to store electrical energy in the form of accumulated charge.
\subsection{How do Capacitors store Energy?}
Capacitors retain electrical energy due to the separation of charge on the plates and the presence of the dielectric material. An external agent must do work to charge a capacitor. We can imagine doing the work ourselves by transferring electrons from one plate to the other, one by one. As the charges pile up, the electric field between the plates increases, which opposes the continued transfer. So, a greater amount of work is required. A battery does all this for us at the expense of its stored chemical energy. This spent energy is stored in the capacitor plate. 

The stored energy can be released rapidly When the capacitor is connected to a circuit. The charge flows from one plate to another, discharging the capacitor and releasing the stored energy. This discharge can occur in a controlled manner, allowing capacitors to provide bursts of energy when needed.

Suppose $q^{'}$ is transferred to the plates in a short period. The potential difference between the plates at that same time is $\displaystyle V^{'}=\frac{q^{'}}{C}$. If an extra small amount of charge $dq$ is transferred, the work increases as follows: 
\begin{align}
    dW&=V^{'}dq^{'}=\frac{q^{'}}{C}dq^{'}.
\end{align}
The total work done is evaluated by integrating both sides
\begin{align}
    W=\bigintss_{0}^{W} dW=\frac{1}{C}\bigintss_{0}^{Q} q^{'}dq^{'}=\frac{Q^2}{2C}.
\end{align}
This work is stored as electric potential energy in the field.
\begin{align}
    U=\frac{1}{2C}Q^2=\frac{1}{2}CV^2.
\end{align}
In a parallel-plate capacitor, neglecting fringing, the electric field has the same value at all points between the plates. Thus, the energy density $u$—the potential energy per unit volume between the plates—should also be uniform.
\begin{align}
    u=\frac{U}{V}=\frac{CV^2}{2Ad}=\frac{1}{2}\epsilon_0\left(\frac{V}{d}\right)^2=\frac{1}{2}\epsilon_0E^2
\end{align}
\section{Capacitors in Circuit}
\subsection{Series Connection}
In a series connection, the voltage across each capacitor is not the same, and the total voltage across the series combination is equal to the sum of the voltages across each individual capacitor.
\begin{figure}[H]
    \centering
    \includegraphics[width=.83\linewidth]{figs/capacitors-series.pdf}
    \caption{(a) Three capacitors connected in series to a battery. The battery maintains a potential difference $V$ between the top and bottom plates of the series combination. (b) The equivalent capacitor, with capacitance $C_\text{eq}$, replaces the series combination.}
    \label{fig:capacitance-series}
\end{figure}
\textbf{When a potential difference $V$ is applied across several capacitors connected in series, the capacitors have identical charge $Q$. The sum of the potential differences across all the capacitors is equal to the applied potential difference $V$.}

From the circuit,
\begin{align}
    V_1&=\frac{Q}{C_1},\qquad V_2=\frac{Q}{C_2},\qquad V_3=\frac{Q}{C_3}\\
    V&=V_1+V_2+V_3\notag\\
    &=\frac{Q}{C_1}+\frac{Q}{C_2}+\frac{Q}{C_3}\notag\\
    V&=Q\left(\frac{1}{C_1}+\frac{1}{C_2}+\frac{1}{C_3}\right)\notag\\
    \frac{V}{Q}&=\frac{1}{C_1}+\frac{1}{C_2}+\frac{1}{C_3}\notag\\
    \frac{1}{C_\text{eq}}&=\frac{1}{C_1}+\frac{1}{C_2}+\frac{1}{C_3}=\sum_i^n\frac{1}{C_i}.
\end{align}
\subsection{Parallel Connection}
In a parallel connection, the voltage across each capacitor is the same, and the total charge stored in the combination is the sum of the charges stored in each individual capacitor.

From the circuit (Fig. \ref{fig:capacitance-parallel}),
\begin{align}
    Q_1&=C_1V,\qquad Q_2=C_2V,\qquad Q_3=C_3V\\
    Q&=Q_1+Q_2+Q_3\notag\\
    &=C_1V+C_2V+C_3V\notag\\
    Q&=V\left(C_1+C_2+C_3\right)\notag\\
    \frac{Q}{V}&=C_1+C_2+C_3\notag\\
    C_\text{eq}&=C_1+C_2+C_3=\sum_i^n{C_i}.
\end{align}
\textbf{When a potential difference $V$ is applied across several capacitors connected in parallel, that potential difference $V$ is applied across each capacitor. The total charge $q$ stored on the capacitors is the sum of the charges stored on all the capacitors.}
\begin{figure}[H]
    \centering
    \includegraphics[width=.92\linewidth]{figs/capacitors-parallel.pdf}
    \caption{(a) Three capacitors connected in parallel to a battery. The battery maintains potential difference $V$ across its terminals and, thus, across each capacitor. (b) The equivalent capacitor, with capacitance $C_\text{eq}$, replaces the parallel combination.}
    \label{fig:capacitance-parallel}
\end{figure}
\subsection{Capacitance with a Dielectric}
\begin{figure}[H]
    \centering
    \includegraphics[width=.82\linewidth]{figs/induced-field.pdf}
    \caption{(a) A nonpolar dielectric slab is placed between the capacitor plates. The circles represent the electrically neutral atoms within the dielectric. (b) The field slightly stretches the atoms, separating the centers of a positive and negative charge. The separation produces surface charges on the slab faces. These charges set up a field in the opposite direction, generating $\vec{E}_i$. $\vec{E}_\text{net}$ is the new, post-dielectric field that works for the capacitor.}
    \label{fig:induced-capacitance}
\end{figure}
If we fill the space between the plates of a capacitor with a dielectric, the induced $\vec{E}_i$ reduces the original $\vec{E}$ field. The net field $\vec{E}_\text{net}$ reduces the potential difference between the capacitor plates. The capacitance must increase since the battery source must still provide $Q$ to the plates. Assume without the dielectric, the capacitance is $C_0$. The increased capacitance is $C$. They are related to a constant called \textit{dielectric constant}.
\begin{align}
    \frac{C}{C_0}=K=\frac{\displaystyle\frac{\epsilon A}{d}}{\displaystyle\frac{\epsilon_0 A}{d}}=\frac{\epsilon}{\epsilon_0}.\label{eq:permittivity}
\end{align}
We can extend the above equation to other important parameters related to capacitance.
\begin{align}
    K&=\frac{C}{C_0}=\frac{\displaystyle\frac{Q}{V}}{\displaystyle\frac{Q}{V_0}}=\frac{V_0}{V}=\frac{E_0d}{{E}_\text{net}d}=\frac{E_0}{{E}_\text{net}}.
\end{align}
In the limit, the plate separation is very small in practice; we can assume the following for the electric fields at work for both cases: capacitors with and without a dielectric.
\begin{align}
    E_0=\frac{\sigma}{\epsilon_0},&\qquad E_i=\frac{\sigma_i}{\epsilon_0}.\\
    E_\text{net}&=E_0-E_i=\frac{1}{\epsilon_0}(\sigma-\sigma_i).\\
    K&=\frac{E_0}{{E}_\text{net}}=\frac{\sigma}{\sigma-\sigma_0}.
    % \sigma_i&=\sigma\left(1-\frac{1}{K}\right).
\end{align}
Using (\ref{eq:permittivity}), we can generate all the equations for a capacitor with a dielectric inside.
\begin{align}
    C&=\frac{\epsilon A}{d}=\frac{K\epsilon_0 A}{d}.\\
    u&=\frac{1}{2}\epsilon E^2=\frac{1}{2}K\epsilon_0E^2.
\end{align}
\section{Electric Current versus Steady Electric Current}
\begin{figure}[H]
    \centering
    \includegraphics[width=.95\linewidth]{figs/electric-current.pdf}
    \caption{(a) A loop of copper in electrostatic equilibrium. The entire loop is at a single potential, and the electric field is zero at all points inside the copper. (b) Adding a battery imposes an electric potential difference between the ends of the loop that are connected to the terminals of the battery. The battery thus produces an electric field within the loop, from terminal to terminal, and the field causes charges to move around the loop. This movement of charges is a current $i$.}
    \label{fig:electric-current-def}
\end{figure}
Current is defined as the rate of flow of electric charge. In the context of steady current, although the individual charges within the conductor are stationary, there is a continuous movement of charge carriers. This movement occurs at a constant rate, resulting in a steady flow of charges in a particular direction.

Mathematically, electric current is defined as:
\begin{align}
    I=\frac{dq}{dt}.
\end{align}
According to this equation, the current is directly proportional to the amount of charge passing through a point and inversely proportional to the time taken for that charge to pass.

We can find the charge that passes through a cross-section in a time interval:
\begin{align}
    Q=\int dq=\int idt
\end{align}

The ampere (A) is the SI unit for the current. One ampere is defined as the flow of one coulomb of charge per second:
\begin{align}
     1\text{ A}=\frac{1\text{ C}}{1\text{ s}}.
\end{align}
Steady current refers to the flow of electric charges in a conductor that remains constant over time. 
\begin{figure}[H]
    \centering
    \includegraphics[width=.75\linewidth]{figs/drift-velocity-def.pdf}
    \caption{(a) If there is no electric field inside a conductor, a positive charge carrier moves randomly from point $P_A$ to point $P_B$ in a time $dt$. (b) If an electric field $\vec{E}$ is present, the electric force $\vec{F}=+q\vec{E}$ imposes a small drift (greatly exaggerated here). (c) This drift velocity takes the positive charge to point $P_B^{'}$, a distance $v_ddt$ from $P_B$ in the direction of the force.}
    \label{fig:drift-velocity-def}
\end{figure}
The charge carriers are typically electrons in a conductor, such as a wire. Under the influence of an electric field generated by a voltage source, electrons experience a force that propels them to move in a specific direction. As electrons move, they transfer their charge to neighboring electrons, creating a \textit{domino effect} of charge transfer throughout the conductor.

Even though the individual charges (electrons) may not have a net motion, the transfer of charge from one electron to another results in a net flow of charge in the conductor. This net flow constitutes the steady current.
\subsection{A Deeper Look at Current: Drift Velocity and Model of Metallic Conduction}
Drift velocity refers to the \textit{average} velocity of charged particles, typically electrons, in a conductor due to an applied electric field. It measures the charge carriers' net motion in the direction of the electric current.

In a conductor (Figure \ref{fig:drift-velocity-def}), electrons cannot move at high speeds due to collisions with atoms or ions in the conductor's lattice structure. These collisions cause the charged particles to scatter and change direction randomly. An electric field exerts a force on the electrons, causing them to move in the direction opposite to the field (from negative to positive potential). As the electrons move, they experience frequent collisions with the lattice, which randomize their direction. However, an overall \textbf{drift motion} emerges under the influence of the applied electric field.

Suppose there are $n$ moving charged particles per unit volume. Assume that all the particles move with the same drift velocity with magnitude $v_d$. In a time interval $dt$, each particle moves a distance $v_d dt$. The particles that flow out of the right end of the shaded cylinder with length $v_d dt$ during $dt$ are the particles within this cylinder at the beginning of the interval $dt$. This interval gets to have a name of its own, \textit{mean time between collisions}, $\tau$. It measures the average time elapsed between two successive collisions.

The volume of the cylinder is $V=Av_d dt$, and the number of particles within it is $n Av_d dt$. If each particle has a charge $q$, the charge $dQ$ that flows out of the end of the cylinder during time $dt$ is $dQ=q(nAv_ddt)=nqv_dAdt$.

The drift velocity $v_d$ is related to the current $I$ and the number of charge carriers $n$ per unit volume, passing through a cross-section $A$ by the equation:
\begin{align}
    I=\frac{dQ}{dt}=nqv_dA.\label{current-with-drift-speed}
\end{align}
This equation shows that the current is directly proportional to the drift velocity. The higher the drift velocity of charge carriers, the greater the current flow. 

However, drift velocity is relatively low, typically on the order of millimeters per second, due to the frequent collisions experienced by the charge carriers. Using \nth{1} laws of motion from mechanics: 
\begin{align}
    v_d&=\left(\frac{qE}{m}\right)\tau.\label{drift-speed-mean-time}
\end{align}
\textsc{Note}: It is important to note that the drift velocity is not the speed at which electrical signals or energy propagate through a circuit. Signals propagate nearly at the speed of light.
\subsection{A Deeper Look at Current: Current Density and its Vector Property}
\begin{figure}[H]
    \centering
    \includegraphics[width=.73\linewidth]{figs/current-density.pdf}
    \caption{Streamlines representing current density in the flow of charge through a constricted conductor.}
    \label{fig:current-density-def}
\end{figure}
Imagine a localized view and study the flow of charge through a cross-section $A$ of the conductor at a particular point, having the current per unit density $J$. This density is a vector quantity and points in the direction of the drift velocity of the charge carrier producing the current. The current through a small cross-sectional element is
\begin{align}
    dI&=\vec{J}\cdot d\vec{A}\\
    I=\int dI&=\int{J}d{A}={J}\int d{A}=JA\notag\\
    J&=\frac{I}{A}=nqv_d.
\end{align}
\section{Ohm's Law}
From the previous section, we found from Eqs. (\ref{current-with-drift-speed}) and (\ref{drift-speed-mean-time}):
\begin{align}
    I&=nqA\left(\frac{qE}{m}\right)\tau\\
    J&=\frac{I}{A}=nq\left(\frac{qE}{m}\right)\tau\notag\\
    J&=\left(\frac{nq^2\tau}{m}\right)E\label{ohms-law-with-conductivity}
\end{align}
The current density $\vec{J}$ in a conductor depends on the electric field $\vec{E}$ and on the properties
of the material (everything inside the parentheses). In general, this dependence can be quite complex. But for some materials, especially metals, at a given temperature, $\vec{J}$ is nearly directly proportional to $\vec{E}$. The proportionality constant is called Conductivity (measured in Siemens $\sigma=\unit{\per\ohm\per\meter}$).

Taking the reciprocal of Eq. (\ref{ohms-law-with-conductivity}), gives us a different look at the same relation:
\begin{align}
    E&=\left(\frac{m}{nq^2\tau}\right)J\label{ohms-law-with-resistiity}
\end{align}
The terms in the parentheses also form a new proportionality constant called \textit{Resistivity} $\rho$, measured in $\unit{\ohm\meter}$, which is simply the reciprocal of $\sigma$.

This relationship (\ref{ohms-law-with-conductivity} and \ref{ohms-law-with-resistiity}) is known as Ohm's Law. The more famous form can be easily derived from them.
\begin{align}
    {E}&=\rho{J}\\
    \frac{\Delta V}{L}&=\rho\frac{I}{A}\\
    V&=\left(\frac{\rho L}{A}\right)I\label{eq:ohm-law}
\end{align}
The term in parenthesis is called \textit{Resistance}:
\begin{align}
    R=\rho\frac{L}{A},
\end{align}
where $\rho$ is the specific resistance of the material making up the conductor. Resistance is not necessarily constant for each wire or conductor. It can vary depending on temperature, length, cross-sectional area, and the material used.

Resistance refers to the property of a material or component to resist the flow of electric current. It measures how effectively a material hinders the movement of charged particles (usually electrons) when an electric potential difference is applied across it. Resistance is denoted by the $R$ symbol and measured in ohms $\Omega$.

Mathematically, Ohm's Law (Eq. (\ref{eq:ohm-law})) can be expressed as:
\begin{align}
    \Delta V=IR.
\end{align}
Ohm's Law can be understood using an analogy of water flowing through a pipe. Imagine the flow of water as the flow of electric current, the pressure difference in the pipe as the voltage, and the narrowness of the pipe as the resistance. If you increase the pressure difference (voltage) in the pipe, more water will flow through it. Similarly, if you increase the voltage across a conductor, more current will flow through it. On the other hand, if the pipe becomes narrower (increased resistance), the water flow will decrease. Likewise, if the resistance of a conductor increases, the current flow will decrease.

A more realistic form of Ohm's Law is:
\begin{align}
    I=\frac{\Delta V}{R}
\end{align}
Since current is produced in the first place due to potential differences, it should convince You of this form. Remember, $\Delta V$ is the cause, and $I$ is the effect, in practice. Materials that follow Ohm's law are called Ohmic materials.
\subsection{When Ohm's Law no Longer Works}
While Ohm's Law holds true for many materials and circuits under certain conditions, it is not universally valid in all situations. Here are some scenarios where Ohm's Law may break down:
\begin{enumerate}
    \item \textbf{Non-Ohmic Materials}: Ohm's Law specifically applies to ohmic materials, which have a linear relationship between voltage and current. Ohmic materials obey Ohm's Law within a specific range of voltages and currents. However, some materials, like diodes or transistors, exhibit non-linear behavior and do not follow Ohm's Law.
    \item \textbf{Temperature Dependency}: Ohm's Law assumes a constant temperature. However, in some cases, temperature changes can lead to variations in resistance, which can affect the accuracy of Ohm's Law. For example, as the temperature increases, the resistance of some materials may also increase, causing deviations from Ohm's Law.
    \item \textbf{Nonlinear Circuit Elements}: Circuits that include nonlinear elements, such as transistors or diodes, cannot accurately be described by Ohm's Law since these elements exhibit nonlinear voltage-current relationships.
    \item \textbf{High Frequencies}: At high frequencies, especially in circuits involving inductors and capacitors, the reactance (frequency-dependent resistance) may come into play and alter the behavior of the circuit. Ohm's Law may not fully describe the circuit behavior in such cases.
\end{enumerate}
\section{Electric Circuits}
Electric circuits are interconnected paths through which electric current can flow. They consist of electrical components such as resistors, capacitors, inductors, and switches connected by conductive wires. From previous chapters, we know that moving a charge $q$ within an electric field $\vec{E}$ causes it to gain/lose electric potential energy $\Delta U_E$. Let this current pass through a circuit element, and we can receive this electric potential energy in a more usable, practical form, say a light bulb, cell phone, or computer. Electric circuits are the foundation of modern electrical and electronic devices, providing a means to control and utilize electrical energy for various applications.
\subsection{Electromotive Force of a Power Source}
Electromotive force (EMF) is the electrical potential difference across the terminals of a power source, like a battery or a generator, \textbf{when it is not connected to any external circuit}. It represents the ability of the source to drive or push charges through a circuit. EMF is measured in volts $V$. 
\begin{figure}[H]
    \centering
    \includegraphics[width=.76\linewidth]{figs/EMF-def.pdf}
    \caption{(a) An electric field $\vec{E}$ produced inside an isolated conductor causes a current. (b) The current causes charge to build up at the ends. The charge build-up produces an opposing field $\vec{E}_\text{ind}$, thus reducing the current. (c) After a very short time, $\vec{E}_\text{ind}$ has the same magnitude as $\vec{E}$; then the net field is $\vec{E}_\test{net}=\vec{E}-\vec{E}_\text{ind}$=0, and the current stops completely.}
    \label{fig:EMF-def}
\end{figure}
When a source has an EMF, it creates an electric field within the source that pushes charges from one terminal to another, creating a persistent potential difference between the terminals. To understand the necessity of an EMF, we must learn the scenario without it.

Consider a small segment of the current-carrying conductor, and an electric field is applied across it. Due to this field, the charge carriers will gain a net drift motion and produce a flow of current. But inside the conductor segment, soon, charge carriers of opposite signs pile up in the opposite end, creating an induced electric field $\vec{E}_\text{ind}$ that opposes the original field. The current $I$ now owes its strength to the net surviving electric field $\vec{E}_\test{net}=\vec{E}-\vec{E}_\text{ind}$. Sooner or later, this induced field matches the strength of the original field, and the current stops flowing, beating the whole purpose of a continuous current flow by means of an electric circuit. This is where EMF comes in.

The presence of EMF establishes an electric potential difference that can make charges move in a circuit against the battery's polarity. A persistent EMF hinders the build-up of an induced field and keeps the circuit \textit{on}.

Define the EMF $\mathcal{E}$ as the voltage difference between the two terminals of an ideal battery source:
\begin{align}
    \mathcal{E}=\frac{dW}{dq}=V_{ab}.
\end{align}
\subsection{A Tale of Real Batteries}
In practice, however, batteries have internal resistance $r$ due to the inherent characteristics of their chemical composition and construction. The physical intuition behind the internal resistance can be understood in terms of the obstacles the charge carriers face as they travel within the battery. When an external load (e.g., a light bulb or a motor) is connected to the battery, a current is drawn from the battery to power the load. However, as the current flows through the battery, it encounters resistance within the battery itself. This internal resistance acts like a \textbf{bottleneck} that hinders the smooth flow of current.

The internal resistance effectively reduces the voltage available at the battery terminals compared to its open-circuit voltage $\mathcal{E}$ (voltage without any load connected). The net available potential difference a battery can then supply is:
\begin{align}
    V_{ab}=\mathcal{E}-Ir.\label{eq:available-voltage}
\end{align}
Comparing this to Ohm's Law, the net current provided by the battery when an external load $R$ is connected would become:
\begin{align}
    V_{ab}=IR&=\mathcal{E}-Ir\notag\\
    I&=\frac{\mathcal{E}}{r+R}.\\
    r&=\frac{\mathcal{E}-IR}{I}
\end{align}
%--------------------------------------------------------
%--------------------------------------------------------
%--------------------------------------------------------
\section{Resistors in Electric Circuits}
\subsection{Series Connection}
In a series connection, the voltage across each resistor is not the same, and the total voltage across the series combination is equal to the sum of the voltages across each individual resistor.
\begin{figure}
    \centering
    \includegraphics[width=.7\linewidth]{figs/resistors-series.pdf}
    \caption{(a) Three resistors connected in series to a battery. The battery maintains a potential difference $V$ given by Eq. (\ref{eq:available-voltage}) between the top and bottom plates of the series combination. (b) The equivalent resistor, with resistance $R_\text{eq}$, replaces the series combination.}
    \label{fig:resistance-series}
\end{figure}
\textbf{When a potential difference $V$ is applied across several capacitors connected in series, the resistors allow the identical current flow $I$. The sum of the potential difference drops across all the resistors is equal to the applied potential difference $V$.}

From the circuit in Figure (\ref{fig:resistance-series}),
\begin{align}
    V_1&=IR_1,\qquad V_2=IR_2,\qquad V_3=IR_3\\
    V&=V_1+V_2+V_3\notag\\
    &=IR_1+IR_2+IR_3\notag\\
    V&=I(R_1+R_2+R_3)\notag\\
    \frac{V}{I}&=R_1+R_2+R_3\notag\\
    R_\text{eq}&=R_1+R_2+R_3=\sum_i^n{R_i}.
\end{align}
\subsection{Parallel Connection}
In a parallel connection, the voltage across each resistor is the same, and the current, however, splits up in the junction to a different value for each individual resistor.
\begin{figure}[H]
    \centering
    \includegraphics[width=.8\linewidth]{figs/resistors-parallel.pdf}
    \caption{(a) Three resistors connected parallel to a battery. The battery maintains a potential difference $V$ given by Eq. (\ref{eq:available-voltage}) between the top and bottom plates of the parallel combination. (b) The equivalent resistor, with resistance $R_\text{eq}$, replaces the parallel combination.}
    \label{fig:resistance-parallel}
\end{figure}
From the circuit in Figure (\ref{fig:resistance-parallel}),
\begin{align}
    I_1&=\frac{V}{R_1},\qquad I_2=\frac{V}{R_2},\qquad I_3=\frac{V}{R_3}\\
    I&=I_1+I_2+I_3\notag\\
    &=\frac{V}{R_1}+\frac{V}{R_2}+\frac{V}{R_3}\notag\\
    I&=V\left(\frac{1}{R_1}+\frac{1}{R_2}+\frac{1}{R_3}\right)\notag\\
    \frac{I}{V}&=\frac{1}{R_1}+\frac{1}{R_2}+\frac{1}{R_3}\notag\\
    \frac{1}{R_\text{eq}}&=\frac{1}{R_1}+\frac{1}{R_2}+\frac{1}{R_3}=\sum_i^n\frac{1}{R_i}.
\end{align}
\textbf{When a potential difference $V$ is applied across several resistors connected in parallel, the total current $I$ breaks into different parts for each resistor, and the sum of all the currents recombines at the junction at the other end of the connection.}
\section{Power in Electric Circuits}
In an electric circuit, power is the rate at which electrical energy $U$ is transferred or consumed. It represents how quickly electrical energy is converted into other forms of energy (such as heat, light, or mechanical work) within the circuit.
\begin{align}
    P&=\frac{dW_{ab}}{dt}=\frac{-dU_{ba}}{dt}=\frac{-dqV_{ba}}{dt}=V_{ab}\frac{dq}{dt}=V_{ab}I.\label{power-circuit-def}
\end{align}
If the potential at $a$ is lower than at $b$, then $V_{ab}$ is negative, and there is a net transfer of energy (and power) out of the circuit element. The element then acts as a source, delivering electrical energy (and power) into the circuit to which it is attached. This is the usual situation for a battery, which converts chemical energy into electrical energy and delivers it to the external circuit. Thus $qV_{ab}$ can denote a quantity of energy that is either delivered to a circuit element or extracted from that element, say $R$.
\begin{align}
    P=I^2R=\frac{V_{ab}^2}{R}.
\end{align}
The unit of electrical power is Watt (W). 

\begin{align*}
    1\text{W}=\frac{1\text{J}/1\text{C}}{1\text{C}/1\text{s}}=\frac{1\text{J}}{1\text{s}}.
\end{align*}
\subsection*{What Happens to this Power?}
The moving charges collide with atoms in the resistor and transfer some of their energy to these atoms, increasing the internal energy of the
material. Either the temperature of the resistor increases, or there is a flow of heat out of it, or both. Every resistor has a power rating, the maximum power the device can dissipate without becoming overheated and damaged.
\subsection{Power Output of a Circuit}
The power output of a battery source with an electromotive force (EMF) and internal resistance is the electrical power the battery delivers to an external circuit. The power output is the rate at which the battery transfers energy to the load connected to it.

Consider a battery with an EMF $\mathcal{E}$ (voltage across its terminals when no current is flowing) and an internal resistance $r$. The battery is connected to an external load with resistance $R$, forming a complete circuit.

When a current $I$ flows through the circuit, there is a voltage drop across the internal resistance $r$, and the voltage across the external load $V_{ab}$ (terminal voltage) is given by Eq. (\ref{eq:available-voltage}). Merging that with Eq. (\ref{power-circuit-def}) gives us:
\begin{align}
    P_\text{out}&=V_{ab}I=\mathcal{E}I-I^2r=P_B-P_R.\\
    P_B&=\mathcal{E}I.\\
    P_R&=I_2R.
\end{align}
$P_B$ is the rate at which work is done on the circulating charges by whatever agency causes the nonelectrostatic force in the source, i.e., the battery. This term represents the rate of conversion of
nonelectrical energy to electrical energy within the source.

$P_R$ is the rate at which electrical energy is dissipated in the internal resistance of the source. 

The difference $P_\text{out}=\mathcal{E}I-I^2r$ is the net electrical power output of the source—that is, the rate at which
the source delivers electrical energy to the remainder of the circuit.
\subsection{Power Input of a Circuit}
To calculate the power input to a circuit, we only need to flip the direction of the current $I$ by applying a battery with higher EMF to the lower-EMF battery present in the circuit.

Because of this reversal of current, instead of Eq. (\ref{eq:available-voltage}), we have for the power receiving source:
\begin{align}
    V_{ab}=\mathcal{E}+Ir.
\end{align}
Merging that with Eq. (\ref{power-circuit-def}) gives us:
\begin{align}
    P_\text{in}&=V_{ab}I=\mathcal{E}I+I^2r=P_B+P_R.\\
    P_B&=\mathcal{E}I.\\
    P_R&=I_2R.
\end{align}
This is what happens when a rechargeable battery (a storage battery) is connected to a charger. The charger supplies electrical energy to the battery; part of it is converted to chemical energy, to be reconverted later, and the remainder is dissipated (wasted) in the battery’s internal resistance, warming the battery and causing a heat flow out of it. If you have a power tool or laptop computer with a rechargeable battery, you may have noticed that it gets warm while it is charging.
\section{Solving DC Circuits: Kirchoff's Laws}
Named after German physicist Gustav Kirchhoff (\textit{Kir-koff}), critical in electrical circuit analysis and governs the current and voltage behavior in a circuit. These laws, known as Kirchhoff's current law (KCL) and Kirchhoff's voltage law (KVL), are essential tools for understanding and analyzing complex electrical circuits.
\subsection{Kirchhoff's Current Law (KCL)}
KCL states that the algebraic sum of currents at any junction (or node) in a circuit is zero. In other words, the total current flowing into a node equals the total current flowing out. This law is based on the principle of charge conservation.
\begin{figure}[H]
    \centering
    \includegraphics[width=.56\linewidth]{figs/kcl-current-rule.pdf}
    \caption{Kirchhoff’s Current law or junction rule states that as much current flows into a junction as flows out of it.}
    \label{fig:kcl-def}
\end{figure}
Imagine a junction in a circuit as a meeting point where different paths for current flow converge, much like a traffic intersection. The sum of all the incoming currents $I$ at that junction must equal the sum of all the outgoing currents $I_1+I_2$. It's like cars entering an intersection from different directions and then dispersing along different roads. The total number of cars entering the intersection must equal the number of cars leaving it.

We will use \textbf{Nodal analysis} to implement KCL in our circuit-solving problems. Nodal analysis, also known as the node-voltage method, is a technique that uses Kirchhoff's Current Law (KCL) to determine the voltage at each node in a circuit.

\textsc{Takeaway}: KCL establishes the Conservation of Electric Charge principle for electric circuits.
\subsection{Kirchhoff's Voltage Law (KVL)}
KVL states that the sum of the voltage drops (or rises) around any closed loop in a circuit equals the sum of the voltage sources (such as batteries) in that loop. This law is based on the principle of energy conservation.
\begin{figure}[H]
    \centering
    \includegraphics[width=\linewidth]{figs/single-loop-circuit-kvl.pdf}
    \caption{Kirchhoff’s Voltage law or loop rule states that the algebraic sum of the changes in potential encountered in a complete traversal of any loop of a circuit must be zero.}
    \label{fig:kcl-def}
\end{figure}
\begin{align*}
    \mathcal{E}-iR-IR=\sum V=0
\end{align*}
Think of a closed loop in a circuit as a path you can travel, much like a loop in a roller coaster ride. You may encounter uphill climbs and downhill descents as you move along the loop. The sum of all the changes in height (potential) you experience during the loop must be zero. The energy you gain during uphill climbs must equal the energy you lose during downhill descents.

We will use \textbf{Loop analysis} to implement KVL. Loop analysis, also known as the mesh current method, uses Kirchhoff's Voltage Law (KVL) to determine the currents flowing in each circuit loop.

\textsc{Takeaway}: KCL establishes the Conservation of Energy principle for electric circuits and the fact that \textbf{Electric Force is conservative}.
\subsection*{Sign Convention for the KVL}
\textbf{Resistance Rule}: For a move through resistance in the direction of the current, the change in potential is $-iR$; in the opposite direction, it is $+iR$.

\textbf{EMF Rule}: For a move through an ideal EMF device in the direction of the EMF arrow ($-$ to $+$), the change in potential is $+\mathcal{E}$; in the opposite direction, it is $-\mathcal{E}$.

\section{Transient Series $RC$ Circuits}
A transient current, also known as a \textit{transient response} or transient behavior, refers to \textbf{the temporary, short-lived flow of electric current in an electrical circuit that occurs in response to a sudden change or disturbance in the circuit}. 
\begin{figure}
    \centering
    \includegraphics[width=.75\linewidth]{figs/rc-circuit-def.pdf}
    \caption{When switch $S$ is closed on $a$, the capacitor is charged through the resistor. When the switch is afterward closed on $b$, the capacitor discharges through the resistor.}
    \label{fig:rc-circuit-def}
\end{figure}
Transient currents are characterized by their relatively brief duration, and they occur during the transient period when the circuit transitions from one steady-state condition to another. We will look at two different transient response systems in this course. The first one is the $RC$ circuit discussed below. The second one, the $RL$ circuit, will be discussed after we are introduced to \textit{Inductors}.
\subsection{Charging of a Capacitor: Growth of Charge}
Consider the circuit with the battery, resistor, and capacitor connected in series. The voltage across the capacitor at any time $t$ is denoted as $V$, and the charge stored in the capacitor at time $t$ is $q(t)$.

Use the KVL for the circuit. The instantaneous potential difference $V_{ab}$ and $V_{bc}$ are: $v_{ab}=iR$ and $\displaystyle v_{bc}=\frac{q}{C}$.
\begin{align}
    \mathcal{E}-iR-\frac{q}{C}&=0\notag\\
    i&=\frac{\mathcal{E}}{R}-\frac{q}{RC}\\
    \frac{dq}{dt}&=\frac{\mathcal{E}}{R}-\frac{q}{RC}\notag\\
    \frac{dq}{dt}&=-\frac{1}{RC}\left(q-C\mathcal{E}\right)\notag\\
    \frac{dq}{q-C\mathcal{E}}&=-\frac{dt}{RC}.\notag
\end{align}
The integration limits are $q=0$ at $t=0$ to $q=q$ at $t$,
\begin{align}
    \bigintss_{0}^{q}\frac{dq}{q-C\mathcal{E}}&=-\bigintss_{0}^{t}\frac{dt}{RC}\notag\\
    \ln\left(\frac{q-C\mathcal{E}}{-C\mathcal{E}}\right)&=-\frac{t}{RC}\notag\\
    \frac{q-C\mathcal{E}}{-C\mathcal{E}}&=e^{\displaystyle-\frac{t}{RC}}\notag\\
    q(t)&=C\mathcal{E}\left(1-e^{\displaystyle-\frac{t}{RC}}\right).\label{eq:charge-eq-rc-circuit}
\end{align}
This is the charge equation for a capacitor while charging in a series RC circuit. It shows how the charge in the capacitor increases with time as it charges up to the maximum charge allowed by the circuit and the capacitor.
\begin{figure}[H]
    \centering
    \includegraphics[width=.92\linewidth]{figs/rc-circuit-charging-cycle.pdf}
    \caption{Charge $q$ increases exponentially with time $t$. At the same time, Current $i$ decreases exponentially.}
    \label{fig:rc-charging-cycle}
\end{figure}
At $t=0$, Eq. (\ref{eq:charge-eq-rc-circuit}) leads to:
\begin{align*}
    q(t)=C\mathcal{E}\left(1-e^{\displaystyle-\frac{0}{RC}}\right)=0.
\end{align*}
At $t\to\infty$, Eq. (\ref{eq:charge-eq-rc-circuit}) leads to:
\begin{align*}
    q(t)=C\mathcal{E}\left(1-e^{\displaystyle-\frac{\infty}{RC}}\right)=C\mathcal{E}=Q_f.
\end{align*}
\subsection{Charging of a Capacitor: Decay of Current}
Differentiating Eq. (\ref{eq:charge-eq-rc-circuit}) with respect to $t$ once gives the transient current equation during the charging period of the capacitor.
\begin{align}
    i(t)=\frac{dq}{dt}=\frac{\mathcal{E}}{R}e^{\displaystyle-\frac{t}{RC}}.\label{eq:charging-current-eq-rc-circuit}
\end{align}
At $t=0$, Eq. (\ref{eq:charging-current-eq-rc-circuit}) leads to:
\begin{align*}
    i(t)=\frac{\mathcal{E}}{R}e^{\displaystyle-\frac{0}{RC}}=\frac{\mathcal{E}}{R}=I_0.
\end{align*}
At $t\to\infty$, Eq. (\ref{eq:charging-current-eq-rc-circuit}) leads to:
\begin{align*}
    i(t)=\frac{\mathcal{E}}{R}e^{\displaystyle-\frac{\infty}{RC}}=0.
\end{align*}
\subsection{Charging of a Capacitor: Decay of Resistive Voltages}
During the charging phase, the resistive voltage $V_R$ can be found by manipulating Eq. (\ref{eq:charge-eq-rc-circuit}):
\begin{align}
    V_R^\text{charging}(t)=i(t)R&=\mathcal{E}e^{\displaystyle-\frac{t}{RC}}.\label{eq:charge-rest-voltage-eq-rc-circuit}
\end{align}
At $t=0$, Eq. (\ref{eq:charge-rest-voltage-eq-rc-circuit}) leads to:
\begin{align*}
    V_R^\text{charging}(t)=\mathcal{E}\left(1-e^{\displaystyle-\frac{0}{RC}}\right)=\mathcal{E}.
\end{align*}
At $t\to\infty$, Eq. (\ref{eq:charge-rest-voltage-eq-rc-circuit}) leads to:
\begin{align*}
    V_R^\text{charging}(t)=\mathcal{E}\left(1-e^{\displaystyle-\frac{\infty}{RC}}\right)=0.
\end{align*}
\begin{figure}[H]
    \centering
    \includegraphics[width=\linewidth]{figs/rc-circuit-charging-cycle-voltage.pdf}
    \caption{Resistive voltage $V_R$ decreases exponentially with time $t$. At the same time, the Capacitive voltage $V_C$ increases exponentially during charging.}
    \label{fig:rc-circuit-charging-cycle-voltage}
\end{figure}
\subsection{Charging of a Capacitor: Growth of Capacitive Voltages}
During the charging phase, the voltage can be found by manipulating Eq. (\ref{eq:charge-eq-rc-circuit}):
\begin{align}
    V_C^\text{charging}(t)=\frac{q(t)}{C}&=\mathcal{E}\left(1-e^{\displaystyle-\frac{t}{RC}}\right).\label{eq:charge-capa-voltage-eq-rc-circuit}
\end{align}
At $t=0$, Eq. (\ref{eq:charge-capa-voltage-eq-rc-circuit}) leads to:
\begin{align*}
    V_C^\text{charging}(t)=\mathcal{E}\left(1-e^{\displaystyle-\frac{0}{RC}}\right)=0.
\end{align*}
At $t\to\infty$, Eq. (\ref{eq:charge-capa-voltage-eq-rc-circuit}) leads to:
\begin{align*}
    V_C^\text{charging}(t)=\mathcal{E}\left(1-e^{\displaystyle-\frac{\infty}{RC}}\right)=\mathcal{E}.
\end{align*}
\subsection{Discharging of a Capacitor: Decay of Charge}
Consider the same $RC$ circuit above but have the capacitor discharge now from an arbitrary maximum charge $Q_0$ across the load $R$. Only the EMF is absent now in the circuit.
\begin{align}
    iR-\frac{q}{C}&=0\notag\\
    i&=-\frac{q}{RC}\\
    \frac{dq}{dt}&=-\frac{q}{RC}\notag\\
    \frac{dq}{q}&=-\frac{dt}{RC}\notag
\end{align}
The integration limits are $q=Q_0$ at $t=0$ to $q=q$ at $t$,
\begin{align}
    \bigintss_{Q_0}^{q}\frac{dq}{q}&=-\bigintss_{0}^{t}\frac{dt}{RC}\notag\\
    \ln\left(\frac{q}{Q_0}\right)&=-\frac{t}{RC}\notag\\
    \frac{q}{Q_0}&=e^{\displaystyle-\frac{t}{RC}}\notag\\
    q(t)&=Q_0e^{\displaystyle-\frac{t}{RC}}.\label{eq:discharge-eq-rc-circuit}
\end{align}
\begin{figure}
    \centering
    \includegraphics[width=.92\linewidth]{figs/rc-circuit-discharging-cycle.pdf}
    \caption{Charge $q$ decreases exponentially with time $t$. At the same time, the Current $i$ increases exponentially.}
    \label{fig:rc-discharging-cycle}
\end{figure}
At $t=0$, Eq. (\ref{eq:discharge-eq-rc-circuit}) leads to:
\begin{align*}
    q(t)=Q_0e^{\displaystyle-\frac{t}{RC}}=Q_0.
\end{align*}
At $t\to\infty$, Eq. (\ref{eq:discharge-eq-rc-circuit}) leads to:
\begin{align*}
    q(t)=Q_0e^{\displaystyle-\frac{\infty}{RC}}=0.
\end{align*}
This is the discharge equation for a capacitor in a series RC circuit. It shows how the charge in the capacitor decreases with time as it discharges to the minimum charge allowed by the circuit and the capacitor.
\subsection{Discharging of a Capacitor: Growth of Current}
Differentiating Eq. (\ref{eq:discharge-eq-rc-circuit}) with respect to $t$ once gives the transient current equation during the discharging period of the capacitor.
\begin{align}
    i=\frac{dq}{dt}=-\frac{Q_0}{RC}e^{\displaystyle-\frac{t}{RC}}.\label{eq:discharging-current-eq-rc-circuit}
\end{align}
At $t=0$, Eq. (\ref{eq:discharging-current-eq-rc-circuit}) leads to:
\begin{align*}
    i=-\frac{Q_0}{RC}e^{\displaystyle-\frac{0}{RC}}=-\frac{Q_0}{RC}=I_0.
\end{align*}
At $t\to\infty$, Eq. (\ref{eq:discharging-current-eq-rc-circuit}) leads to:
\begin{align*}
    i=-\frac{Q_0}{RC}e^{\displaystyle-\frac{\infty}{RC}}=0.
\end{align*}
\subsection{Discharging of a Capacitor: Growth of Resistive Voltages}
During the discharging phase, the voltage can be found by manipulating Eq. (\ref{eq:discharge-eq-rc-circuit}):
\begin{align}
    V_R^\text{discharging}(t)=\frac{q(t)}{C}&=\frac{Q_0}{C}e^{\displaystyle-\frac{t}{RC}}.\label{eq:discharge-rest-voltage-eq-rc-circuit}
\end{align}
At $t=0$, Eq. (\ref{eq:discharge-rest-voltage-eq-rc-circuit}) leads to:
\begin{align*}
    V_R^\text{discharging}(t)=-\frac{Q_0}{C}e^{\displaystyle-\frac{0}{RC}}=-\frac{Q_0}{C}.
\end{align*}
At $t\to\infty$, Eq. (\ref{eq:discharge-rest-voltage-eq-rc-circuit}) leads to:
\begin{align*}
    V_R^\text{discharging}(t)=-\frac{Q_0}{C}e^{\displaystyle-\frac{\infty}{RC}}=0.
\end{align*}
\begin{figure}
    \centering
    \includegraphics[width=.8\linewidth]{figs/rc-circuit-discharging-cycle-voltage.pdf}
    \caption{Resistive voltage $V_R$ increases exponentially with time $t$ to zero. At the same time, the Capacitive voltage $V_C$ decreases exponentially during discharging.}
    \label{fig:rc-circuit-discharging-cycle-voltage}
\end{figure}
\subsection{Discharging of a Capacitor: Decay of Capacitive Voltages}
During the discharging phase, the voltage can be found by manipulating Eq. (\ref{eq:discharge-eq-rc-circuit}):
\begin{align}
    V_C^\text{discharging}(t)=\frac{q(t)}{C}&=\frac{Q_0}{C}e^{\displaystyle-\frac{t}{RC}}.\label{eq:discharge-capa-voltage-eq-rc-circuit}
\end{align}
At $t=0$, Eq. (\ref{eq:discharge-capa-voltage-eq-rc-circuit}) leads to:
\begin{align*}
    V_C^\text{discharging}(t)=\frac{Q_0}{C}e^{\displaystyle-\frac{0}{RC}}=\frac{Q_0}{C}=V_0=\mathcal{E}.
\end{align*}
At $t\to\infty$, Eq. (\ref{eq:discharge-capa-voltage-eq-rc-circuit}) leads to:
\begin{align*}
    V_C^\text{discharging}(t)=\frac{Q_0}{C}e^{\displaystyle-\frac{\infty}{RC}}=0.
\end{align*}
In practice, both these cycles occur one after the other. We may attach each characteristic curve with others in sequence to get the whole picture of charging and discharging.
\begin{figure}[H]
    \centering
    \includegraphics[width=.95\linewidth]{figs/capacitor-charge-discharge-cycle.pdf}
    \caption{Voltage across the capacitor plates rises and decays following the Eqs. (\ref{eq:charge-capa-voltage-eq-rc-circuit}) and (\ref{eq:discharge-capa-voltage-eq-rc-circuit}) with time.}
    \label{fig:rc-charging-discharging-cycle}
\end{figure}
\clearpage
%%%%%%%%%%%%%%%%%%%%%%%%%%%%%%%%%%%%%%%%%%%%%%%%%%%%%%%%%%%%%%%%%%%%%
\bigskip \
\vspace{10cm}
\vskip 20pt
\begin{center}
    \thispagestyle{empty}
    \huge{\selectfont \bf \sffamily Part II: $\boldsymbol{\vec{B}}$-fields and all that}
\end{center}
\addcontentsline{toc}{section}{\Large Part II: $\boldsymbol{\vec{B}}$-fields and all that}
\vfill
\clearpage
%%%%%%%%%%%%%%%%%%%%%%%%%%%%%%%%%%%%%%%%%%%%%%%%%%%%%%%%%%%%%%%%%%%%%
\section{Magnetic Field}
\thispagestyle{empty}
\textbf{The magnetic field}, we shall heavily opt for calling it $\vec{B}$ from now on, is a physical vector field. It is produced by several sources, including \textbf{moving charges} (such as current-carrying wires) and \textbf{permanent magnets}. When charges move, they generate a magnetic field around them. Additionally, magnetic materials, such as iron or nickel, have their own intrinsic magnetic fields due to the alignment of their atomic or molecular magnetic moments.
\begin{figure}[H]
    \centering
    \includegraphics[width=.8\linewidth]{figs/geomagnet.pdf}
    \caption{Due to the convectional current produced by molten lava inside Earth's core generates a magnetic field ($5\times10^{-5}\,\unit{\tesla}$) around the Earth. This field shields us from harsh solar radiation, keeping our atmosphere sustainable for life. The geomagnetic north pole almost coincides with the geographic south pole and vice versa for the geomagnetic south pole.}
    \label{fig:geomagnet}
\end{figure}
Mathematically, the magnetic field is defined in terms of \textbf{the magnetic force, called Lorentz force, experienced by a moving charge in the presence of a magnetic field}. The magnetic field is represented by the symbol $B$ and is measured in units of Tesla $T$.
\subsection{Permanent Magnet versus Magnetic Substance versus Electromagnet}
\textbf{Permanent magnets} are objects that possess their own magnetic fields. They can attract or repel other magnets or magnetic materials. A magnet has two poles, namely the North Pole and the South Pole, and the magnetic field lines form closed loops from the North Pole to the South Pole.

\textbf{Magnetic substances} are materials that can be magnetized in the presence of an external magnetic field, but \textbf{do not retain their magnetism} once the external field is removed. Examples include iron, nickel, and cobalt.

An \textbf{electromagnet} is a type of magnet that is created by passing an electric current through a coil of wire. The current flowing through the wire produces a magnetic field, and the strength of the magnetic field can be controlled by adjusting the current. Unlike permanent magnets, the magnetic properties of an electromagnet can be turned on and off by controlling the current.
\subsection{$\boldsymbol{\vec{B}}$-Field Lines}
Magnetic field lines are a visual representation used to depict the direction and strength of the magnetic field at different points in space. They are used to provide an understanding of the magnetic field's behavior and its interaction with other magnetic objects or current-carrying wires.

Magnetic field lines are different from electric field lines in a few ways. While electric field lines start from positive charges and end on negative charges (or go to infinity for isolated charges), \textbf{magnetic field lines form closed loops that extend from the North Pole of a magnet to its South Pole}. This is because there are no isolated magnetic charges (i.e., magnetic monopoles) known to exist in nature.
\begin{figure}[H]
    \centering
    \includegraphics[width=\linewidth]{figs/mag-field-iron-fillings.pdf}
    \caption{Magnetic field lines around a bar magnet. (a) Iron fillings orient themselves in the directions of the magnetic field. (b) Each iron filling can be thought of as tiny bar magnets with their own North and south pole. They orient themselves accordingly to $\vec{B}$. (c) The same field lines are drawn in terms of conventional lines and arrows.}
    \label{fig:mag-field-filings}
\end{figure}
The direction of the magnetic field at any point is tangential to the field line at that point. The density of the magnetic field lines is indicative of the strength of the magnetic field, with a higher field line density representing a stronger field.
\section{Gauss's Law for Magnetism}
Analogous to the case for $\vec{E}$-field, we may try to associate the $\vec{B}$-field origin with a magnetic charge. Gauss's Law for Magnetism states that \textbf{the net magnetic flux through any closed surface is always zero}. Why might that be?

Magnetic field lines always form closed loops, and magnetic field lines always have both a north and south pole. \textbf{There are no magnetic monopoles in nature, and magnetic dipoles produce all magnetic fields} (combinations of north and south poles) \textbf{or magnetic materials}. 
\subsection{Magnetic Flux}
Magnetic flux is a measure of the total magnetic field passing through a given surface. It quantifies the amount of magnetic field lines penetrating an area. The magnetic flux $\Phi_B$ through a closed surface is defined as the surface integral of the magnetic field vector $\vec{B}$ over that closed surface. In mathematical notation, it can be expressed as:
\begin{align}
    \Phi_B=\bigintss\vec{B}\cdot d\vec{A}.\label{gauss-integral-mag}
\end{align}
\begin{figure}[H]
    \centering
    \includegraphics[width=0.6\linewidth]{figs/mag-flux-def.pdf}
    \caption{The magnetic flux through an area element $dA$ is defined to be $d\Phi_B = B_\perp dA$.}
    \label{fig:mag-field-flux}
\end{figure}
Consider that a magnetic monopole would produce magnetic field lines that radiate outward from the monopole or converge into it, similar to how electric field lines behave around electric charges. In such a case, the net magnetic flux through any closed surface surrounding the magnetic monopole would be non-zero because some magnetic field lines would enter or leave the closed surface.
\begin{align}
    \oint\Phi_B=\oint \vec{B}\cdot d\vec{A}=0.
\end{align}
\subsection{Vector Calculus with $\boldsymbol{\vec{B}}$-Field}
\subsection*{Outflow of $\boldsymbol{\vec{B}}$-Field}
The divergence of $\vec{B}$-field calculates the out(in)flow out of a closed surface.
\begin{align}
    \vec{\nabla}\cdot\vec{B} &= 0.\qquad\text{(Differential Form)}\label{gauss-differential-mag}
\end{align}
\textsc{Note}: We have used the Divergence/Gauss's theorem to the integral form of Gauss's law (\ref{gauss-integral-mag}) to get the differential form (\ref{gauss-differential-mag}).
\subsection*{Circulation of $\boldsymbol{\vec{B}}$-Field}
In magnetostatics, where $\Phi_B$ is not varying in time, the magnetic fields have zero curl.
\begin{align}
    \vec{\nabla}\times\vec{B}=0.
\end{align}
This directly implies, from Stokes's theorem, that in magnetostatics,
\begin{align}
    \oint \vec{E}\cdot d\vec{l}=0.
\end{align}
However, when $\Phi_B$ varies, some new Physics arises. We shall revisit this part. Put a pin on that.
\section{Biot-Savart Law}
Consider a straight current-carrying conductor. We'll take a small segment of the conductor with length $d\vec{l}$ at a point $P$, carrying current $I$ (a small amount of current). We want to find the magnetic field $d\vec{B}$ produced by this small current element at point $P$.
\begin{figure}[H]
    \centering
    \includegraphics[width=.98\linewidth]{figs/bio-savart-law-def.pdf}
    \caption{(a) A current-length element $I$ produces a differential magnetic field at point $P$. The blue (the tail of an arrow) at the dot for point P indicates that it is directed to the page there. (b) The magnetic field vector at any point is tangent to a circle for a wire with the current into (out of) the page.}
    \label{fig:bio-savart-law-def}
\end{figure}
The magnetic field $d\vec{B}$ produced by the small current element $d\vec{l}$ at point 
$P$ can be given by a mathematical expression:
\begin{align}
    |d\vec{B}|&=\frac{\mu_0}{4\pi}\frac{dQ(\vec{v}\times\hat{r})}{r^2}\notag\\
    &=\frac{\mu_0}{4\pi}\frac{(nqAv_d)(dl\sin\phi)}{r^2}\notag\\
    dB&=\frac{\mu_0}{4\pi}\frac{Idl\sin\phi}{r^2}\\
    \therefore\vec{B}&=\frac{\mu_0}{4\pi}\bigintss\frac{Id\Vec{l}\times\hat{r}}{r^2}
\end{align}
This law describes the magnetic field generated by a \textbf{steady current}. It mathematically relates the magnetic field at a point to the current flowing through a small wire segment.

The physical significance of the Biot-Savart (\textit{Bee-o Savah}) law lies in its ability to explain how a steady current produces a magnetic field. It shows that the magnetic field at a point is proportional to the current and inversely proportional to the distance from the current element. The law also illustrates that the magnetic field generated by a current-carrying wire circulates the wire in concentric circles, with the strength of the field decreasing as the distance from the wire increases.

The Biot-Savart law is valid for steady currents and in situations where the size of the current-carrying elements is much smaller compared to the distance from the observation point. However, it breaks down under certain conditions, such as changing electric fields or time-varying currents. In those cases, more general equations, such as Ampère's law or Faraday's law, need to be considered.
\subsection{$\boldsymbol{\vec{B}}$-Field of a Straight Current Carrying Conductor}
Consider a straight current-carrying conductor with current $I$ and length $2L$. Take a small segment of the conductor with length $d\vec{l}$ at a point $P$, carrying current $i$ (small amount of current). The current $I$ is uniformly distributed along the conductor, so $I$ is proportional to $d\vec{l}$. 
\begin{figure}[H]
    \centering
    \includegraphics[width=.66\linewidth]{figs/wire-segment-charge-mag-field.pdf}
    \caption{Calculating the magnetic field $\vec{B}$ produced by a current $I$ in a long straight wire. The field at P associated with the current-length element $id\vec{l}$ is directed into the page, as shown.}
    \label{fig:wire-segment-charge-mag-field}
\end{figure}
The magnetic field $d\vec{B}$ at point $P$ due to the small current element $I$ is given by Biot-Savart's law:
\begin{align}
    \vec{B}=\frac{\mu_0}{4\pi}\bigintss\frac{Id\vec{l}\times\hat{r}}{r^2}
\end{align}
The distance $r$ from the current element $d\vec{l}$ to point P is related to the distance $x$ away, perpendicularly from the conductor: $r=\sqrt{x^2+y^2}$.
\begin{align}
    B&=\frac{\mu_0I}{4\pi}\bigintss_{-L}^L\frac{xdy}{\left(x^2+y^2\right)^{\frac{3}{2}}}\notag\\
    &=\frac{\mu_0I}{4\pi}\frac{2L}{x\sqrt{x^2+L^2}}.
\end{align}
For a more realistic scenario where $L\gg x$, $L\to\infty$ renders the terms inside the square root to $x^2+L^2\sim L^2$. In this limit:
\begin{align}
    B=\frac{\mu_0I}{2\pi x}.
\end{align}
A more general case would be to denote $x$ as $r$, where $r$ is the distance perpendicularly away from the conductor where the magnetic field is being measured.
\begin{align}
    B=\frac{\mu_0I}{2\pi r}.
\end{align}
\subsection{$\boldsymbol{\vec{B}}$-Field of a Curved Current Carrying Conductor (At the Center)}
Consider a circular arc of current-carrying wire with a radius $R$ and a current $I$. We'll take a small segment of the arc with an angular extent $d\phi$ at a point $P$, carrying current. We want to find the magnetic field $d\vec{B}$ produced by this small current element at $P$.
\begin{figure}[H]
    \centering
    \includegraphics[width=\linewidth]{figs/curved-wire-mag-field.pdf}
    \caption{(a) A wire in the shape of a circular arc with center $C$ carries current $I$. (b) For any element of wire along the arc, the angle between the directions of and is $90^\circ$. (c) The field is out of the page, as indicated by the colored dot at $C$.}
    \label{fig:curved-wire-mag-field}
\end{figure}
The magnetic field $d\vec{B}$ produced by the small current element $d\vec{l}$ at point $P$ can be given by a modified version of the Biot-Savart law, considering the circular geometry:
\begin{align}
    \vec{B}&=\frac{\mu_0I}{4\pi}\bigintss\frac{d\vec{s}\times\hat{R}}{R^2}\notag\\
    &=\frac{\mu_0I}{4\pi}\bigintss\frac{Rd\phi\sin90^\circ}{R^2}\notag\\
    &=\frac{\mu_0IR}{4\pi}\bigintss_0^\phi\frac{d\phi}{R^2}\notag\\
    &=\frac{\mu_0I\phi}{4\pi R}.
\end{align}
For a complete circular wire, this becomes:
\begin{align}
    B=\frac{\mu_0I\times2\phi}{4\pi R}=\frac{\mu_0I}{2R}.\label{eq:complete-circular-wire}
\end{align}
This $\vec{B}$ is measured in the center of the circular arc wire or the circular wire.
\subsection{$\boldsymbol{\vec{B}}$-Field of a Looped Current Carrying Conductor (Axial)}
Consider a looped current-carrying conductor with a radius $R$ and a constant current $I$. We want to find the magnetic field $\vec{B}$ at a point $P$ located at a distance $r$ away from the center of the conductor (perpendicularly).
\begin{figure}
    \centering
    \includegraphics[width=.72\linewidth]{figs/looped-wire-mag-field.pdf}
    \caption{(Magnetic field on the axis of a circular loop. The current in the segment $d\vec{S}$ causes the field $d\vec{B}$ at $P$, which lies in the $xy$-plane. The currents in other $d\vec{S}$’s cause $d\vec{B}$’s with different components perpendicular to the x-axis, adding to zero. The x-components of the $d\vec{B}$’s combine to give the total $\vec{B}$ field at point $P$.}
    \label{fig:looped-wire-mag-field}
\end{figure}
Since the conductor is circular, we can express the differential magnetic field at a point $x$ distance away from the center up in the following form of Biot-Savart Law:
\begin{align}
    B&=\frac{\mu_0I}{4\pi}\bigintss\frac{dS}{\left(x^2+R^2\right)^{\frac{3}{2}}}
\end{align}
This $dB$ will then have to be resolved into its two components:
\begin{align}
    B_x&=\frac{\mu_0I}{4\pi}\bigintss\frac{dS}{\left(x^2+R^2\right)^{\frac{3}{2}}}\cos\theta\\
    B_y&=\frac{\mu_0I}{4\pi}\bigintss\frac{dS}{\left(x^2+R^2\right)^{\frac{3}{2}}}\sin\theta.
\end{align}
We can neglect the vertical component since they cancel out during the integration due to mirror symmetry. Take the orientation of the coil's axis in an arbitrary direction $\hat{r}$.
\begin{align}
    B=\frac{\mu_0I}{4\pi}\bigintss\frac{dS}{\left(r^2+R^2\right)}\cos\theta&=\frac{\mu_0I}{4\pi}\bigintss\frac{dS}{\left(r^2+R^2\right)}\frac{R}{\sqrt{r^2+R^2}}\notag\\
    &=\frac{\mu_0IR}{4\pi}\bigintss\frac{dS}{\left(r^2+R^2\right)^{\frac{3}{2}}}\notag\\
    &=\frac{\mu_0IR}{4\pi}\frac{1}{\left(r^2+R^2\right)^{\frac{3}{2}}}\bigintss_0^{2\pi R} dS\notag\\
    &=\frac{\mu_0IR^2}{2}\frac{1}{\left(r^2+R^2\right)^{\frac{3}{2}}}
\end{align}
In the limit $R\gg r$:
\begin{align}
    B=\frac{\mu_0I}{2R}.
\end{align}
This is identical to Eq. (\ref{eq:complete-circular-wire}).
\section{Magnetic Force Caused by $\vec{B}$-Fields}
The force on a moving charge $q$ in a magnetic field $B$ can be described by the equation:
\begin{align}
    \vec{F}_B=qvB\sin\phi=q\Vec{v}\times\Vec{B},\label{eq:mag-field-force}
\end{align}
where $F_B$ is the Magnetic force, $v$ is the velocity of the charge, and $\phi$ is the angle between the velocity vector and the magnetic field vector.

An external electric charge will also experience an electrostatic force $\vec{F}_E=q\vec{E}$.
\begin{figure}[H]
    \centering
    \includegraphics[width=0.7\linewidth]{figs/mag-force-def.pdf}
    \caption{The magnetic force $\vec{F}_B$ acting on a positive charge $q$ moving with velocity $\vec{v}$ is perpendicular to both $\vec{v}$ and the magnetic field $\vec{B}$. For given values of speed v and magnetic field strength $B$, the force is greatest when $\vec{v}$ and $\vec{B}$ are perpendicular.}
    \label{fig:mag-force-def}
\end{figure}
The total force on the external charge is then
\begin{align}
    \vec{F}&=\vec{F}_B+\vec{F}_E\notag\\
    &=q\Vec{v}\times\Vec{B}+q\vec{E}\notag\\
    &=q\left(\Vec{v}\times\Vec{B}+\vec{E}\right).
\end{align}
This is called the Lorentz force.
\subsection{Magnetic Force on a Current-Carrying Conductor}
Consider a current-carrying conductor with a current $I$ (Figure \ref{fig:mag-force-def}). We'll take a small segment of the conductor with a length $d\vec{l}$ at a point $P$. The current segment $d\vec{l}$ experiences a magnetic force $\vec{F}_B$ due to the magnetic field $\vec{B}$ at that point.
\begin{align}
    \vec{F}_B&=qv_dB\sin\phi\notag\\
    &=(nAl)(qv_d B)\notag\\
    &=(nqv_dA)(lB_\perp)\notag\\
    &=IlB_\perp\notag\\
    &=IlB\sin\phi\notag\\
    &=I\vec{l}\times\vec{B}.
\end{align}
\begin{figure}[H]
    \centering
    \includegraphics[width=0.7\linewidth]{figs/mag-force-current-wire.pdf}
    \caption{(a) A close-up view of a section of the wire. The current direction is upward, which means that electrons drift downward. A magnetic field that emerges from the plane of the page causes the electrons and the wire to be deflected to the right. (b) A flexible wire passes between the pole faces of a magnet (only the farther pole face is shown). (b1) Without current in the wire, the wire is straight. (b2) With an upward current, the wire is deflected rightward. (b3) With a downward current, the deflection is leftward. The connections for getting the current into the wire at one end and out of it at the other end are not shown.}
    \label{fig:mag-force-def}
\end{figure}
Magnetic force on an infinitesimal
wire segment:
\begin{align}
    d\vec{F}_B=Id\vec{l}\times\vec{B}.
\end{align}
Remember, the current $I$ is not a vector. The direction of the current flow is described by $d\vec{l}$, not $I$. If the conductor is curved, $I$ is the same at all points along its length, but $d\vec{l}$ changes direction\textemdash. It is always tangent to the conductor.
\subsection{Magnetic Force on a Current-Carrying Loop}
\begin{figure}[H]
    \centering
    \includegraphics[width=0.75\linewidth]{figs/mag-torque-current-wire-loop.pdf}
    \caption{A rectangular loop, of length $a$ and width $b$ and carrying a current $I$, is located in a uniform magnetic field. A torque $\vec{\tau}$ acts to align the normal vector with the direction of the field. (a) The loop as seen by looking in the direction of the magnetic field. (b) A perspective of the loop showing how the right-hand rule gives the direction of $\vec{\tau}$, perpendicular to the loop's plane. (c) A side view of the loop from side 2. The loop rotates as indicated. (d) The elements of an electric motor. A rectangular loop of wire, carrying a current and free to rotate about a fixed axis, is placed in a magnetic field. Magnetic forces on the wire produce a torque that rotates it. A commutator (not shown) reverses the direction of the current every half-revolution so that the torque always acts in the same direction.}
    \label{fig:mag-torque-current-wire-loop}
\end{figure}
Consider a rectangular loop of wire with side lengths $a$ and $b$. A line perpendicular to the plane of the loop (i.e., a normal to the plane) makes an angle $\theta$ with the direction of the magnetic field $\vec{B}$, and the loop carries a current $I$. The wires leading the current into and out of the loop and the source of EMF are omitted to keep the diagram simple.

The force $\vec{F}$ on the right side of the loop (length $a$) is to the right, in the $+x$ direction, as
shown. On this side, $\vec{B}$ is perpendicular to the current direction, and the force on this side
has magnitude:
\begin{align}
    F_1=IaB.
\end{align}
A force $\vec{F}_3=-\vec{F}_1$ with the same magnitude but in opposite directions acts on the opposite side of the loop. The sides with length $b$ make an angle $90^\circ - \theta$ with the direction of $\vec{B}$. The forces on these sides are the vectors $\vec{F}_3$ and $\vec{F}_4=-\vec{F}_2$; their magnitude $\vec{F}_2$ and $\vec{F}_4$ are given by
\begin{align}
    F_3=IbB\sin(90^\circ-\theta)=IdB\cos\theta
\end{align}
The lines of action of both forces lie along the $y$ axis. The total force on the loop is \textit{zero} because the forces on opposite sides cancel out in pairs.

\textsc{Takeaway}: \textbf{The net force on a current loop in a uniform magnetic field is zero. However, the net torque is not, in general, equal to zero.}

\subsection{Magnetic Torque on a Current-Carrying Loop and Electromagnet}
The two forces $\vec{F}_1$ and $\vec{F}_3$ in Figure \ref{fig:mag-torque-current-wire-loop}(c) lie along the same line and so give rise to zero net torque with respect to any point. The two forces $\vec{F}_2$ and $\vec{F}_4$ lie along different lines, and each gives rise to a torque about the $y$-axis. 
\begin{figure}[H]
    \centering
    \includegraphics[width=.8\linewidth]{figs/mag-moment-in-current-loop.pdf}
    \caption{The magnetic moment represents the strength and orientation of a magnetic source, such as a magnet or a current loop. It helps us understand how \textit{magnetically strong} and in which direction a magnetic object or system influences its surroundings. We use the right-hand palm rule to find the direction of $\vec{\mu}$. Curl the four fingers in the current direction, and the thumb will point along $\vec{\mu}$. The magnetic field lines work as if the moment was a bar magnet from the North Pole to the South. $\vec{\mu}$ points from the South Pole to the North.}
    \label{fig:mag-moment}
\end{figure}
According to the right-hand rule for determining the direction of torques, the vector torques due to $\vec{F}_2$ and $\vec{F}_4$ are both in the $+y$ direction; hence the net vector torque $\vec{\tau}$ is in the $+y$ direction as well. The moment arm for each of these forces (equal to the perpendicular distance from the rotation axis to the line of action of the force) is $\displaystyle \frac{b}{2}\sin\theta$, so the torque due to each force has magnitude $\displaystyle F\frac{b}{2}\sin\theta$. 

The net torque is then given by the following:
\begin{align}
    |\vec{\tau}|&=2F\frac{b}{2}\sin\theta\notag\\
    &=(IBa)(b\sin\theta)\notag\\
    \tau&=IBA\sin\theta\notag\\
    \vec{\tau}&=I\vec{B}\times\vec{A}.
\end{align}
The product $IA$ is defined as the \textbf{magnetic dipole moment} or \textbf{magnetic moment} $\vec{\mu}$ of the loop. It is analogous to the electric dipole moment.

In terms of $\mu$, the net torque can be expressed as
\begin{align}
    |\vec{\tau}|&=\mu B\sin\theta\\
    \vec{\tau}&=\vec{\mu}\times\vec{B}.
\end{align}
where $\theta$ is the angle between the normal to the loop (the direction of the vector area $\vec{A}$) and $\vec{B}$. A current loop, or any other object that experiences a magnetic torque, is also called a magnetic dipole. $\vec{\mu}$ points in the direction of the normal vector to the plane of the current loop.
\section{Charged Particles in $\vec{B}$-Field}
\subsection{Circular Deflection of Charged Particles}
\subsection{The Thompson Experiment: Discovery of Electron}
The cross-field cathode ray tube experiment, also known as the Thomson experiment, was conducted by J.J. Thomson in 1897. This experiment played a crucial role in the discovery of the electron and provided strong evidence for the existence of subatomic particles.

Thomson used a cathode ray tube, a sealed glass tube containing a vacuum with two metal electrodes. The cathode, a negatively charged electrode, emitted a stream of electrons, known as cathode rays, towards the anode, a positively charged electrode. Additionally, perpendicular to the path of the cathode rays, Thomson introduced a magnetic field and an electric field.

The basic idea behind the experiment was to observe the deflection of the cathode rays in the presence of both electric and magnetic fields. The magnetic force $\vec{F}_B$ on a charged particle moving through a magnetic field $\vec{B}$ is given by Eq. (\ref{eq:mag-field-force}):
\begin{align}
    F_B=qvB.
\end{align}
\begin{figure}
    \centering
    \includegraphics[width=0.82\linewidth]{figs/thomson-device-cross-field.pdf}
    \caption{A modern version of J. J. Thomson’s apparatus for measuring the ratio of mass to charge for the electron. An electric field is established by connecting a battery across the deflecting plate terminals. The magnetic field is set up by means of a current in a system of coils (not shown).}
    \label{fig:cathode-ray-tube}
\end{figure}
The electric force $\vec{F}_E$ on a charged particle moving through an electric field $\vec{E}$ is given by:
\begin{align}
    F_E=qE.
\end{align}
Thompson first applied the $\vec{B}$ field and recorded a deflection by the charged particle. He then applied an $\vec{E}$ field perpendicularly to $\vec{B}$ and tuned it till the deflection was nullified. Such effect claims:
\begin{align}
    F_B&=F_E\notag\\
    qvB&=qE\notag\\
    v&=\frac{E}{B}.
\end{align}
The accelerating voltage $V$ that was applied to the particle stored the gained electric potential energy in the form of kinetic energy due to acceleration.
\begin{align}
    qV&=\frac{1}{2}mv^2\notag\\
    \frac{q}{m}&=\frac{1}{2}\frac{v^2}{V}\notag\\
    &=\frac{1}{2}\frac{E^2}{B^2V}.\\
    &=\frac{1}{2}\frac{V^2}{B^2VL^2}=\frac{1}{2}\frac{V}{B^2L^2}
\end{align}
\subsection{Particle Accelerators: The Cyclotron}
A particle accelerator is a device that accelerates charged particles, such as protons or electrons, to high speeds and energies. The cyclotron is a type of circular accelerator. It consists of two hollow, D-shaped metal electrodes called dees, which are placed in a strong magnetic field. The dees are connected to a high-frequency alternating voltage source, creating an electric field that keeps the particle accelerating.
\begin{figure}
    \centering
    \includegraphics[width=0.82\linewidth]{figs/cyclotron.pdf}
    \caption{The elements of a cyclotron, showing the particle source S and the dees. A uniform magnetic field is directed up from the plane of the page. Circulating protons spiral outward within the hollow dees, gaining energy every time they cross the gap between the dees.}
    \label{fig:cyclotron}
\end{figure}
The charged particles, say, a proton, are injected at the center of the cyclotron, and the electric field between the dees accelerates them. As the particles gain energy, they move in a circular path due to the magnetic field. \textbf{The magnetic field acts as a constant force perpendicular to the particle's velocity}, \textbf{causing it to undergo circular motion according to the Lorentz force equation}:
\begin{align}
    \vec{F}_B&=q\vec{v}\times\vec{B}.
\end{align}
As the particles complete each circular orbit in the dees, the electric field polarity switches at precisely the right moment to accelerate the particles further. This process repeats with each half-cycle of the alternating voltage, allowing the particles to gain energy with each pass through the dees.

The particle has a centripetal force given by:
\begin{align}
    F_c=\frac{mv^2}{R}.
\end{align}
Since the centripetal force $F_c$ is provided by the magnetic force $F_B$, we can equate these forces:
\begin{align}
    q|v|B&=\frac{mv^2}{R}\notag\\
    R&=\frac{mv}{|q|B}.
\end{align}
The period $T$ (the time for one full revolution) is equal to the circumference divided by the speed:
\begin{align}
    T=\frac{2\pi R}{v}=\frac{2\pi}{v}\frac{mv}{|q|B}=\frac{2\pi m}{|q|B}.
\end{align}
The frequency $f$ (the number of revolutions per unit time) is
\begin{align}
    f=\frac{1}{T}=\frac{|q|B}{2\pi m}.
\end{align}
It is also known as the \textit{Cyclotron Frequency}.

The angular frequency $\omega$ of the motion is then
\begin{align}
    \omega=2\pi f=\frac{|q|B}{m}.
\end{align}
\textsc{Note}: The quantities $T$, $f$, and $v$ do not depend on the speed of the particle (provided the speed is much less than the speed of light). Fast particles move in large circles and slow ones in small circles, but all particles with the same charge-to-mass ratio $\displaystyle\frac{q}{m}$ take the same time $T$ (the period) to complete one round trip.
\subsection{The Hall Effect}
The Hall Effect occurs when a current-carrying conductor (usually a metal or semiconductor) is placed in a magnetic field perpendicular to the direction of the current flow. When this happens, an electric field is induced across the conductor, perpendicular to both the current and magnetic field directions. This leads to a voltage difference, known as the Hall voltage, across the width of the conductor.
\begin{figure}[H]
    \centering
    \includegraphics[width=0.82\linewidth]{figs/hall-effect.pdf}
    \caption{A strip of copper carrying a current i is immersed in a magnetic field. (a) The situation immediately after the magnetic field is turned on. The curved path that an electron will then take is shown. (b) The situation is at equilibrium, which quickly follows. Note that negative charges pile up on the right side of the strip, leaving uncompensated positive charges on the left. Thus, the left side has a higher potential than the right one. (c) For the same current direction, if the charge carriers were positively charged, they would pile up on the right side, and the right side would be at a higher potential.}
    \label{fig:hall-effect}
\end{figure}
Imagine a flat metal (Figure \ref{fig:hall-effect}) (charge carriers are electrons) plate with length $l$, width $d$, and thickness $t$ through which an electric current $I$ flows from one edge to the other. Suppose we apply a magnetic field $B$ perpendicular to the plate. In that case, the moving charge carriers (electrons in metals) experience a magnetic force $\vec{F}_\text{magnetic}=q\vec{v}\times\vec{B}$ that acts perpendicular to both the current direction and the magnetic field direction. 

Due to this force, the electrons are pushed to one side of the plate, creating an excess of negative charges on one edge and a deficit of negative charges on the other edge. This separation of charges establishes an electric field $\vec{E}$ within the plate, pointing from the negative charge side to the positive charge side. It exerts an electric force $\vec{F}_\text{electric}=q\vec{E}$ on the charge carriers opposite to the magnetic force $\vec{F}_\text{magnetic}$: The induced electric field (Hall Field) will eventually balance the magnetic force on the electrons, resulting in a steady-state condition with a constant Hall voltage $V_\text{Hall}$ across the plate. This voltage difference can be measured and is proportional to the magnetic field strength and the current in the conductor.

Since the system reaches a steady state, the net force on the charge carriers is zero:
\begin{align}
    \vec{F}_\text{magnetic}+\vec{F}_\text{electric}&=0\notag\\
    q\vec{v}\times\vec{B}&=q\vec{E}_\text{Hall}\notag\\
    \vec{E}_\text{Hall}&=-\vec{v}_d\vec{B}_\perp\notag\\
    &=-\frac{\vec{J}\,\vec{B}_\perp}{nq}\quad \because \vec{J}=\frac{I}{A}=\frac{nq\vec{v}_dA}{A}=nq\vec{v}_d.
\end{align}
When $q$ is positive, $\vec{E}$ is negative. When $q$ is negative, $\vec{E}$ is positive (depending on the dimension). 

When positive charge carriers are present, they accumulate at the upper edge, creating a potential difference opposite to negative charges. Shortly after the Hall effect was discovered in 1879, it was noted that some materials, particularly semiconductors, exhibited a Hall EMF that was opposite to that of metals, as if their charge carriers were positively charged. It is now understood that these materials are conducted through a process called \textit{hole conduction}. Within these materials, there are locations known as holes that would usually be occupied by electrons but are actually vacant. The absence of a negative charge is equivalent to the presence of a positive charge. When an electron moves to fill a hole, it leaves another hole behind it, causing the hole to migrate in the opposite direction of the electron. 

The Hall voltage $V_\text{Hall}$ across the width of the plate is given by the electric field multiplied by the width $d$:
\begin{align}
    V_\text{Hall}&=E_\text{Hall}d=-\left(\frac{1}{nq}\right)(\vec{J}\,\vec{B}_\perp)d=-R_H\vec{J}\,\vec{B}_\perp d,
\end{align}
where $R_H$ is the Hall coefficient for the specific material.

\textbf{Note}: The $\perp$ sign with $\vec{B}$ implies that $\vec{B}$ is always perpendicular to the vector ($\vec{E}_\text{Hall}$ or $\vec{J}$) in front of it.
\section{Ampère's Law}
Ampère's Law is a fundamental principle in electromagnetism that relates the magnetic field around a closed loop to the electric current passing through the loop. 
\begin{figure}[H]
    \centering
    \includegraphics[width=.9\linewidth]{figs/amprere-law.pdf}
    \caption{Using Ampère’s Law to find the magnetic field that a current $i$ produces outside a long straight wire of circular cross-section. The Amperian loop is a concentric circle that lies outside the wire.}
    \label{fig:ampere-law}
\end{figure}
It states that the line integral of the magnetic field along a closed loop is equal to the product of the permeability of free space and the total current enclosed by the loop.

Mathematically, Ampère's Law is expressed as:
\begin{align}
    \oint\Vec{B}\cdot d\Vec{l}=\mu_0 I_\text{enclosed}.\label{eq:ampere-law}
\end{align}
Ampère's Law is a powerful tool for calculating the magnetic field in cases with a high degree of symmetry, such as when dealing with cylindrical or spherical symmetry. It allows us to determine the magnetic field produced by a current-carrying wire, a solenoid, or other geometries with symmetry.

This demonstrates the intimate relationship between electricity and magnetism, indicating that an electric current produces a magnetic field, and the strength of the field is directly proportional to the current.

Once again, just like Gauss' Law, Ampère's Law is only valid when the electric current is steady and the magnetic field is constant over time. It breaks down when there are time-varying electric fields or changing magnetic fields, as in cases involving rapidly changing currents or electromagnetic induction. In such scenarios, Maxwell's extension of Ampère's Law, known as Maxwell's equations, must accurately describe the electromagnetic phenomena. We will get to it shortly.
\subsection*{Convention for Current while using Ampère's Law}
Curl your right hand around the Amperian loop, with the fingers pointing in the direction of integration. A current through the loop in the general direction of your outstretched thumb is assigned a plus sign, and a current generally in the opposite direction is assigned a minus sign.
\begin{align*}
    \oint B_\parallel dl=\mu_0 (i_1+i_2-i_3).
\end{align*}
\subsection*{Condition to Choose Appropriate Amperian Loop}
Choosing the appropriate Amperian loop is crucial when applying Ampère's Law.

To ensure accurate and meaningful results when using Ampère's Law, selecting an Amperian loop that satisfies certain conditions is important. Here are the key considerations:
\begin{itemize}
\setlength{\itemsep}{0pt}
    \item Symmetry
    \item Uniform $\vec{B}$-Field at all points on the loop
    \item Avoid Singularities
\end{itemize}
\subsection{$\boldsymbol{\vec{B}}$-Field Outside a Long Straight Current-Carrying Wire}
Consider a long straight wire carrying a current $I$. We want to find the magnetic field $\vec{B}$ at a distance $|\vec{r}|$ from the wire, where $r$ is greater than the radius of the wire. Choose an Amperian loop of radius $r$ concentric with the wire's center. Due to symmetry, the field $\vec{B}$ will remain constant across the Amperian loop.

Apply Ampère's Law:
\begin{align}
    \oint B_\parallel dA&=\mu_0 I_\text{enc}=B\oint dl=B\times 2\pi r\notag\\
    B\times 2\pi r&=\mu_0 I_\text{enc}\notag\\
    B&=\frac{\mu_0 I_\text{enc}}{2\pi r}
\end{align}
\subsection{$\boldsymbol{\vec{B}}$-Field Inside a Long Straight Current-Carrying Wire}
Consider a long straight wire carrying a current $I$. We want to find the magnetic field $\vec{B}$ at a distance $|\vec{r}|$ from the wire, where $r$ is smaller than the radius of the wire. Choose an Amperian loop of radius $r$ concentric with the wire's center. Due to symmetry, the field $\vec{B}$ will remain constant across the Amperian loop.

Here,
\begin{align*}
    I_\text{enc}=I_\text{actual}\times\frac{\pi r^2}{\pi R^2}.
\end{align*}
\begin{figure}[H]
    \centering
    \includegraphics[width=\linewidth]{figs/straight-wire-ampere-law.pdf}
    \caption{To find the magnetic field at radius $r<R$, we apply Ampère’s Law to the circle enclosing the gray area. The current through the gray area is $\displaystyle I_\text{actual}\times\frac{\pi r^2}{\pi R^2}$. To find the magnetic field at radius $r>R$, we apply Ampère’s Law to the circle enclosing the entire conductor.}
    \label{fig:straight-wire-ampere-law}
\end{figure}
Apply Ampère's Law:
\begin{align}
    \oint B_\parallel dA&=\mu_0 I_\text{enc}=B\oint dl=B\times 2\pi r\notag\\
    B\times 2\pi r&=\mu_0 I_\text{actual}\times\frac{\pi r^2}{\pi R^2}\notag\\
    B&=\left(\frac{\mu_0 I_\text{enc}}{2\pi R^2}\right)r
\end{align}
\begin{center}
\begin{tabular}{cl}
    \textbf{Region} & \textbf{$\vec{B}$-Field} \\[3pt]
    $r<R$ & $\displaystyle B=\left(\frac{\mu_0 I_\text{enc}}{2\pi R^2}\right)r$ \\[5pt]
    $r=R$ & $\displaystyle B=\frac{\mu_0 I_\text{enc}}{2\pi R}$ \\[8pt]
    $r>R$ & $\displaystyle B=\frac{\mu_0 I_\text{enc}}{2\pi r}$
\end{tabular}
\end{center}
\subsection{$\boldsymbol{\vec{B}}$-Field of a Solenoid}
Consider a solenoid, which is a long cylindrical coil of wire with $N$ turns per unit length and carrying a current $I$. We want to find the magnetic field $\vec{B}$ inside the solenoid.

Choose a rectangular loop that passes through the center of the solenoid and is parallel to the axis of the solenoid. The magnetic field $\vec{B}$ inside the solenoid is uniform and directed along the axis of the solenoid. The length of the closed path $l$ of the rectangular loop equals the length of the solenoid $L$.
\begin{figure}[H]
    \centering
    \includegraphics[width=.9\linewidth]{figs/solenoid-mag-field-ampere-law.pdf}
    \caption{(a) A vertical cross-section through the central axis of a “stretched-out” solenoid. The back portions of five turns are shown, as are the magnetic field lines due to a current through the solenoid. Each turn produces circular magnetic field lines near itself. Near the solenoid’s axis, the field lines combine into a net magnetic field that is directed along the axis. The closely spaced field lines there indicate a strong magnetic field. Outside the solenoid the field lines are widely spaced; the field there is very weak. (b) Magnetic field lines for a real solenoid of finite length. The field is strong and uniform at interior points such as $P_1$ but relatively weak at external points such as $P_2$. (c) Application of Ampère’s Law to a section of a long ideal solenoid carrying a current $I$.The Amperian loop is the rectangle $abcda$.}
    \label{fig:solenoid-mag-field-ampere-law}
\end{figure}
Here,
\begin{align*}
    I_\text{enc}=NI_\text{actual}L.
\end{align*}
Apply Ampère's Law:
\begin{align}
    \oint B_\parallel dA&=\mu_0 I_\text{enc}=B\oint dl=BL\notag\\
    B L&=\mu_0 \times NI_\text{actual}L\notag\\
    B&=\mu_0NI_\text{actual}.
\end{align}
\subsection{$\boldsymbol{\vec{B}}$-Field of a Toroid}
Consider a toroid, which is a hollow, circular, ring-shaped coil of wire with multiple turns. The toroid has a mean radius $R$ and a cross-sectional area $A$. The current flowing through the toroid is $I$. We want to find the magnetic field $\vec{B}$ inside the toroid.

Choose a circular loop that lies in the same plane as the toroid and passes through the center of the toroid. The magnetic field $\vec{B}$ inside the toroid is circular and concentric with the toroid's circular cross-section. A toroid is clearly may describe as a (hollow) solenoid that has been curved until its two ends meet, forming a sort of hollow bracelet.
\begin{figure}[H]
    \centering
    \includegraphics[width=.85\linewidth]{figs/toroid-mag-field-ampere-law.pdf}
    \caption{(a) A toroidal solenoid. For clarity, only a few turns of the winding are shown. (b) Integration paths (black circles) are used to compute the magnetic field $\vec{B}$ set up by the current (shown as dots and crosses).}
    \label{fig:toroid-mag-field-ampere-law}
\end{figure}
Apply Ampère's Law:
\begin{align}
    \oint B_\parallel dA&=\mu_0 I_\text{enc}=B\oint dl=B\times2\pi r\notag\\
    B\times2\pi r&=\mu_0NI_\text{actual}\notag\\
    B&=\frac{\mu_0NI_\text{actual}}{2\pi r}.
\end{align}
\clearpage
%%%%%%%%%%%%%%%%%%%%%%%%%%%%%%%%%%%%%%%%%%%%%%%%%%%%%%%%%%%%%%%%%%%%%
\bigskip \
\vspace{10cm}
\vskip 20pt
\begin{center}
    \thispagestyle{empty}
    \huge{\selectfont \bf \sffamily Part III: Uniting $\boldsymbol{\vec{E}}$ and $\boldsymbol{\vec{B}}$-fields}
\end{center}
\addcontentsline{toc}{section}{\Large Part III: Uniting $\boldsymbol{\vec{E}}$ and $\boldsymbol{\vec{B}}$-fields}
\vfill
\clearpage
%%%%%%%%%%%%%%%%%%%%%%%%%%%%%%%%%%%%%%%%%%%%%%%%%%%%%%%%%%%%%%%%%%%%%
\section{Electromagnetic Induction}
\thispagestyle{empty}
Electromagnetic induction refers to the process by which an electric current is induced in a conductor when exposed to a changing magnetic field. This phenomenon was first discovered by Michael Faraday in the early \nth{19} century and is governed by Faraday's Law of electromagnetic induction and Lenz's Law.

A British scientist, Michael Faraday, conducted a famous experiment in the early \nth{19} century to demonstrate electromagnetic induction. He observed that when he moved a magnet (Experiment \#1) through a coil of wire or changed the magnetic field strength near the coil, a current was induced in the wire. The same thing happens when a current-carrying loop is brought near a coil (Experiment \#2). There appears a current only when the incoming circuit opens and/or closes. These groundbreaking experiments showed that a changing magnetic field could create an electric current in a conductor, revealing the relationship between electricity and magnetism.
\subsection{Induced EMF and Faraday's Law}
Induced EMF refers to the electromotive force or voltage generated in a conductor due to electromagnetic induction. When a conductor is exposed to a changing magnetic field or experiences a relative motion with respect to a magnetic field, the magnetic flux linking the conductor changes, leading to the generation of an induced EMF. This phenomenon is the basis for generating electrical power in generators and is essential in operating transformers. The induced EMF $\mathcal{E}$ in the conductor is given by:
\begin{align}
    \mathcal{E} = - \frac{d\Phi_B}{dt}, \label{eq:faraday's-law}
\end{align}
\begin{figure}[H]
    \centering
    \includegraphics[width=\linewidth]{figs/em-induction.pdf}
    \caption{(a) An ammeter registers a current in the wire loop when the magnet is moving with respect to the loop. (b) An ammeter registers a current in the left-hand wire loop just as switch $S$ is closed (to turn on the current in the righthand wire loop) or opened (to turn off the current in the right-hand loop). No motion of the coils is involved.}
    \label{fig:em-induction}
\end{figure}
Eq. (\ref{eq:faraday's-law}) is known as \textbf{Faraday's Law} of electromagnetic induction. It states that \textbf{the magnitude of the induced electromotive force (EMF) in a circuit is directly proportional to the rate of change of magnetic flux passing through the circuit}. 

Faraday's Law tells us that \textbf{a changing magnetic field near a conductor induces an electric field within the conductor, which in turn causes an electric current to flow}. This principle forms the basis for operating devices such as electric generators and transformers.
\subsection{Lenz's Law}
Lenz's Law, which is a consequence of Faraday's Law, states that the direction of the induced current in a conductor is such that it opposes the change in magnetic flux that produced it. In other words, the induced current creates a magnetic field that opposes the original change in a magnetic field. \textbf{This Law can be understood as a consequence of the conservation of energy}. Lenz's Law ensures that the induced current and magnetic field always acts in a way that opposes the original cause of the change in flux.
\subsubsection*{How to Find the Direction of the Induced EMF Using Lenz's Law}
\begin{figure}[H]
    \centering
    \includegraphics[width=\linewidth]{figs/lenz's-law.pdf}
    \caption{The direction of the current $I$ induced in a loop is such that the current’s magnetic field $\vec{B}_\text{ind}$ opposes the change in the magnetic field $\vec{B}$ inducing $I$. The field $\vec{B}_\text{ind}$ is always directed opposite an increasing field $\vec{B}$ (a, c) and in the same direction as a decreasing field $B$ (b, d). The curled–straight right-hand rule gives the direction of the induced current based on the direction of the induced field.}
    \label{fig:lenz's-law}
\end{figure}
\begin{enumerate}
    \item Substitute the \textit{incoming} current loop (Experiment \#2) or bar magnet (Experiment \#1) with a magnetic moment $\vec{\mu}$ aligned in the direction from the south pole of a bar magnet towards its north pole (similar to the orientation of the bar magnet). 
    \item In accordance with Lenz's Law, which necessitates the induced magnetic field to \textit{oppose} the change introduced by the incoming components, introduce a current in the induced loop such that the poles of $\vec{\mu}_\text{ind}$ are oriented in a direction \textit{opposite} to the main magnetic moment $\vec{\mu}$ (in contrast to the direction of the main magnetic field $\vec{B}$).
\end{enumerate}
\subsection{A Reformulation of Faraday’s Law}
Consider a charge carrier of charge $q_0$ moving around the looped conducting wire. The work done $W$ done on it in one complete revolution by the induced electric field is $W=\mathcal{E}q$.
\begin{align}
    W&=(q_0E)(2\pi r)\notag\\
    &=q_0\int(2\pi rE)=\oint\vec{F}\cdot d\vec{s}\notag\\
    &=q_0\oint\vec{E}\cdot d\vec{l}\notag\\
    \frac{W}{q_0}&=\mathcal{E}=\oint\vec{E}\cdot d\vec{l}\notag\\
    \therefore\oint\vec{E}\cdot d\vec{l}&=-\frac{d\Phi_B}{dt}.\label{eq:faraday-law-integral}
\end{align}
Apply Stoke's theorem to (\ref{eq:faraday-law-integral}) to get the differential form of Faraday's Law.
\begin{align}
    \oint_S(\vec{\nabla}\times\vec{E})\cdot d\vec{A}&=-\oint_S\cdot\frac{\partial}{\partial t}(\vec{B}\cdot d\vec{A})\notag\\
    \vec{\nabla}\times\vec{E}&=-\frac{\partial\vec{B}}{\partial t}.\label{eq:faraday-law-differential}
\end{align}
Eq. (\ref{eq:faraday-law-differential}) is the \nth{3} Maxwell's Equation that says there should exist a divergence of the electric field if the magnetic field is varying. No variations in $\vec{B}$ give zero divergence.
\subsection{Magnetic Properties of Matter}
Magnetic properties of matter arise from the interactions between the magnetic fields generated by the atomic or molecular structure of materials and external magnetic fields. There are three main types of magnetic behavior exhibited by materials: diamagnetism, paramagnetism, and ferromagnetism.
\begin{itemize}
\setlength{\itemsep}{0pt}
    \item Diamagnetism: Diamagnetic materials exhibit a weak repulsion when placed in an external magnetic field. This behavior arises due to the induced magnetic fields that oppose the applied field. In diamagnetic materials, the electron orbital motion within atoms or molecules creates small, individual magnetic fields that align in the opposite direction to the external field. As a result, the net magnetic moment of the material is reduced, leading to repulsion. Examples of diamagnetic substances include water, wood, and most organic compounds.
    \item Paramagnetism: Paramagnetic materials are weakly attracted to an external magnetic field. Unlike diamagnetic materials, paramagnetic substances possess unpaired electrons in their atomic or molecular orbitals. When exposed to a magnetic field, these unpaired electrons align their magnetic moments in the direction of the field, resulting in a net attraction. However, the alignment is weak and disappears when the external field is removed. Paramagnetic behavior is observed in aluminum, platinum, and oxygen materials.
    \item Ferromagnetism: Ferromagnetic materials are strongly attracted to an external magnetic field and can retain magnetization even after removing the field. In ferromagnetic substances, groups of atoms or ions called magnetic domains have aligned magnetic moments that reinforce each other, creating a macroscopic magnetization. The alignment is a result of interactions between neighboring atomic magnetic moments. This spontaneous alignment gives rise to the characteristic properties of ferromagnetic materials, such as permanent magnetism and the ability to attract other magnetic materials. Iron, Nickel, and Cobalt are examples of ferromagnetic materials.
\end{itemize}
\begin{figure}[H]
    \centering
    \includegraphics[width=\linewidth]{figs/magnetism.pdf}
    \caption{A magnified diagram of the inner workings of magnetic dipoles, the arrows show the directions of magnetization in the domains of a single crystal of Nickel. Domains that are magnetized in the direction of an applied magnetic field grow larger.}
    \label{fig:magnetism}
\end{figure}
The underlying principles that govern these magnetic behaviors involve the interaction of magnetic moments at the atomic or molecular level. In diamagnetism, the magnetic moments of individual atoms or molecules oppose the applied field due to the electron orbital motion. In paramagnetism, the unpaired electrons align with the external field, resulting in a net attraction. In ferromagnetism, the alignment of magnetic domains leads to strong and persistent magnetization.
\subsection{Motional EMF}
The motional EMF is the electromotive force induced in a circuit due to the motion of a conductor through a magnetic field. When the conductor moves through the field, the motion creates a change in magnetic flux, which induces a voltage across the circuit. This induced voltage can drive a current in the circuit, even if the circuit is initially open.

Consider a conducting rod of length $L$ moving with a constant velocity $\vec{v}$ perpendicular to a magnetic field $\vec{B}$. The rod is connected to a circuit that forms a complete loop with a total resistance $R$. As the rod moves through the magnetic field, an EMF is induced in the circuit due to the motion of the rod cutting across the magnetic field lines.
\begin{figure}[H]
    \centering
    \includegraphics[width=.86\linewidth]{figs/motional-EMF.pdf}
    \caption{You pull a closed conducting loop out of a magnetic field at constant velocity $\vec{V}$. While the loop is moving, a clockwise current $i$ is induced in the loop, and the loop segments still within the magnetic field experience forces $\vec{F}_1$, $\vec{F}_2$ and $\vec{F}_3$.}
    \label{fig:emotional-EMF-def}
\end{figure}
The EMF induced in the circuit is given by Faraday's Law of electromagnetic induction:
\begin{align}
    \mathcal{E}&=-\frac{d\Phi_B}{dt}=-\frac{d(BA)}{dt}\notag\\
    &=-B\frac{dA}{dt}=-BL\frac{dx}{dt}\notag\\
    \therefore\mathcal{E}&=-BLv.
\end{align}
The induced current in the circuit is:
\begin{align}
    i&=\frac{\left|\mathcal{E}\right|}{R}=\frac{BLv}{R}\label{motional-current}
\end{align}
The unbalanced force in the current loop with side length $L$ is the magnetic force:
\begin{align}
    |\vec{F}|=i\vec{L}\times\vec{B}=iLB\label{motional-force}
\end{align}
Mixing Eq. (\ref{motional-current}) and (\ref{motional-force}) we get the following:
\begin{align}
    F=\frac{B^2L^2v}{R}.
\end{align}
The rate of work, i.e., Power on the loop as you move it in the magnetic field:
\begin{align}
    P=Fv=\frac{B^2L^2v^2}{R}.
\end{align}
This will also be the thermal energy appearance in the loop when moving at a constant speed $v$.
\begin{align}
    P=i^2R=\frac{B^2L^2v^2}{R^2}R=\frac{B^2L^2v^2}{R}.
\end{align}
% \clearpage
\section{Inductance}
Inductance $L$ is a fundamental property of an electrical circuit or component that describes its ability to resist changes in current. It quantifies how much magnetic flux is produced per unit of current change. Inductance is typically measured in henries (\unit{\henry}), and its value depends on the geometry and arrangement of the circuit elements.
\subsection{Self Inductance}
Self-inductance refers to the property of a circuit or component to induce an electromotive force (EMF) in itself when the current through it changes. It occurs due to the changing magnetic field generated by the current.

Consider a solenoid (a coil of wire) with $N$ turns and length $L$, carrying a current $I$. When a current flows through the solenoid, it generates a magnetic field $\vec{B}$ inside the coil. The magnetic flux $\Phi_B$ is linked with each turn of the coil.

The induced magnetic flux $\Phi_B$ (One unit of it per turn of the coil) is proportional to the current $i$ in the coil.
\begin{align}
    N\Phi_B&\propto i\notag\\
    N\Phi_B&=Li\notag\\
    \therefore L&=\frac{N\Phi_B}{i}.\label{eq:self-inductance}
\end{align}
Eq. (\ref{eq:self-inductance}) defines the self-inductance or simply inductance of a current-carrying coil.
\begin{figure}[H]
    \centering
    \includegraphics[width=.83\linewidth]{figs/inductor-property.pdf}
    \caption{If the current in a coil is changed by varying the contact position on a variable resistor, a self-induced EMF $L$ will appear in the coil while the current is changing. The current $i$ is increasing, and the self-induced EMF $L$ appears along the coil in a direction such that it opposes the increase. The arrow representing $L$ can be drawn along a turn of the coil or alongside the coil. Both are shown. The current $i$ is decreasing, and the self-induced EMF appears in a direction such that it opposes the decrease.}
    \label{fig:self-inductance}
\end{figure}
\subsubsection{Induced EMF and Current Due to an Inductor}

The induced EMF $\mathcal{E}$ can now be defined using Faraday's Law of induction.
\begin{align}
    \mathcal{E}_L&=-\frac{d(N\Phi_B)}{dt}=-L\frac{di}{dt}.
\end{align}
When you increase the current $i$ in a coil, the rate of change of current $\displaystyle\frac{di}{dt}$ represents the change that self-inductance resists. Consequently, a self-induced EMF appears in the coil, opposing this increase in current, striving to maintain the initial state. Conversely, if the current decreases over time, the self-induced EMF resists this decrease, attempting to sustain the initial state. This behavior aligns with Lenz's Law.

Furthermore, a self-induced potential difference $V_L$ can be defined across an inductor assumed to be between its terminals outside the changing flux region. In the case of an ideal inductor (with negligible wire resistance), the magnitude of $V_L$ equals the magnitude of the self-induced EMF $\mathcal{E}_L$. However, if the inductor's wire has resistance $r$, we mentally separate it into a resistance $r$ (assumed outside the changing flux region) and an ideal inductor with self-induced EMF $L$. Similar to a real battery with both EMF and internal resistance $r$, the potential difference across a real inductor's terminals differs from its EMF. It's important to note that unless stated otherwise, we consider inductors to be ideal.
\subsection{Mutual Inductance}
Mutual inductance $M$ is a fundamental property of two nearby electrical circuits or components that describes how changes in the current in one circuit induce an electromotive force (EMF) in the other circuit.
\begin{figure}[H]
    \centering
    \includegraphics[width=.88\linewidth]{figs/mutual-induction.pdf}
    \caption{Mutual induction. (a) The magnetic field produced by current $i_1$ in Coil 1 extends through Coil 2. If $i_1$ is varied (by varying resistance $R$), an EMF is induced in Coil 2, and current registers on the meter connected to Coil 2. (b) The roles of the coils are interchanged.}
    \label{fig:mutual-inductance}
\end{figure}
Consider two separate coils, Coil 1 and Coil 2, placed close to each other. When the current in Coil 1 changes, it generates a changing magnetic field $\vec{B}_1$ that links with Coil 2. The magnetic flux $\Phi_{21}$ due to Coil 1's changing current linked with Coil 2 is given by the following equation, mirroring Eq. (\ref{eq:self-inductance}):
\begin{align}
    M_{21}=\frac{N_2\Phi_{12}}{i_1}.\label{eq:mutual-inductance-21}
\end{align}
Similarly, when the current in Coil 2 changes, it generates a changing magnetic field $\vec{B}_2$ that links with Coil 1. The magnetic flux $\Phi_{12}$ due to Coil 2's changing current linked with Coil 1 is given by:
\begin{align}
    M_{12}=\frac{N_1\Phi_{21}}{i_2}.\label{eq:mutual-inductance-12}
\end{align}
Combining Eqs. (\ref{eq:mutual-inductance-21}) and (\ref{eq:mutual-inductance-21}) gives us:
\begin{align}
    M&=M_{21}=M_{12}=\frac{N_1\Phi_{21}}{i_2}=\frac{N_2\Phi_{12}}{i_1}.
\end{align}
\section{Energy Stored in $\vec{B}$-Field}
The energy stored in a magnetic field can be derived using the formula for the energy stored in an inductor. When a current flows through an inductor, it generates a magnetic field, and energy is stored in this field. We begin with the power stored in the circuit.

The power in the $RL$ circuit $P_L$ can be calculated using the following steps:
\begin{align}
    P_L&=\mathcal{E}i=Li\frac{di}{dt}+i^2R,\label{eq:power-rl-circuit}
\end{align}
where the first term is the power stored in the inductor and the second term is the power dissipated across the wire resistance. We can consider the wire resistance $R$ out of the jurisdiction of the inductor.

Differentiating Eq. (\ref{eq:power-rl-circuit}), we get the energy in the circuit:
\begin{align}
    \frac{dU_B}{dt}=P_L&=Li\frac{di}{dt}\notag\\
    U_B=\int_0^{U_B}dU_B&=L\int_0^iidi\notag\\
    U_B&=\frac{1}{2}Li^2.
\end{align}
This is the magnetic field energy stored in the inductor.
\subsection{Energy Density of $\vec{B}$-Field}
Consider a length $l$ near the middle of a long solenoid of cross-sectional area $A$ carrying current $i$; the volume associated with this length is $Al$. The energy $U_B$ stored by the length $l$ of the solenoid must lie entirely within this volume because the magnetic field outside such a solenoid is approximately zero.

Moreover, the stored energy must be uniformly distributed within the solenoid because the magnetic field is (approximately) uniform everywhere inside. Thus, the energy stored per unit volume of the field is
\begin{align}
    u_B&=\frac{U}{V}=\frac{\displaystyle \frac{1}{2}Li^2}{Al}=\frac{Li^2}{2Al}=\frac{L}{l}\frac{i^2}{2A}\notag\\
    &=\frac{i^2}{2Al}\frac{N\Phi_B}{i}\notag\\
    &=\frac{i}{2Al}(nl)(BA)\notag\\
    &=\frac{i}{2Al}(nl)(\mu_0in)A\notag\\
    &=\frac{n^2i^2\mu_0}{2}=\frac{(\mu_0ni)^2}{2\mu_0}\notag\\
    \therefore u_B&=\frac{1}{2\mu_0}B^2.
\end{align}
\section{Transient Series $RL$ Circuits}
Similar to RC circuits, we define $RL$ circuits where the capacitor is replaced with an inductor with inductance $L$.
\begin{figure}[H]
    \centering
    \includegraphics[width=.75\linewidth]{figs/rl-circuit-def.pdf}
    \caption{When switch $S$ is closed on $a$, the inductor is activated through the resistor when the current changes only. When the switch is afterward closed on $b$, the inductor does not let the current drop sharply through the resistor.}
    \label{fig:rc-circuit-def}
\end{figure}
\subsection{Turn on the Circuit: Growth of Current}
Consider the circuit with the battery, resistor, and inductor connected in series. The voltage across the capacitor at any time $t$ is denoted as $V_L(t)$. The transient current in this circuit is $i(t)$.

Use the KVL for the circuit. The instantaneous potential difference $V_{ab}$ and $V_{bc}$ are: $v_{ab}=iR$ and $\displaystyle v_{bc}=-L\frac{di}{dt}$.
\begin{align}
    \mathcal{E}-iR-L\frac{di}{dt}&=0\notag\\
    \frac{di}{dt}&=\frac{\mathcal{E}-iR}{L}\notag\\
    \frac{di}{dt}&=\frac{\mathcal{E}}{L}-\frac{iR}{L}\notag\\
    \frac{di}{\displaystyle i-\frac{\mathcal{E}}{R}}&=-\frac{R}{L}dt.
\end{align}
The integration limits are $i=0$ at $t=0$ to $i=i$ at $t$,
\begin{align}
    \int_0^i\frac{di}{\displaystyle i-\frac{\mathcal{E}}{R}}&=-\int_0^t\frac{dt}{\displaystyle\frac{L}{R}}\notag\\
    \ln\left(\frac{\displaystyle i-\frac{\mathcal{E}}{R}}{\displaystyle-\frac{\mathcal{E}}{R}}\right)&=-\frac{t}{\displaystyle\frac{L}{R}}\notag\\
    \displaystyle i-\frac{\mathcal{E}}{R}&=e^{\displaystyle-\frac{t}{\displaystyle\frac{L}{R}}}\notag\\
    i&=\frac{\mathcal{E}}{R}\left(1-e^{\displaystyle-\frac{t}{\displaystyle\frac{L}{R}}}\right)\notag\\
    i&=\frac{\mathcal{E}}{R}\left(1-e^{\displaystyle-\frac{t}{\tau_L}}\right).\label{eq:increasing-current-eq-rl-circuit}
\end{align}
This is the current equation for an inductor during a current surge in a series $RL$ circuit. It shows how the current in the inductor increases with time as it charges up to the maximum charge allowed by the circuit.

The inductive time constant is given by $\displaystyle\tau_L=\frac{L}{R}$. This has the same dimensions of time similar to the capacitive time constant $\tau_C=RC$.
\begin{figure}
    \centering
    \includegraphics[width=.92\linewidth]{figs/rl-circuit-charging-cycle.pdf}
    \caption{Both currents $i$ (following the Eqs. (\ref{eq:increasing-current-eq-rl-circuit})) and voltage $V_L$ (following the Eqs. (\ref{eq:increasing-voltage-eq-rl-circuit})) increases exponentially with time $t$.}
    \label{fig:rl-charging-cycle}
\end{figure}
At $t=0$, Eq. (\ref{eq:increasing-current-eq-rl-circuit}) leads to:
\begin{align*}
    i=\frac{\mathcal{E}}{R}\left(1-e^{\displaystyle-\frac{0}{\tau_L}}\right)=0.
\end{align*}
At $t\to\infty$, Eq. (\ref{eq:increasing-current-eq-rl-circuit}) leads to:
\begin{align*}
    i=\frac{\mathcal{E}}{R}\left(1-e^{\displaystyle-\frac{\infty}{\tau_L}}\right)=\frac{\mathcal{E}}{R}.
\end{align*}
\subsection{Turn on the Circuit: Growth of Resistive Voltage}
Multiplying Eq. (\ref{eq:increasing-current-eq-rl-circuit}) with the resistor $R$ gives the voltage equation during the current surge period of the resistor.
\begin{align}
    V_R^\text{circuit-on}=iR=\mathcal{E}\left(1-e^{\displaystyle-\frac{t}{\tau_L}}\right).\label{eq:increasing-voltage-eq-rl-circuit}
\end{align}
At $t=0$, Eq. (\ref{eq:increasing-voltage-eq-rl-circuit}) leads to:
\begin{align*}
    V_R^\text{circuit-on}=\mathcal{E}\left(1-e^{\displaystyle-\frac{0}{\tau_L}}\right)=0.
\end{align*}
At $t\to\infty$, Eq. (\ref{eq:increasing-voltage-eq-rl-circuit}) leads to:
\begin{align*}
    V_R^\text{circuit-on}=\mathcal{E}\left(1-e^{\displaystyle-\frac{\infty}{\tau_L}}\right)=\mathcal{E}.
\end{align*}
\subsection{Turn on the Circuit: Decay of Inductive Voltage}
Differentiating Eq. (\ref{eq:increasing-current-eq-rl-circuit}) with time once and multiplying it by $L$ gives the voltage equation during the current surge period of the inductor.
\begin{align}
    V_L^\text{circuit-on}=L\frac{di}{dt}=\frac{L\mathcal{E}}{R}\times\frac{1}{\tau_L}\left(e^{\displaystyle-\frac{t}{\tau_L}}\right)=\mathcal{E}e^{\displaystyle-\frac{t}{\tau_L}}.\label{eq:decreasing-ind-voltage-eq-rl-circuit}
\end{align}
At $t=0$, Eq. (\ref{eq:decreasing-ind-voltage-eq-rl-circuit}) leads to:
\begin{align*}
    V_L^\text{circuit-on}=\mathcal{E}e^{\displaystyle-\frac{0}{\tau_L}}=\mathcal{E}.
\end{align*}
At $t\to\infty$, Eq. (\ref{eq:decreasing-ind-voltage-eq-rl-circuit}) leads to:
\begin{align*}
    V_L^\text{circuit-on}=\mathcal{E}e^{\displaystyle-\frac{\infty}{\tau_L}}=0.
\end{align*}
\subsection{Turn off the Circuit: Decay of Current}
Consider the same $RL$ circuit above but have the battery turn off. The current then decays slowly instead of a sharp, instantaneous drop.
\begin{align}
    -iR-L\frac{di}{dt}&=0\notag\\
    i&=-\frac{L}{R}\frac{di}{dt}\\
    \frac{di}{i}&=-\frac{dt}{\displaystyle\frac{L}{R}}\notag
\end{align}
The integration limits are $i=I_0$ at $t=0$ to $i=i$ at $t$,
\begin{align}
    \int_{I_0}^i\frac{di}{i}&=-\int_0^t\frac{dt}{\displaystyle\frac{L}{R}}\notag\\
    \ln\left(\frac{i}{I_0}\right)&=-\frac{t}{\tau_L}\notag\\
    \frac{i}{I_0}&=e^{\displaystyle-\frac{t}{\tau_L}}\notag\\
    i&=I_0e^{\displaystyle-\frac{t}{\tau_L}}.\label{eq:decreasing-current-eq-rl-circuit}
\end{align}
\begin{figure}
    \centering
    \includegraphics[width=.92\linewidth]{figs/rl-circuit-discharging-cycle.pdf}
    \caption{Both currents $i$ (following the Eqs. (\ref{eq:decreasing-current-eq-rl-circuit})) and voltage $V_L$ (following the Eqs. (\ref{eq:decreasing-voltage-eq-rl-circuit})) decreases exponentially with time $t$.}
    \label{fig:rl-discharging-cycle}
\end{figure}
At $t=0$, Eq. (\ref{eq:decreasing-voltage-eq-rl-circuit}) leads to:
\begin{align*}
    i&=I_0e^{\displaystyle-\frac{0}{\tau_L}}=I_0.
\end{align*}
At $t\to\infty$, Eq. (\ref{eq:decreasing-voltage-eq-rl-circuit}) leads to:
\begin{align*}
    i&=I_0e^{\displaystyle-\frac{t}{\tau_L}}=0.
\end{align*}
This is the decay current equation for an inductor in a series $RL$ circuit. It shows how the current in the inductor decreases with time as it decays to the minimum charge allowed by the circuit.
\subsection{Turn off the Circuit: Decay of Resistive Voltage}
Multiplying Eq. (\ref{eq:decreasing-current-eq-rl-circuit}) with the resistor $R$ gives the voltage equation during the current surge period of the inductor.
\begin{align}
    V_R^\text{circuit-off}=iR&=(I_0R)e^{\displaystyle-\frac{t}{\tau_L}}.\label{eq:decreasing-voltage-eq-rl-circuit}
\end{align}
At $t=0$, Eq. (\ref{eq:decreasing-voltage-eq-rl-circuit}) leads to:
\begin{align*}
    V_R^\text{circuit-off}=iR&=(I_0R)e^{\displaystyle-\frac{0}{\tau_L}}=I_0R.
\end{align*}
At $t\to\infty$, Eq. (\ref{eq:decreasing-voltage-eq-rl-circuit}) leads to:
\begin{align*}
    V_R^\text{circuit-off}=iR&=(I_0R)e^{\displaystyle-\frac{\infty}{\tau_L}}=0.
\end{align*}
\subsection{Turn off the Circuit: Growth of Inductive Voltage}
Differentiating Eq. (\ref{eq:decreasing-current-eq-rl-circuit}) with time once and multiplying it by $L$ gives the voltage equation during the current surge period of the inductor.
\begin{align}
    V_L^\text{circuit-off}=L\frac{di}{dt}=-L\frac{I_0}{\tau_L}\left(e^{\displaystyle-\frac{t}{\tau_L}}\right)=-\frac{LI_0}{\displaystyle \frac{L}{R}}\left(e^{\displaystyle-\frac{t}{\tau_L}}\right)=-V_0e^{\displaystyle-\frac{t}{\tau_L}}.\label{eq:decreasing-ind-voltage-eq-rl-circuit}
\end{align}
At $t=0$, Eq. (\ref{eq:decreasing-ind-voltage-eq-rl-circuit}) leads to:
\begin{align*}
    V_L^\text{circuit-off}=-V_0e^{\displaystyle-\frac{0}{\tau_L}}=-V_0.
\end{align*}
At $t\to\infty$, Eq. (\ref{eq:decreasing-ind-voltage-eq-rl-circuit}) leads to:
\begin{align*}
    V_L^\text{circuit-off}=-V_0e^{\displaystyle-\frac{\infty}{\tau_L}}=0.
\end{align*}
\section{Electromagnetic Oscillations I: Series $LC$ Circuits}
LC oscillations refer to the periodic exchange of energy between an inductor $L$ and a capacitor $C$ in an electrical circuit. This phenomenon occurs in circuits known as \textit{LC circuit}s. LC oscillations are a form of simple harmonic motion where the electric and magnetic fields in the circuit store and exchange energy at a specific resonant frequency.
\begin{figure}
    \centering
    \includegraphics[width=.75\linewidth]{figs/lc-circuit-def.pdf}
    \caption{When switch $S$ is closed on $a$, the capacitor charges only. When the switch is afterward closed on $b$, the inductor exchanges the stored energy back and forth, creating an \textit{electromagnetic oscillation}.}
    \label{fig:lc-circuit-def}
\end{figure}
Consider a series LC circuit consisting of an inductor (L) and a capacitor (C) connected in series. The current $I$ flowing through the circuit at any time $t$ charges the capacitor and creates a magnetic field in the inductor. The voltage across the capacitor $V_C$ and the voltage across the inductor $V_L$ contribute to the total voltage $V$ across the circuit. Use KVL to find:
\begin{align}
    V=V_C+V_L&=\frac{q}{C}+L\frac{di}{dt}\notag\\
    L\frac{d^2q}{dt^2}+\frac{q}{C}&=V.
\end{align}
This is a \nth{2} order linear differential equation that describes the behavior of an LC circuit. Its solution gives the charge $q$ as a function of time $t$ and relates to the LC oscillations in the circuit.
\begin{align}
    q(t)=Q\cos{(\omega t + \phi)},\label{lc-charge-equation}
\end{align}
where $Q$ is the maximum charge, $\omega$ is the angular frequency of the oscillation, and $\phi$ is the phase angle.
\begin{figure}[H]
    \centering
    \includegraphics[width=.6\linewidth]{figs/current-charge-lc-circuit.pdf}
    \caption{When switch $S$ is closed on $a$, the capacitor charges only. When the switch is afterward closed on $b$, the inductor exchanges the stored energy back and forth, creating an \textit{electromagnetic oscillation}; while the resistor dissipates the energy with an exponentially decreasing amplitude.}
    \label{fig:current-charge-lc-circuit}
\end{figure}
If we turn on the voltage and allow the capacitor to store some charges, thus introducing electrical energy into the circuit. Turning off the source will lead to the following:
\begin{align}
    L\frac{d^2q}{dt^2}+\frac{q}{C}&=0.
\end{align}
This leads to a harmonic $LC$ oscillations of energy transfer with essentially no one to damp the oscillation. The charge oscillation (and thus, by extension, the energy oscillates) following Eq. (\ref{lc-charge-equation}).

This sinusoidal variation in charge corresponds to oscillations in the current and voltage across the inductor and capacitor, leading to LC oscillations in the circuit. The resonant frequency $\omega_0$ of these oscillations is determined by:
\begin{align}
    \omega_0=\frac{1}{\sqrt{LC}}.
\end{align}
We can measure the sinusoidal current in the circuit by differentiating Eq. (\ref{lc-charge-equation}) once:
\begin{align}
    i=\frac{dq}{dt}=-\omega Q\sin{(\omega t+\phi)}=-I\sin{(\omega t+\phi)}.
\end{align}
The maximum amplitude of this oscillating current is $I=\omega Q$.

Theoretically, this oscillation is undamped, just like a simple pendulum, because no one is present in the circuit to dissipate any energy.
\subsection{Electrical and Magnetic Energy Oscillations in $LC$ Circuits}
The electrical energy stored in the $LC$ circuit at time $t$ is,
\begin{align}
    U_E(t)=\frac{q^2}{2C}=\frac{Q^2}{2C}\cos^2{(\omega t+\phi)}.
\end{align}
The magnetic energy stored in the $LC$ circuit at time $t$ is,
\begin{align}
    U_B(t)=\frac{1}{2}Li^2=\frac{1}{2}L\omega^2Q^2\sin^2{(\omega t+\phi)}=\frac{Q^2}{2C}\sin^2{(\omega t+\phi)}.
\end{align}
\begin{figure}[H]
    \centering
    \includegraphics[width=.75\linewidth]{figs/lc-oscillation-total-energy.pdf}
    \caption{Energy oscillation graph for an LC circuit. At any instant, the total energy is constant (dashed line).}
    \label{fig:lc-circuit-undamped-energy}
\end{figure}
The total circuit energy (undamped) would then be
\begin{align}
    E_\text{no-damping}=U_E(t)+U_B(t)=\frac{Q^2}{2C}\left[\sin^2{(\omega t+\phi)}+\cos^2{(\omega t+\phi)}\right]=\frac{Q^2}{2C}.
\end{align}
Figure \ref{fig:lc-circuit-undamped-energy} plots of $U_E(t)$ and $U_B(t)$ for the case of $\phi=0$. Note that
\begin{itemize}
\setlength{\itemsep}{0pt}
    \item The maximum values of $U_E$ and $U_B$ are both $\displaystyle\frac{Q^2}{2C}$
    \item At any instant, the total energy $\displaystyle\frac{Q^2}{2C}$ is constant
    \item When $U_E$ is maximum, and $U_B$ is minimum, and vice versa.
\end{itemize}
\section{Electromagnetic Oscillations II: Series $RLC$ Circuits}
$RLC$ oscillations refer to the periodic oscillatory behavior that occurs in circuits with a resistor $R$, an inductor $L$, and a capacitor $C$ in an electrical circuit. This phenomenon occurs in circuits known as \textit{RLC circuit}s. $RLC$ oscillations are a form of damped simple harmonic motion where the electric and magnetic fields in the circuit store and exchange energy at a specific resonant frequency while the energy is dissipated with an exponential envelope.
\begin{figure}[H]
    \centering
    \includegraphics[width=.75\linewidth]{figs/rlc-circuit-def.pdf}
    \caption{When switch $S$ is closed on $a$, the capacitor charges only. When the switch is afterward closed on $b$, the inductor exchanges the stored energy back and forth, creating an \textit{electromagnetic oscillation}; while the resistor dissipates the energy with an exponentially decreasing amplitude.}
    \label{fig:rlc-circuit-def}
\end{figure}
Consider a series $RLC$ circuit consisting of a resistor $R$, an inductor $L$, and a capacitor $C$ connected in series. The current $I$ flowing through the circuit at any time $t$ charges the capacitor and creates a magnetic field in the inductor. The voltage across the resistor $V_R$, voltage across the capacitor $V_C$, and the voltage across the inductor $V_L$ contribute to the total voltage $V$ across the circuit. Use KVL to find:
\begin{align}
    V=V_R+V_C+V_L&=iR+\frac{q}{C}+L\frac{di}{dt}\notag\\
    \frac{dq}{dt}+L\frac{d^2q}{dt^2}+\frac{q}{C}&=V.
\end{align}
This is a \nth{2} order linear differential equation that describes the behavior of an RLC circuit. Its solution gives the charge $q$ as a function of time $t$ identical to Eq. (\ref{lc-charge-equation}) with a damping exponential amplitude.
\begin{align}
    q(t)=Qe^{-\alpha t}\cos{(\omega_\text{damping}\,t + \phi)},\label{eq:rlc-charge-equation}
\end{align}
where $Q$ is the maximum charge, $\displaystyle\omega_\text{damping}=\sqrt{\omega_0^2-\left(\frac{R}{2L}\right)^2}$ is the damping angular frequency of the oscillation, $\alpha$ is the damping parameter and $\phi$ is the phase angle.

If we turn on the voltage and allow the capacitor to store some charges, thus introducing electrical energy into the circuit. Turning off the source will lead to the following:
\begin{align}
    \frac{dq}{dt}+L\frac{d^2q}{dt^2}+\frac{q}{C}&=0.
\end{align}
\begin{figure}[H]
    \centering
    \includegraphics[width=.72\linewidth]{figs/current-charge-rlc-circuit.pdf}
    \caption{When switch $S$ is closed on $a$, the capacitor charges only. When the switch is afterward closed on $b$, the inductor exchanges the stored energy back and forth, creating an \textit{electromagnetic oscillation}; while the resistor dissipates the energy with an exponentially decreasing amplitude.}
    \label{fig:current-charge-lc-circuit}
\end{figure}
This leads to a damped harmonic $RLC$ oscillations of circuit energy with $R$ essentially being the damping component of the oscillation. The charge oscillation (and thus, by extension, the energy oscillates) following Eq. (\ref{eq:rlc-charge-equation}).

We can measure the sinusoidal current in the circuit by differentiating Eq. (\ref{lc-charge-equation}) once:
\begin{align}
    i=\frac{dq}{dt}&=-\omega_\text{damping}Qe^{-\alpha t}\sin{(\omega_\text{damping}\,t+\phi)}-\alpha Qe^{-\alpha t}\cos{(\omega_\text{damping}\,t + \phi)}\notag\\
    &=-\omega_\text{damping}Qe^{-\alpha t}\sin{(\omega_\text{damping}\,t+\phi)}-\alpha q(t) \notag\\[1.75pt]
    &=-\omega_\text{damping}Qe^{-\alpha t}\sin{(\omega_\text{damping}\,t+\phi)}-\alpha\int idt.\notag\\
    i-\alpha\int idt&=-\omega_\text{damping}Qe^{-\alpha t}\sin{(\omega_\text{damping}\,t+\phi)}\notag\\
    \therefore i&=-\left(\frac{1}{1-\alpha t}\right)\omega_\text{damping}Qe^{-\alpha t}\sin{(\omega_\text{damping}\,t+\phi)}\\
    &=-Ie^{-\alpha t}\sin{(\omega_\text{damping}\,t+\phi)}.
\end{align}
The decaying (maximum) amplitude of this oscillating current is $\displaystyle I=\left(\frac{1}{1-\alpha t}\right)\omega_\text{damping}Q$.

This oscillation is damped, just like a practical simple pendulum.
\subsection{Electrical and Magnetic Energy Oscillations in $RLC$ Circuits}
The electrical energy stored in the $LC$ circuit at time $t$ is,
\begin{align}
    U_E(t)=\frac{q^2}{2C}=\frac{Q^2}{2C}e^{-2\alpha t}\cos^2{(\omega_\text{damping}\,t+\phi)}.
\end{align}
The magnetic energy stored in the $LC$ circuit at time $t$ is,
\begin{align}
    U_B(t)&=\frac{1}{2}Li^2=\frac{1}{2}LI^2\sin^2{(\omega_\text{damping}\,t+\phi)}\notag\\
    &=\frac{Q^2}{2C}\left[\frac{\omega_\text{damping}}{\omega_0\left(1-\alpha t\right)}\right]^2e^{-2\alpha t}\sin^2{(\omega_\text{damping}\,t+\phi)}.
\end{align}
The total circuit energy (damped) would then be
\begin{align}
    E_\text{damping}=U_E(t)+U_B(t)&=\frac{Q^2}{2C}e^{-2\alpha t}\left[\cos^2{(\omega_\text{damping} t+\phi)-\left\{\frac{\omega_\text{damping}}{\omega_0\left(1-\alpha t\right)}\right\}^2\sin^2{(\omega_\text{damping} t+\phi)}}\right]\notag\\
    &=\frac{Q^2}{2C}e^{-2\alpha t}\cos^2(\omega_\text{damping} t+\phi)\left[1-\left\{\frac{\omega_\text{damping}}{\omega_0\left(1-\alpha t\right)}\right\}^2\tan^2{(\omega_\text{damping} t+\phi)}\right]\notag\\
    &=U_E\left[1-\left\{\frac{\omega_\text{damping}}{\omega_0\left(1-\alpha t\right)}\right\}^2\tan^2{(\omega_\text{damping} t+\phi)}\right].
\end{align}

\section{Alternating Current \& Forced Electromagnetic Oscillation}
The oscillations in an $RLC$ circuit will not damp out if an external EMF device supplies enough energy to \textit{compensate} for the energy dissipated as thermal energy in the resistance $R$. This is called \textit{Forced Oscillation}. Forced oscillation is the phenomenon of an external alternating voltage source being applied to the circuit, driving it to oscillate at a frequency determined by the source. This external voltage source is often sinusoidal and leads to generating an alternating current (AC) in the circuit. The circuit's natural frequency of oscillation $\displaystyle\omega_0=\frac{1}{\sqrt{LC}}$ and the frequency of the external driving source $\omega_d$ play a significant role in determining the response.

\textbf{Alternating current} periodically changes its direction and magnitude. It is the type of current commonly found in power grids and household electrical systems. AC is generated by power stations and transmitted over long distances. It can easily be converted to different voltages using transformers.
\begin{figure}[H]
    \centering
    \includegraphics[width=.75\linewidth]{figs/rlc-circuit-def.pdf}
    \caption{When switch $S$ is closed on $a$, the capacitor charges only. When the switch is afterward closed on $b$, the inductor exchanges the stored energy back and forth, creating an \textit{electromagnetic oscillation}, while the resistor dissipates this energy as thermal energy.}
    \label{fig:rlc-circuit-def}
\end{figure}
Consider a series $RLC$ circuit driven by an external sinusoidal EMF source $\mathcal{E}_\text{source}(t)=\mathcal{E}_m\sin(\omega_d t+\phi)$, where $\mathcal{E}_m$ is the amplitude of the EMF, $\omega_d$ is the driving angular frequency of the source, and $\omega_dt$ is the phase of the EMF. An alternator can produce this type of alternating EMF to the damping $RLC$ circuit.

The current would then be of the following form:
\begin{align}
    i(t)=I_m\sin(\omega_d t+\phi),
\end{align}
where $I_m$ is the magnitude of the driven alternating current.

\textsc{Note}: Whatever the natural angular frequency $\omega_0$ of a circuit may be, forced oscillations of charge, current, and potential difference in the circuit always occur at the driving angular frequency $\omega_d$.
\subsection{A Quick Detour through Phasors}
Phasors are a geometrical representation that simplifies the analysis of sinusoidal oscillations in alternating current (AC) circuits and wave phenomena. They allow us to analyze complex oscillatory behaviors using simpler mathematical operations and concepts. Phasors are especially useful when dealing with the amplitude, phase, and frequency relationships of oscillating quantities. We will better understand them once we see them in action.

Phasors have the following properties:
\begin{itemize}
\setlength{\itemsep}{0pt}
    \item \textbf{Angular speed}: Phasors rotate counterclockwise (positive) about the origin with an angular speed equal to the driving angular frequency $\omega_d$
    \item \textbf{Length}: The length of each phasor represents the amplitude of the alternating quantities, $v,\,I$.
    \item \textbf{Projection}: The projection of each phasor on the vertical axis represents the value of the alternating quantity at time $t$
    \item \textbf{Rotation angle}: The rotation angle of each phasor is equal to the phase of the alternating quantity at time $t$
\end{itemize}
\textsc{Note}: Phasors are not exactly vectors in the traditional sense, but they are often treated as vectors in mathematical analysis due to their similarity in behavior and algebraic operations. Phasors are complex numbers that represent the magnitude and phase of a sinusoidal oscillation. While complex numbers have both real and imaginary components, they can also be represented as vectors in a complex plane. \textbf{Whatever You do, do not draw them in a Cartesian coordinate system with $XY$ axis}.
\subsection{Foced Oscillations in Three Simple Circuits}
Before we get to the whole $RLC$ circuit being operated under the influence of an alternating source, it is good practice to see how this EMF affects each circuit element independently.
\subsection*{A Resistive Load}
\begin{figure}[H]
    \centering
    \includegraphics[width=.75\linewidth]{figs/resistive-ac-circuit-def.pdf}
    \caption{A resistor is connected across an alternating-current generator.}
    \label{fig:resistive-ac-circuit-def}
\end{figure}
Consider the following circuit where an alternating EMF $\mathcal{E}$ is connected in series with a resistor $R$. Use KVL to find:
\begin{align}
    \mathcal{E}-v_R&=0,
\end{align}
$v_R$ is the resistive voltage drop.

This gives us 
\begin{align}
    v_R&=\mathcal{E}_m\sin(\omega_d t+\phi)=V_R\sin(\omega_d t+\phi),\label{eq:voltage-drop-resistor}
\end{align}
$V_R$ is the maximum amplitude of the resistive voltage drop.

From Ohm's Law,
\begin{align}
    i_R&=\frac{v_R}{R}=\frac{V_R}{R}\sin(\omega_d t+\phi)\notag\\
    i_R&=I_R\sin(\omega_d t+\phi),\label{eq:current-drop-resistor}
\end{align}
where $I_R$ is the maximum amplitude of the alternating current in the resistive circuit, giving us:
\begin{align}
    V_R=I_RR.
\end{align}
Comparing Eqs. (\ref{eq:voltage-drop-resistor}) and (\ref{eq:current-drop-resistor}), we see that voltage and current in a resistive circuit are in phase.
\begin{figure}[H]
    \centering
    \includegraphics[width=.85\linewidth]{figs/resistive-signals.pdf}
    \caption{(a) The current $i_R$ and the potential difference $v_R$ across the resistor are plotted on the same graph, both versus time $t$. They are in phase and complete one cycle in one period $T$. (b) A phasor diagram shows the same thing as (a).}
    \label{fig:resistive-signals}
\end{figure}
\subsection*{An Inductive Load}
\begin{figure}[H]
    \centering
    \includegraphics[width=.75\linewidth]{figs/inductive-ac-circuit-def.pdf}
    \caption{An inductor is connected across an alternating-current generator.}
    \label{fig:inductive-ac-circuit-def}
\end{figure}
Consider the following circuit where an alternating EMF $\mathcal{E}$ is connected in series with an inductor $L$. Use KVL to find:
\begin{align}
    \mathcal{E}-v_L&=0,
\end{align}
$v_L$ is the inductive voltage drop.

This gives us 
\begin{align}
    v_L&=\mathcal{E}_m\sin(\omega_d t+\phi)=V_L\sin(\omega_d t+\phi),\label{eq:voltage-drop-inductor}
\end{align}
$V_L$ is the maximum amplitude of the inductive voltage drop.

The voltage drop across an inductor,
\begin{align}
    v_L&=L\frac{di}{dt}\notag\\
    \frac{di}{dt}&=\frac{V_L}{L}\sin(\omega_d t+\phi)\notag\\
    i_L=\int di_L&=\frac{V_L}{L}\int \sin(\omega_d t+\phi)\notag\\
    &=-\left(\frac{V_L}{\omega_dL}\right)\cos(\omega_d t+\phi)\notag\\
    &=-\left(\frac{V_L}{X_L}\right)\cos(\omega_d t+\phi)\notag\\
    \therefore i_L&=\left(\frac{V_L}{X_L}\right)\sin[(\omega_d t+\phi)-90^\circ],\label{eq:current-drop-inductor}
\end{align}
where $X_L=\omega_dL$ is the \textit{inductive reactance}, giving us:
\begin{align}
    V_L=I_LX_L.
\end{align}
Comparing Eqs. (\ref{eq:voltage-drop-inductor}) and (\ref{eq:current-drop-inductor}) we see that voltage and current in an inductive circuit are out of phase. The inductive current is $90^\circ$ behind in phase with the inductive voltage. The current $I_L$ \textit{lags} the voltage $V_L$.
\begin{figure}[H]
    \centering
    \includegraphics[width=.85\linewidth]{figs/inductive-signals.pdf}
    \caption{(a) The current $i_L$ and the potential difference $v_L$ across the inductor are plotted on the same graph, both versus time $t$. The current in the inductor lags the voltage by $90^\circ$ in one period $T$. (b) A phasor diagram shows the same thing as (a).}
    \label{fig:inductive-signals}
\end{figure}
\subsection*{A Capacitive Load}
\begin{figure}[H]
    \centering
    \includegraphics[width=.75\linewidth]{figs/capacitive-ac-circuit-def.pdf}
    \caption{A capacitor is connected across an alternating-current generator.}
    \label{fig:capcitive-ac-circuit-def}
\end{figure}
Consider the following circuit where an alternating EMF $\mathcal{E}$ is connected in series with a capacitor $C$. Use KVL to find:
\begin{align}
    \mathcal{E}-v_C&=0,
\end{align}
$v_C$ is the capacitive voltage drop.

This gives us 
\begin{align}
    v_C&=\mathcal{E}_m\sin(\omega_d t+\phi)=V_C\sin(\omega_d t+\phi),\label{eq:voltage-drop-capacitor}
\end{align}
$V_C$ is the maximum amplitude of the capacitive voltage drop.

From the Capacitance definition,
\begin{align}
    q_C&=Cv_C\notag\\
    &=CV_C\sin(\omega_d t+\phi)\notag\\
    i_C=\frac{dq}{dt}&=\omega_dCV_C\cos(\omega_d t+\phi)\notag\\
    i_C&=\frac{V_C}{\left(\displaystyle\frac{1}{\omega_dC}\right)}\cos(\omega_d t+\phi)\notag\\
    i_C&=\frac{V_C}{X_C}\cos(\omega_d t+\phi)\notag\\
    i_C&=I_C\cos[(\omega_d t+\phi)+90^\circ],\label{eq:current-drop-capacitor}
\end{align}
where $I_C$ is the maximum amplitude of the alternating current in the capacitive circuit, giving us:
\begin{align}
    V_C=I_CX_C.
\end{align}
Comparing Eqs. (\ref{eq:voltage-drop-capacitor}) and (\ref{eq:current-drop-capacitor}), we see that voltage and current in an inductive circuit are out of phase. The capacitive current is $90^\circ$ ahead in phase with the capacitive voltage. The current $I_C$ \textit{leads} the voltage $V_C$.
\begin{figure}[H]
    \centering
    \includegraphics[width=.85\linewidth]{figs/capcitive-signals.pdf}
    \caption{(a) The current $i_C$ and the potential difference $v_C$ across the capacitor are plotted on the same graph, both versus time $t$. The current in the capacitor leads the voltage by $90^\circ$ in one period $T$. (b) A phasor diagram shows the same thing as (a).}
    \label{fig:capcitive-signals}
\end{figure}
\subsection{Phasor Diagrams of the Three Simple Circuits}
We try to represent the following table full of information in their respective phasor diagrams.
\begin{table}[H]
    \centering
    \begin{tabular}{c|c|c|c|c|c}
        Components & Voltage & Current & Reactance & Phase Relations & Phase Angle ($I$ with $V$)\\\hline
        Resistor & $V_R$ & $I_R$ & $R$ & $I_R$ in phase with $v_R$ & $0^\circ$ \\[6pt]\hline
        Inductor & $V_L$ & $I_L$ & $X_L=\omega_dL$ & $I_L$ \textit{lags} $V_L$ & $+90^\circ$ \\[6pt]\hline
        Capacitor & $V_C$ & $I_C$ & $\displaystyle X_C=\frac{1}{\omega_dC}$ & $I_C$ \textit{leads} $V_C$ & $-90^\circ$ \\\hline 
    \end{tabular}
    % \caption{Caption}
    \label{tab:phasor-info}
\end{table}
\begin{figure}[H]
    \centering
    \includegraphics[width=.8\linewidth]{figs/reactance-graph.pdf}
    \caption{Circuit element reactances $X$ variations with time $t$ in an alternating circuit.}
    \label{fig:reactance-graph}
\end{figure}

\subsection{The Series $RLC$ Circuits in Alternating Source}
We are now finally ready to see what happens when all three circuit elements come together in a circuit where the source has an alternating EMF.

Consider the very first circuit of this chapter (Figure \ref{fig:rlc-circuit-def}), and use KVL to find:
\begin{align}
    \mathcal{E}(t)&=V_R+V_L+V_C\notag\\
    \mathcal{E}_m\sin(\omega_dt)&=IR+L\frac{di}{dt}+\frac{q}{C}\label{eq:alternating-circuit-kvl-apply}
\end{align}
Differentiating Eq. (\ref{eq:alternating-circuit-kvl-apply}) once with $t$ we get the following:
\begin{align}
    -\omega_d\mathcal{E}_m\sin(\omega_dt)&=R\frac{dI}{dt}+L\frac{d^2i}{dt^2}+\frac{1}{C}\int Idt\label{eq:alternating-circuit-2nd-order}
\end{align}
This is a \nth{2} order linear ordinary differential equation with variable coefficients. To solve it, we assume a solution of the form $I=I_m\cos(\omega_dt+\phi)$, where $I_0$ is the amplitude of the alternating current and $\phi$ is the phase angle.

Substituting this into Eq. (\ref{eq:alternating-circuit-2nd-order}) and simplifying:
\begin{align}
    -\omega_d\mathcal{E}_m\sin(\omega_dt)&=-\omega_dRI_m\sin(\omega_dt+\phi)-\omega_d^2LI_m\cos(\omega_dt-\phi)+\frac{I_m}{C}
\end{align}
A simple matching of coefficients gives us the following relationship:
\begin{align}
    \omega_dRI_m&=\omega_d\mathcal{E}\\
    \omega_d^2LI_m&=\omega^2CI_m.
\end{align}
Solving for $I$ gets us:
\begin{align}
    I_m=I=\frac{\mathcal{E}}{\displaystyle\sqrt{R^2+\left(\omega_dL-\frac{1}{\omega_dC}\right)^2}}
\end{align}
This equation describes the amplitude of the current through the circuit in response to the external sinusoidal voltage source. The current amplitude depends on the impedance $Z$ of the circuit, given by:
\begin{align}
    Z&=\sqrt{R^2+\left(\omega_dL-\frac{1}{\omega_dC}\right)^2}=\sqrt{R^2+\left(X_L-X_C\right)^2}.
\end{align}
\subsection*{Special Cases}
\begin{itemize}
    \item $X_L=X_C$: $Z=R\to$ circuit is resistive
    \item $X_L>X_C$: $R\leq(X_L-X_C)\to$ circuit is inductive
    \item $X_L<X_C$: $R\geq(X_L-X_C)\to$ circuit is capacitive
\end{itemize}
We can overlap the (b) diagrams (\ref{fig:resistive-signals}, \ref{fig:inductive-signals}, \ref{fig:capcitive-signals}) to get a single phasor diagram representing the series $RLC$ circuit altogether.
\begin{figure}[H]
    \centering
    \includegraphics[width=.95\linewidth]{figs/phasor-diagram-complete-ac-circuit.pdf}
    \caption{A capacitor is connected across an alternating-current generator.}
    \label{fig:phasor-diagram-complete-ac-circuit}
\end{figure}
\subsection*{The Phase Constant}
For an $RLC$ circuit or any oscillatory system, the phase constant in the phasor diagram is crucial in determining the behavior of the circuit. It determines the phase relations between the source EMF and the resistive voltage. It influences how the current and voltage across different circuit elements are related to each other, affecting the circuit's overall response. 
\begin{align}
    \tan\phi&=\frac{V_L-V_C}{V_R}=\frac{IX_L-IX_C}{IR}=\frac{X_L-X_C}{R}.
\end{align}
\subsection*{Special Cases}
\begin{itemize}
    \item $X_L=X_C$: $\phi=0^\circ\to$ $V_R$ is in phase with $I$; Resonance; $\omega_d=\omega_0$
    \item $X_L>X_C$: $\phi\equiv+\text{ve}\to$ $V_L$ leads $I$
    \item $X_L<X_C$: $\phi\equiv-\text{ve}\to$ $V_L$ lags $I$
\end{itemize}
\subsection{Power in Alternating Circuits}
Recall the equation relating instantaneous power dissipated across a resistor by a current:
\begin{align}
    P&=i^2R\notag\\
    &=\left[I\sin(\omega_dt-\phi)\right]^2R\notag\\
    P&=I^2R\sin^2(\omega_dt-\phi).\label{eq:instantaneous-power-ac-circuit}
\end{align}
The average rate at which energy is dissipated in the resistor, however, is the average of Eq. (\ref{eq:instantaneous-power-ac-circuit}) over time. Over one complete cycle, the average value of $\sin\theta$, where $\theta$ is any variable, is zero. However, the average value of $\sin^2\theta$ is $\displaystyle\frac{1}{2}$.
\begin{align}
    P_\text{avg}&=(I^2R)\times\frac{1}{2}=\left(\frac{I}{\sqrt{2}}\right)^2R.\label{rms-power-def}
\end{align}
The term in the bracket is labelled the \textbf{root-mean-square}, or \textbf{RMS}, value of the current $I$. The root-mean-square (RMS) value is a way of expressing the effective or average value of an alternating current (AC) voltage, current, or power in a way that's comparable to the direct current (DC) values.

Thus, Eq. (\ref{rms-power-def}) can now be written as
\begin{align}
    P_\text{avg}=I_\text{RMS}^2R.
\end{align}
We can also define RMS values of voltages and EMFs for alternating-current
circuits:
\begin{align}
    V_\text{RMS}&=\frac{V}{\sqrt{2}}.\label{rms-voltage-def}\\
    \mathcal{E}_\text{RMS}&=\frac{\mathcal{E}_m}{\sqrt{2}}.\label{rms-emf-def}
\end{align}
Because the proportionality factor in Eqs. (\ref{rms-power-def}) and (\ref{rms-voltage-def}-\ref{rms-emf-def}) is the same
for all three variables, we can write
\begin{align}
    I_\text{RMS}=\frac{\mathcal{E}_\text{RMS}}{\displaystyle\sqrt{R^2+\left(\omega_dL-\frac{1}{\omega_dC}\right)^2}}.
\end{align}
\clearpage
%%%%%%%%%%%%%%%%%%%%%%%%%%%%%%%%%%%%%%%%%%%%%%%%%%%%%%%%%%%%%%%%%%%%%
\bigskip \
\vspace{10cm}
\vskip 20pt
\begin{center}
    \thispagestyle{empty}
    \huge{\selectfont \bf \sffamily Part IV: A \small{teensy} \huge{Glimpse into Electrodynamics}}
\end{center}
\addcontentsline{toc}{section}{\Large Part IV: A teensy Glimpse into Electrodynamics}
\vfill
\clearpage
%%%%%%%%%%%%%%%%%%%%%%%%%%%%%%%%%%%%%%%%%%%%%%%%%%%%%%%%%%%%%%%%%%%%%
\section{Maxwell's Equations and Where to Find Them}
Maxwell's four fundamental equations describe the behaviour of electric and magnetic fields and their interactions. These equations, formulated by the Scottish physicist James Clerk Maxwell in the \nth{19} century, played a pivotal role in unifying the study of electricity and magnetism and laid the foundation for the field of electromagnetism.

In fact, we have been building them throughout this course. They come in the form of Gauss's Law [\ref{gauss-differential}, \ref{gauss-differential-mag}], Faraday's Law [\ref{eq:faraday-law-differential}], and Ampère' Law [\ref{eq:ampere-law}, \ref{eq:ampere-maxwell-law}]. Four equations there are, two for $\vec{E}$-fields and two for $\vec{B}$-fields. Two measure the outflow of the two fields out of a closed surface, and the other two measure the circulation of the fields. Literally, everything we derived and \textit{suffered} through this course can be boiled down to these four elegant equations. 

But the story is not over yet. These equations hold keys to yet another paradigm-shifting outcome. You would be surprised. I know I was when I had first found out.
\subsection{Displacement Current and Ampère-Maxwell Law}
Displacement current is a concept in electromagnetism that was introduced by James Clerk Maxwell to explain certain behaviors in electric circuits. It's not an actual flow of charge like current in a wire, but rather, it's a mathematical term that accounts for changing electric fields in the context of Maxwell's equations.
\begin{figure}[H]
    \centering
    \includegraphics[width=\linewidth]{figs/displacement-current-def.pdf}
    \caption{(a) Before and (d) after the plates are charged, there is no magnetic field. (b) During the charging, a magnetic field is created by both the real and the (fictional) displacement currents. (c) The same right-hand rule works for both currents to give the direction of the magnetic field.}
    \label{fig:displacement-current}
\end{figure}
Consider a circular parallel plate capacitor setup in a circuit. When the capacitor charges, the field in between the plates $\vec{E}$ also increases. The charge $q$ on the plates at any time is related to the magnitude $E$ of the
field between the plates at that time, and the plate area $A$ by
\begin{align}
    q=\epsilon_0AE.
\end{align}
Differentiating this equation once with time gives the variation of $\vec{E}$ and a fictitious current between the plates. This current is fictitious because there is no actual flow of charge carriers in between the plates, unlike the conventional current. This fictitious current is called the \textit{displacement current}.
\begin{align}
    i_\text{disp}=\frac{dq}{dt}=\epsilon_0\frac{d(EA)}{dt}=\epsilon_0\frac{d\Phi_E}{dt}.
\end{align}
This current is generated when there's a variation in $\vec{E}$ without the involvement of moving charge carriers. Maxwell added this current contribution to \textit{fix} the Ampère's law.

With this Maxwell correction, Ampère's Law states:
\begin{align}
    \oint\vec{B}\cdot d\vec{l}=\mu_0\left(i_\text{free}+i_\text{disp}\right)
\end{align}
\subsection{Then God said the Beautiful Maxwell's Equations}
\begin{itemize}
\setlength{\itemsep}{0pt}
    \item \textbf{Gauss's Law for Electric Fields}: Electric fields are generated by isolated electric charges. Electric field lines originate from positive charges and terminate on negative charges. The differential form claims that the total electric flux (flow of electric field lines) through a closed surface is proportional to the enclosed electric charge. The integral form claims that field lines diverge from a positive charge source and converge for a negative charge.
    \begin{align}
        \vec{\nabla}\cdot\vec{E}&=\frac{\rho_\text{enc}}{\varepsilon_0}\qquad\text{(Differential Form)}\\
        \oint\vec{E}\cdot d\vec{A}&=\frac{Q_\text{enc}}{\varepsilon_0}\qquad\text{(Integral Form)}
    \end{align}
    \item \textbf{Gauss's Law for $\vec{B}$-Fields}: Unlike electric charges, there are no magnetic monopoles (sinks or sources) in nature. Magnetic field lines always form closed loops, as they cannot start or end on a single magnetic pole. The total magnetic flux through any closed surface is always zero, emphasizing that magnetic field lines always form closed loops.
    \begin{align}
        \vec{\nabla}\cdot\vec{B}&=0\qquad\text{(Differential Form)}\\
        \oint\vec{B}\cdot d\vec{A}&=0\qquad\text{(Integral Form)}
    \end{align}
    \item \textbf{Faraday's Law of Electromagnetic Induction}: A changing magnetic field creates an electric field. When the magnetic field through a closed loop changes, it induces an electric field that drives current in the loop. When a moving magnet nears a coil of wire, the changing magnetic field induces a flow of electrons (current) in the wire.
    \begin{align}
        \vec{\nabla}\times\vec{E}&=-\frac{\partial \vec{B}}{\partial t}\qquad\text{(Differential Form)}\label{eq:ampere-maxwell-law}\\
        \oint\vec{E}\cdot d\vec{l}&=-\frac{\partial}{\partial t}\oint\vec{B}\cdot d\vec{A}\qquad\text{(Integral Form)}
    \end{align}
    \item \textbf{Ampère's Law with Maxwell's Correction}: Both electric currents and changing electric fields contribute to the creation of magnetic fields.
    \begin{align}
         \vec{\nabla}\times\vec{B}&=\mu_0\vec{J}+\mu_0\epsilon_0\frac{\partial \vec{E}}{\partial t}\qquad\text{(Differential Form)}\\
         \oint\vec{B}\cdot d\vec{l}&=\mu_0I_\text{free}+\mu_0\epsilon_0\oint\vec{E}\cdot d\vec{A}\qquad\text{(Integral Form)}
    \end{align}
\end{itemize}
\subsection{And There Was Light}
Electromagnetic waves are self-propagating waves that consist of electric and magnetic fields oscillating perpendicular to each other and perpendicular to the direction of wave propagation. These waves are produced by the acceleration of charged particles or by varying electric and magnetic fields. \textbf{They are Light} as we know them to be. The light that we can see and the light we can't see.

Electromagnetic waves span a wide range of frequencies and wavelengths, forming the electromagnetic spectrum. This spectrum includes various types of waves, such as radio waves, microwaves, infrared radiation, visible light, ultraviolet radiation, X-rays, and gamma rays. Each type of wave has a different frequency and wavelength range, and they all exhibit the same basic properties of electromagnetic waves.

By combining Maxwell's equations, we can obtain a wave equation that describes the propagation of electromagnetic waves. For the electric field, we get the following when we take the Laplacian of the field $\displaystyle\nabla^2\vec{F}=\vec{\nabla}\left(\vec{\nabla}\cdot\vec{F}\right)-\vec{\nabla}\times\left(\vec{\nabla}\times\vec{F}\right)$:
\begin{align}
    \nabla^2\vec{E}-\varepsilon_0\mu_0\frac{\partial^2 \vec{E}}{\partial t^2}=0.\label{electric-field-wave-equation}
\end{align}
For the magnetic field, we get the following:
\begin{align}
    \nabla^2\vec{B}-\varepsilon_0\mu_0\frac{\partial^2 \vec{B}}{\partial t^2}=0.
\end{align}
These equations have plane wave solutions of the form:
\begin{align}
    \vec{E}(\vec{r},t)&=\vec{E}_\text{max}\sin(\vec{k}\cdot\vec{r}-\omega t)\label{electric-field-solution}\\
    \vec{B}(\vec{r},t)&=\vec{B}_\text{max}\sin(\vec{k}\cdot\vec{r}-\omega t)\label{magnetic-field-solution}
\end{align}
Substituting it into the wave equation, we get:
\begin{align}
    (-k^2+\varepsilon_0\mu_0\omega^2)\vec{E}(\vec{r},t)=0
\end{align}
For this equation to hold true for all $\vec{r}$ and $t$, the term inside the parentheses must be equal to zero:
\begin{align}
    k^2&=\varepsilon_0\mu_0\omega^2\notag\\
    k&=\pm\omega\sqrt{\varepsilon_0\mu_0}
\end{align}
Now, the speed of a wave is given by the ratio of the angular frequency $\omega$ to the wave number $k$:
\begin{align}
    v&=\frac{\omega}{k}=\frac{\omega}{\pm\omega\sqrt{\varepsilon_0\mu_0}}\notag\\
    v&=\frac{1}{\sqrt{\varepsilon_0\mu_0}}=c.
\end{align}
This is the speed of electromagnetic waves in a vacuum, equal to light speed. Electromagnetic waves travel at the speed of light because light itself is an electromagnetic wave.
\begin{figure}[H]
    \centering
    \includegraphics[width=.95\linewidth]{figs/EM_spectrum.pdf}
    \label{fig:em-spectrum}
\end{figure}
The journey is far from over here. If You ask a Physics major, actually this is where the story begins. But, I now must pull the curtain. We really ran out of time. Namárië.
\vspace{1in}
\begin{center}
    \Large \textbf{The End!}
\end{center}
\end{document}